%% ════════════════════════════════════════════════════════════════════════
%% CRF_Complete_v17_15_Part_K.tex
%% Part K — ε-Ontology Closure Arc and the Temporal Bridge
%% Contains: Parts LIII–LVI (ZC83–87 + Cross-Part Connection Map)
%% Governance: Style Guide v3.6, Upgrade Protocol v1.3,
%%             Integration Execution Script v1.2 (TASK A), Lexicon Lock v4
%% Primary audience tier: 1 (Practitioner)  |  On-ramp for tier below (0): yes
%% Communication Layer (v3.6 §31): Reader's On-Ramp present after TOC;
%%   Bridge Prose (Extension Freedom narrative) at each ZC embed (5 blocks,
%%   marked % [BRIDGE-PROSE]); \physmeaning pilot on FIND-ZC83-01;
%%   Map-Not-Reality pointer in On-Ramp. Box content unchanged (No-Loss).
%% ════════════════════════════════════════════════════════════════════════
%%
%% Part K of the v17.15 multi-part single-volume edition.
%%   Part J ends at Part LII (Connection Map of Part J)
%%               → Part K begins at \setcounter{part}{52} → LIII
%%
%% Papers integrated (5):
%%   ZC83  Floor Lattice Ordering Theorem            (opens GAP-ZC83-01)
%%   ZC84  Complete Sector Synthesis                 (synthesis; no new gap)
%%   ZC85  Observed-Bridge-Bracket Theorem           (GAP-ZC83-01 → gated)
%%   ZC86  Strong Running Bridge Theorem             (CLOSES GAP-ZC81-01;
%%                                                    unlocks GAP-ZC83-01 mag.)
%%   ZC87  Temporal Continuum Bridge (v2)            (CLOSES GAP-ZC87-01;
%%                                                    GAP-ZC76-01 structural)
%%
%% Gap lifecycle (self-contained within Part K):
%%   GAP-ZC83-01 (ε_SC ontology)  opened ZC83 → ontology CLOSED ZC85
%%       → magnitude gated on GAP-ZC81-01 (COR-ZC85-01)
%%       → GAP-ZC81-01 CLOSED ZC86 → ε_SC magnitude CLOSED (FIND-ZC86-01)
%%   GAP-ZC87-01 (joint-limit well-definedness)  CLOSED in ZC87 v2
%%   GAP-ZC76-01 (τ_sweep)  structural bridge CLOSED ZC87; seconds-calib OPEN
%%
%% Gap closures in this Part:
%%   GAP-ZC81-01 (Strong running bridge)      → CLOSED  (ZC86/THM-ZC86-01)
%%   GAP-ZC83-01 (ε_SC ontology + closed form)→ CLOSED  (ZC85+ZC86)
%%   GAP-ZC87-01 (joint-limit well-defined)   → CLOSED  (ZC87/THM-ZC87-01)
%%   GAP-ZC76-01 (τ_sweep structural)         → CLOSED  (structural only)
%%
%% Phase Misalignments introduced:
%%   PM-ZC83-01           (AM ≠ ε_SC probe; pattern: false-pattern test)
%%   PM-ZC84-SCOPE-01     (synthesis — all identities inherited, none new)
%%   PM-ZC85-01           (ε_SC = GM(bridge) is RESTATEMENT, not derivation)
%%   PM-ZC86-01           (universal-coupling propagation; failed path)
%%   PM-ZC87-01           (Dirac-large-number probe that fed its own input)
%%   PM-ZC87-02           (v1 spurious dt = τ_sweep/T_avg; corrected v2)
%%
%% Notation overrides (No-Loss, edition-level — same pattern as Part J):
%%   PM-V17_15-EPSFAMILY: the v17.6/7 chain defines several ε-family macros
%%     with a zero subscript or sector-tagged form (\epsEM=ε_0^{EM},
%%     \epscan=ε_0^{can}, \epsSC=ε_{can}^{SC}, \epsS=ε_{0,Spont}).
%%     ZC83–87 (Style Guide v3.5) use the bare sector form
%%     (ε_{EM}, ε_{can}, ε_{SC}, ε_{strong}). Local \renewcommand applied in
%%     this Part's preamble; chain + ZC sources unmodified. Extends the
%%     PM-V17_13-EPSEM override of Part J to the full ε family.
%%
%% Discovery origins: all 5 papers carry %% Discovery origin: fields.
%% ════════════════════════════════════════════════════════════════════════

\documentclass[12pt, a4paper]{article}

%% ── PACKAGES ────────────────────────────────────────────────────────────
\usepackage{amsmath, amssymb, amsthm}
\usepackage{mathtools}
\usepackage{geometry}
\usepackage{xcolor}
\usepackage{parskip}
\usepackage[protrusion=false,expansion=false]{microtype}
\usepackage{hyperref}
\usepackage{tcolorbox}
\usepackage{booktabs}
\usepackage{longtable}
\usepackage{array}
\usepackage[T1]{fontenc}
\usepackage{graphicx}
\usepackage{enumitem}
\usepackage{needspace}
\usepackage{tocloft}

%% Unicode-engine font support; morewrites prevents \write exhaustion.
\usepackage{iftex}
\iftutex
  \usepackage{morewrites}
  \usepackage{fontspec}
  \setmainfont{Latin Modern Roman}
  \usepackage{polyglossia}
  \setdefaultlanguage{english}
  \setotherlanguage{thai}
  \newfontfamily\thaifont[Scale=0.95]{Loma}
\fi

\geometry{margin=2.5cm}

%% TOC spacing
\setlength{\cftbeforesecskip}{0pt}
\setlength{\cftbeforesubsecskip}{0pt}
\setlength{\cftbeforepartskip}{6pt}
\setlength{\cftsubsecindent}{2.5em}
\setlength{\cftsubsubsecindent}{4.5em}
\renewcommand{\cftsecfont}{\normalfont}
\renewcommand{\cftsubsecfont}{\small}
\renewcommand{\cftsecpagefont}{\normalfont}
\renewcommand{\cftsubsecpagefont}{\small}

%% ── COLOR PALETTE (Style Guide v3.5) ────────────────────────────────────
\definecolor{deepblue}{RGB}{10,30,80}
\definecolor{crfgreen}{RGB}{0,100,60}
\definecolor{softgreen}{RGB}{180,230,210}
\definecolor{physgold}{RGB}{140,100,0}
\definecolor{softgold}{RGB}{255,240,180}
\definecolor{loopred}{RGB}{160,30,30}
\definecolor{softred}{RGB}{255,210,210}
\definecolor{mutedgray}{RGB}{90,90,90}
\definecolor{midblue}{RGB}{30,80,160}
\definecolor{softblue}{RGB}{180,210,240}
\definecolor{layerblue}{RGB}{20,60,120}
\definecolor{genealogygold}{RGB}{120,80,0}
\definecolor{softgenealogy}{RGB}{255,248,220}
\definecolor{conmapbg}{RGB}{240,245,255}
\definecolor{conmapframe}{RGB}{50,80,140}
\definecolor{inheritbg}{RGB}{225,240,255}
\definecolor{inheritframe}{RGB}{20,60,160}
\definecolor{correctbg}{RGB}{255,235,220}
\definecolor{correctframe}{RGB}{180,90,30}
\definecolor{giftgold}{RGB}{180,140,0}
\definecolor{softgift}{RGB}{255,252,200}
\definecolor{nobelblue}{RGB}{10,40,100}
\definecolor{softnobel}{RGB}{200,220,255}

\hypersetup{colorlinks=true, linkcolor=deepblue, urlcolor=midblue,
            citecolor=deepblue, bookmarks=false}

%% ── BOX STYLES (Style Guide v3.5 — all 10 types) ────────────────────────
\tcbuselibrary{skins, breakable}

\newtcolorbox{derivedbox}[1][]{colback=softgreen!50,
  colframe=crfgreen!60, fonttitle=\bfseries\color{crfgreen!20!black},
  title={Derived}, breakable, #1}
\newtcolorbox{formalbox}[1][]{colback=softblue!40,
  colframe=midblue!60, fonttitle=\bfseries\color{deepblue},
  title={Formal}, breakable, #1}
\newtcolorbox{openqbox}[1][]{colback=softred!30,
  colframe=loopred!50, fonttitle=\bfseries\color{loopred!20!black},
  title={Open}, breakable, #1}
\newtcolorbox{keybox}[1][]{colback=softgold!50,
  colframe=physgold!60, fonttitle=\bfseries\color{physgold!20!black},
  breakable, #1}
\newtcolorbox{layerbox}[1][]{
  colback=layerblue!8, colframe=layerblue!50,
  fonttitle=\bfseries\color{layerblue},
  title={R-Layer Note}, breakable, #1}
\newtcolorbox{genealogybox}[1][]{
  colback=softgenealogy!60, colframe=genealogygold!70,
  fonttitle=\bfseries\color{genealogygold!30!black},
  title={Genealogy Map}, breakable, #1}
\newtcolorbox{conmapbox}[1][]{
  colback=conmapbg, colframe=conmapframe,
  fonttitle=\bfseries\color{white},
  breakable, #1}
\newtcolorbox{inheritbox}[1][]{
  colback=inheritbg, colframe=inheritframe,
  fonttitle=\bfseries\color{white},
  title={Backward Inheritance}, breakable, #1}
\newtcolorbox{correctionbox}[1][]{
  colback=correctbg, colframe=correctframe,
  fonttitle=\bfseries\color{white},
  title={Forward Correction / Cross-Part PM}, breakable, #1}
\newtcolorbox{contentsbox}[1][]{colback=softblue!20,
  colframe=midblue!50, fonttitle=\bfseries\color{deepblue},
  breakable, #1}
%% Style Guide v3.5: giftbox (ZC65+) and nobelbox (ZC70+)
\newtcolorbox{giftbox}[1][]{
  colback=softgift!70, colframe=giftgold!80,
  fonttitle=\bfseries\color{giftgold!20!black},
  title={Gift}, breakable, #1}
\newtcolorbox{nobelbox}[1][]{
  colback=softnobel!60, colframe=nobelblue!70,
  fonttitle=\bfseries\color{nobelblue},
  title={Nobel-Table Claim}, breakable, #1}

%% ── Theorem environments ────────────────────────────────────────────────
\newtheorem{theorem}{Theorem}[section]
\newtheorem{lemma}{Lemma}[section]
\newtheorem{corollary}{Corollary}[section]
\newtheorem{finding}{Finding}[section]
\newtheorem{proposition}{Proposition}[section]
\newtheorem{definition}{Definition}[section]
\newtheorem{principle}{Principle}[section]
\newtheorem{claim}{Claim}[section]
\newtheorem{conjecture}{Conjecture}[section]
\newtheorem*{remark}{Remark}

%% volumebox (Part overview summary)
\newtcolorbox{volumebox}[1][]{colback=softblue!15,
  colframe=deepblue!50, fonttitle=\bfseries\color{deepblue},
  breakable, #1}


%% ── Title ───────────────────────────────────────────────────────────────
\title{%
  \textbf{\color{deepblue}CRF Complete Edition v17.15 --- Part K}\\[0.4em]
  \large ZC83--87: Floor-Lattice Ordering $\cdot$ Complete Sector Synthesis\\[0.2em]
  \normalsize\color{mutedgray}
  Observed-Bridge Bracket $\cdot$ Strong Running Bridge $\cdot$
  Temporal Continuum Bridge\\[0.2em]
  \textbf{\color{crfgreen}Four closures: floor lattice ordered,
  $\varepsilon_{\rm SC}$ ontology, Strong running bridge, temporal $dt$}
}
\author{Ougit Jarittum\\
\small Independent Researcher, Thailand\quad
ORCID: 0009-0009-4451-6811\\
\small Zenodo: \href{https://doi.org/10.5281/zenodo.18977059}%
{10.5281/zenodo.18977059}\quad (CC-BY-NC 4.0)
}
\date{June 2026 --- Complete Edition v17.15 (Part K of 11)}

%% ── PART COUNTER: continue from Part J (ends at LII = 52) ──────────────
\setcounter{part}{52}

%% ── MACRO PATCH v17.15 ──────────────────────────────────────────────────
\input{CRF_v17_15_macro_patch.tex}

%% ── Part K local notation overrides (No-Loss, edition-level) ──────────────
%% PM-V17_15-EPSFAMILY (see head comment). The chain tags ε-family macros
%% with zero subscripts / sector forms; ZC83–87 (Style Guide v3.5) display
%% them bare. Override here; chain + sources unmodified. \epsW and \epsuniv
%% already match the chain (no override needed).
\renewcommand{\epsEM}{\varepsilon_{\rm EM}}     % chain: ε_0^{EM}
\renewcommand{\epscan}{\varepsilon_{\rm can}}   % chain: ε_0^{can}
\renewcommand{\epsSC}{\varepsilon_{\rm SC}}     % chain: ε_{can}^{SC}
\renewcommand{\epsS}{\varepsilon_{\rm strong}}  % chain first-wins: ε_{0,Spont}

%% ════════════════════════════════════════════════════════════════════════
\begin{document}
\sloppy

%% Governance: Upgrade Protocol v1.2, CRF Style Guide v3.5,
%%             Integration Execution Script v1.1, Lexicon Lock v4

\maketitle
\tableofcontents
\newpage

%% ════ Reader's On-Ramp (Protocol v1.3 §U2; Style Guide v3.6 §31.2) ═══════
%% Primary audience tier: 1   |   On-ramp for tier below (Tier 0): yes
\begin{keybox}[title={If you are just arriving --- Part K in plain language}]
\sloppy
CRF attempts to derive the constants and force structure of physics from a
single primitive called Distinction ($\delta$). Earlier Parts built four
tiny ``instability floors'' --- thresholds where a force sector tips from
one stable state to another --- and noticed they sit very close together.
\textbf{Part K} settles four loose ends about them: it proves \emph{why}
the floors fall in the order they do, shows that one observed value is the
exact geometric mean of two others, confirms the Strong force runs the same
way electromagnetism does, and finally connects CRF's discrete ``ticks'' of
time to ordinary continuous time.

\textbf{Minimum vocabulary:}
\textit{$\delta$ (Distinction)} --- the one starting idea everything is
built from;\;
\textit{$\varepsilon$ floors} --- four close-but-distinct instability
thresholds, one per force sector;\;
\textit{geometric mean} --- $\sqrt{ab}$, the kind of average that turns out
to relate the floors exactly;\;
\textit{running} --- how a coupling changes with scale;\;
\textit{gap / GAP-\dots} --- an open question the chain has not yet closed.
Full symbol list: Master Notation Key (front matter).

These results are pieces of a map under construction, not a claim of final
truth: each carries the conditions under which it would fail (see front
matter, ``Map, not reality'').
\end{keybox}

\newpage

%%=============================================================
\part{ZC83--84 Cluster: Floor-Lattice Ordering and Complete
  Sector Synthesis (v17.15)}
\label{part:LIII}
%%=============================================================

\begin{volumebox}[title={Part LIII --- ZC83--84: The Floor Lattice,
  Ordered and Synthesised (v17.15)}]
\sloppy
\textbf{Part LIII} opens the $\varepsilon$-ontology closure arc with two
papers: the first orders the complete instability-floor lattice and the
second assembles every $\mathbb{Z}_{15}$ sector result into one
falsifiable synthesis.

\textbf{ZC83} (The Floor Lattice Ordering Theorem) proves that the
ordering $\varepsilon_{\rm EM} < \varepsilon_{\rm can} <
\varepsilon_{\rm univ} < \varepsilon_{\rm W}$ is forced by
$\mathbb{Z}_{15}$ sector-weight arithmetic (THM-ZC83-01), and delivers
the central gift FIND-ZC83-01: the bridge floor is the \emph{geometric}
mean of the two sector floors, exactly
($\sqrt{\varepsilon_{\rm EM}\,\varepsilon_{\rm W}} =
\varepsilon_{\rm univ}$), while the arithmetic mean is strictly larger
(FIND-ZC83-02). It opens GAP-ZC83-01 (the ontology of the observed
bracket $\varepsilon_{\rm SC}=0.007605$).

\textbf{ZC84} (The Complete Sector Synthesis) collects eight
simultaneous consequences of $\delta$ alone into THM-ZC84-01 --- the full
four-sector floor lattice, the geometric-mean identity, the mirror
symmetry $\Gamma_{\rm EM}+\Gamma_{\rm strong}=2\Gamma_{\rm W}$, the
cosmological lock $w_{\rm DE}+3\Delta\Gamma=\Gamma_{\rm CRF}$, and the
uniqueness of $\mathbb{Z}_{15}$ --- and presents the four-sector
falsification table (COR-ZC84-01) with zero free parameters.

\textbf{Key results:}
THM-ZC83-01 (floor lattice ordering, forced); FIND-ZC83-01
(geometric-mean identity, exact); THM-ZC84-01 (eight-fold synthesis);
COR-ZC84-01 (four-sector falsification table).

\textbf{Gap opened:} GAP-ZC83-01 ($\varepsilon_{\rm SC}$ ontology and
closed form) --- carried forward to Part LIV.

\textbf{Phase Misalignments preserved:}
PM-ZC83-01 (the ``is $\varepsilon_{\rm SC}$ the arithmetic mean?'' probe;
pattern: false-pattern test, ruled out exactly);
PM-ZC84-SCOPE-01 (synthesis paper --- every identity is inherited, none
newly derived).
\end{volumebox}

\newpage
%%─── VERBATIM EMBED: ZC83 ───────────────────────────────────────────────
\addcontentsline{toc}{section}{ZC83: The Floor Lattice Ordering Theorem}

% [BRIDGE-PROSE: integration-added, Extension Freedom narrative, v3.6 §31.3]
Of all the quantities CRF derives, the four instability floors are the
ones that should make a careful reader uneasy. They are numbers ---
$0.007494$, $0.007550$, $0.007599$, $0.007705$ --- and they sit so close
together that the eye wants to call them the same value measured four
times. They are not. Each belongs to a different force sector, each was
derived through a different sub-chain, and yet they line up in a strict
order with gaps of less than a tenth of a percent between neighbours.

Until this paper, that ordering was an observation, not a result. The
floors had been computed one at a time, in separate papers, for separate
reasons; nobody had gone back and asked whether their arrangement was
forced or merely how things happened to fall out. If it were accidental,
the whole picture would be fragile --- a different choice upstream could
reshuffle them and nothing would forbid it. If it were forced, then the
ordering itself carries information about how the sectors relate.

ZC83 settles the question in the second sense, and in doing so turns up
something sharper than expected: the middle floor is not merely
\emph{between} its neighbours, it is their exact geometric mean. That is
the kind of relation that does not happen by accident. The sections below
state the ordering as a theorem, exhibit the two-line proof of the mean
identity, and record honestly the one quantity this still does not
explain --- the observed bracket that sits a hair above the rest, which
becomes the open question carried into the next cluster.



%%=============================================================
\begin{abstract}
\sloppy
ZC80 collected seven electroweak observables at $R_0$; ZC81 added
the Strong sector floor; but the \emph{ordering} of the five
$\varepsilon$-floor values had never been proved from first principles.
This paper closes that gap.

\textbf{THM-ZC83-01} (Floor Lattice Ordering, Near-Formal):
\[
  \epsEM \;=\; \epsS \;<\; \epscan \;<\; \epsuniv \;<\; \epsW,
\]
proved in three steps from the $\mathbb{Z}_{15}$ sector-weight
ratio $r = C_W/C_{\rm EM}$ and the two-layer gap of THM-ZC71-01.

\textbf{FIND-ZC83-01} (Geometric-Mean Identity, Exact --- Gift):
\[
  \sqrt{\epsEM\cdot\epsW} \;=\; \epsuniv \quad\text{exactly.}
\]
Two-line algebraic proof: the sector-weight exponents cancel.
The bridge floor is the geometric mean of the two sector floors ---
a consequence of the $\pm 1/4$ exponent structure of LEM-ZC67-01.

\textbf{FIND-ZC83-02}: $\AMfloor(\epsEM,\epsW) > \epsuniv$
(strict AM--GM; AM exceeds $\epsuniv$ by $0.0096\%$).
\textbf{PM-ZC83-01} registers that $\epsSC = 0.007605$ is neither
the AM nor the GM of the defined floor pair, opening
\textbf{GAP-ZC83-01} (eps\_SC ontology).
\end{abstract}


%%=============================================================
\subsection{Background and Scope}
\label{sec:background}
%%=============================================================

\begin{keybox}[title={Numerical Summary --- ZC83 (T-canonical, $R_0$)}]
{\small\renewcommand{\arraystretch}{1.25}
\begin{center}
\begin{tabular}{@{}p{5.2cm}p{2.2cm}p{1.6cm}p{3.2cm}@{}}
\toprule
\textbf{Quantity} & \textbf{Value} & \textbf{Gap} & \textbf{Source} \\
\midrule
$\epsEM = \epsS$             & $0.007494$ & exact & ZC70/ZC81 \\
$\epscan = \epsTLS$          & $0.007550$ & exact & ZC79/THM-ZC79-01 \\
$\epsuniv$                   & $0.007599$ & exact & ZC69/COR-ZC69-02 \\
$\epsW$                      & $0.007705$ & exact & ZC70/THM-ZC70-01 \\
\midrule
$r = C_W/C_{\rm EM}$        & $1.056960$ & exact & ZC57/THM-ZC57-01 \\
$x = n\,\epscan$             & $0.060400$ & exact & T-canonical \\
\midrule
$\sqrt{\epsEM\cdot\epsW}$ (NEW) & $0.007599$ & $0.000\%$ & \textbf{FIND-ZC83-01} \\
$\AMfloor(\epsEM,\epsW)$ (NEW)  & $0.007599$ & $+0.0096\%$ & \textbf{FIND-ZC83-02} \\
$\AMfloor - \epsuniv$        & $7.29\times10^{-7}$ & & AM--GM gap \\
$\epsSC$ (observed bracket)  & $0.007605$ & $+0.080\%$ & ZC68 \\
$\epsSC - \AMfloor$          & $5.54\times10^{-6}$ & $+0.073\%$ & PM-ZC83-01 \\
\bottomrule
\end{tabular}
\end{center}}
\end{keybox}

\subsubsection{The Open Item}

ZC81 completed the floor lattice with five values:
\[
  \epsEM = \epsS = 0.007494, \quad \epscan = 0.007550, \quad
  \epsuniv = 0.007599, \quad \epsW = 0.007705.
\]
ZC80/THM-ZC80-01 collected all seven $R_0$-layer electroweak
observables. Yet no paper had \emph{proved} the ordering of these
five values from the $\delta$-chain axioms. The ordering was known
numerically but its algebraic origin was not stated as a theorem.

Two questions arise naturally:
\begin{itemize}[leftmargin=2em, noitemsep]
  \item Why does $\epsEM < \epscan$ hold? (The two quantities come
    from different sub-chains: ZC70 gives $\epsEM$ from the Born
    bridge, ZC79 gives $\epscan$ from the spectral quadratic.)
  \item What is the algebraic relationship between $\epsuniv$ and the
    sector floor pair $(\epsEM, \epsW)$?
\end{itemize}

\subsubsection{What Was Already Known}

\begin{keybox}[title={Pre-ZC83 Status of the Floor Lattice}]
\textbf{Known (exact):}
\begin{itemize}[leftmargin=2em, noitemsep]
  \item $\epsEM = \epsuniv(C_{\rm EM}/C_W)^{1/4}$ and
    $\epsW = \epsuniv(C_W/C_{\rm EM})^{1/4}$ (ZC67/LEM-ZC67-01;
    ZC70/THM-ZC70-01).
  \item $\epscan < \epsuniv$ --- the two-layer gap $0.66\%$,
    irreducible (ZC71/THM-ZC71-01).
  \item $\epsEM = \epsS$ --- Strong floor equals EM floor
    (ZC81/THM-ZC81-01).
  \item All five numerical values (ZC70, ZC79, ZC81).
\end{itemize}

\textbf{Open before ZC83:}
\begin{itemize}[leftmargin=2em, noitemsep]
  \item $\epsEM < \epscan$ --- not proved algebraically.
  \item Complete ordering from first principles --- not stated as theorem.
  \item Relationship between $\epsuniv$ and $(\epsEM, \epsW)$ as a
    pair --- not identified.
\end{itemize}
\end{keybox}

\begin{keybox}[title={What ZC83 Does NOT Do}]
\begin{itemize}[leftmargin=2em, noitemsep]
  \item Does not re-derive $\epsEM$, $\epscan$, $\epsuniv$,
    $\epsW$ (used as exact pillars).
  \item Does not close GAP-ZC83-01 ($\epsSC$ ontology; opened here,
    deferred).
  \item Does not modify THM-ZC71-01 (two-layer gap remains irreducible).
  \item Does not address GAP-ZC81-01 (Strong running bridge).
\end{itemize}
\end{keybox}

\begin{keybox}[title={Companion Papers --- ZC83}]
{\small\renewcommand{\arraystretch}{1.25}
\begin{center}
\begin{tabular}{@{}p{4.8cm}p{7.5cm}@{}}
\toprule
\textbf{Paper} & \textbf{What it provides to ZC83} \\
\midrule
\texttt{CRF\_ZC67\_v1.tex}
  & LEM-ZC67-01: sector-floor formula
    $\varepsilon_i = \epsuniv(C_i/C_{\rm ref})^{1/4}$ \\
\texttt{CRF\_ZC69\_v1.tex}
  & COR-ZC69-02: $\epsuniv$ exact closed form \\
\texttt{CRF\_ZC70\_v1.tex}
  & THM-ZC70-01: $\epsEM$, $\epsW$ exact values \\
\texttt{CRF\_ZC71\_v1.tex}
  & THM-ZC71-01: two-layer gap $\epscan < \epsuniv$,
    irreducible \\
\texttt{CRF\_ZC79\_v1.tex}
  & THM-ZC79-01: $\epscan$ exact from spectral quadratic \\
\texttt{CRF\_ZC81\_v2.tex}
  & THM-ZC81-01: $\epsS = \epsEM$ \\
\bottomrule
\end{tabular}
\end{center}}
\end{keybox}

\begin{derivedbox}[title={Discovery Note}]
This paper began when a cross-paper audit (June 2026) observed that
the five floor values had always been stated as numerical facts
without a first-principles ordering proof. Checking the geometric
mean $\sqrt{\epsEM\cdot\epsW}$ against $\epsuniv$ returned a gap
of $0.000\%$ --- a two-line algebraic consequence of the $\pm 1/4$
exponent structure that had never been stated as a theorem.
The proof direction became clear immediately: the sector-weight
exponents cancel exactly. Checking whether $\epsSC$ equals the
arithmetic mean returned a $0.073\%$ mismatch, registering
PM-ZC83-01 and opening GAP-ZC83-01.
\end{derivedbox}

\subsubsection{Pre-Work Audit (PWA-1 through PWA-10)}

\begin{formalbox}[title={Pre-Work Audit Record --- ZC83}]
Scanned before writing (Upgrade Protocol v1.2):
\begin{itemize}[leftmargin=2em, noitemsep]
  \item ZC67/LEM-ZC67-01 --- \textbf{KEY PILLAR}
    (sector-floor formula with $\pm 1/4$ exponent structure).
  \item ZC70/THM-ZC70-01 --- \textbf{KEY PILLAR}
    ($\epsEM$, $\epsW$ exact; $r = C_W/C_{\rm EM}$).
  \item ZC71/THM-ZC71-01 --- \textbf{KEY PILLAR}
    (two-layer gap $\epscan < \epsuniv$, irreducible; proof
    that $\epscan$ and $\epsuniv$ come from distinct invariants).
  \item ZC79/THM-ZC79-01 --- pillar
    ($\epscan$ exact from spectral quadratic; $x = n\epscan$).
  \item ZC81/THM-ZC81-01 --- pillar ($\epsS = \epsEM$; extends
    ordering to Strong sector without new proof).
  \item ZC57/THM-ZC57-01 --- pillar
    ($C_{\rm EM} = 1+x$, $C_W = 1+2x$; Born-chain weights).
  \item ZC68 --- scanned for $\epsSC = 0.007605$ context
    (THM-ZC68-01b; PM-ZC83-01 builds on this).
\end{itemize}

\textbf{PWA-5} (tautology): FIND-ZC83-01 is not circular.
$\sqrt{\epsEM\cdot\epsW} = \epsuniv$ is an algebraic consequence of
LEM-ZC67-01; it is new content because no prior paper stated it.\\
\textbf{PWA-6} (macro): $\epsS$, $\epsSC$, $\rweight$ added via
\verb|\providecommand|; no conflict.\\
\textbf{PWA-7} (sector-floor): ordering proof uses floor formula ---
PWA-7 satisfied; no cross-sector mixing claimed.\\
\textbf{PWA-8} (layer self-consistency): explicit two-layer
distinction maintained throughout.\\
\textbf{PWA-10} (tautology): PM-ZC83-01 pre-audited as tautology
test (AM probe uses numerical values, not definition).\\
\textbf{No duplicate atomic units found.}
\end{formalbox}

\newpage

%%=============================================================
\subsection{LEM-ZC83-01\quad Sector-Weight Ordering}
\label{sec:lem01}
%%=============================================================

\needspace{10\baselineskip}

\begin{formalbox}[title={LEM-ZC83-01: $\epsEM < \epsuniv < \epsW$
  Follows from $x > 0$ ($R_0$, Exact)}]

\textbf{Statement.}
Let $r := C_W/C_{\rm EM} = (1+2x)/(1+x)$ and
$x = n\,\epscan > 0$. Then:
\begin{equation}
  \epsEM < \epsuniv < \epsW.
  \label{eq:lem01}
\end{equation}

\textbf{Proof.}

\textit{Step 1 ($r > 1$ from $x > 0$).}
$r = (1+2x)/(1+x) > 1$ iff $(1+2x) > (1+x)$ iff $x > 0$.
Since $\epscan > 0$ (from THM-ZC79-01, the TLS quadratic has a positive
root) and $n = 8 > 0$, we have $x = n\epscan > 0$, hence $r > 1$.

\textit{Step 2 (Apply LEM-ZC67-01).}
From the sector-floor formula (ZC67/LEM-ZC67-01):
\begin{equation}
  \epsEM = \epsuniv\left(\frac{C_{\rm EM}}{C_W}\right)^{1/4}
         = \epsuniv\cdot r^{-1/4}, \qquad
  \epsW  = \epsuniv\left(\frac{C_W}{C_{\rm EM}}\right)^{1/4}
         = \epsuniv\cdot r^{+1/4}.
  \label{eq:floor-formula}
\end{equation}
Since $r > 1$: $r^{-1/4} < 1 < r^{+1/4}$, hence
$\epsEM = \epsuniv\cdot r^{-1/4} < \epsuniv < \epsuniv\cdot r^{+1/4} = \epsW$.
\hfill$\square$

\textbf{T-canonical numerical verification:}
\begin{center}
{\small\renewcommand{\arraystretch}{1.25}
\begin{tabular}{@{}lrl@{}}
\toprule
Quantity & Value & Source \\
\midrule
$x = n\,\epscan$                & $0.060400$ & T-canonical \\
$C_{\rm EM} = 1+x$              & $1.060400$ & ZC57 \\
$C_W = 1+2x$                    & $1.120800$ & ZC57 \\
$r = C_W/C_{\rm EM}$            & $1.056960$ & algebra \\
$r^{-1/4}$                      & $0.986246$ & \\
$\epsEM = \epsuniv\cdot r^{-1/4}$ & $0.007494$ & confirms ZC70 \\
$\epsW  = \epsuniv\cdot r^{+1/4}$ & $0.007705$ & confirms ZC70 \\
\bottomrule
\end{tabular}}
\end{center}

\textbf{Classification:} Exact Lemma.
\end{formalbox}

%%=============================================================
\subsection{LEM-ZC83-02\quad Cross-Layer Ordering: $\epsEM < \epscan$}
\label{sec:lem02}
%%=============================================================

\needspace{10\baselineskip}

\begin{formalbox}[title={LEM-ZC83-02: $\epsEM < \epscan$
  ($R_0$, Near-Formal)}]

\textbf{Statement.}
\begin{equation}
  \epsEM \;<\; \epscan.
  \label{eq:lem02}
\end{equation}

\textbf{Proof.}

\textit{Step 1 (Reformulate condition).}
Using $\epsEM = \epsuniv\cdot r^{-1/4}$ from LEM-ZC83-01,
inequality \eqref{eq:lem02} is equivalent to:
\begin{equation}
  \epsuniv\cdot r^{-1/4} < \epscan
  \quad\Longleftrightarrow\quad
  \frac{\epsuniv}{\epscan} < r^{1/4}.
  \label{eq:lem02-cond}
\end{equation}

\textit{Step 2 (Evaluate both sides).}
The left-hand side uses the two-layer gap of THM-ZC71-01:
\[
  \frac{\epsuniv}{\epscan}
  = 1 + \frac{\epsuniv - \epscan}{\epscan}
  = 1 + \frac{0.000049}{0.007550}
  = 1.006455.
\]
The right-hand side uses $r$ from T-canonical ($x = 0.0604$):
\[
  r^{1/4} = \left(\frac{1+2x}{1+x}\right)^{1/4}
  = \left(\frac{1.1208}{1.0604}\right)^{1/4}
  = 1.013945.
\]
Since $1.006455 < 1.013945$, condition \eqref{eq:lem02-cond} holds.
\hfill$\square$

\textbf{Physical reading.}
The EM sector floor $\epsEM$ lies below the spectral floor $\epscan$
because the sector-weight suppression factor $r^{-1/4} < 1$ is larger
in magnitude than the fractional two-layer gap $(\epsuniv -
\epscan)/\epscan \approx 0.65\%$. The two quantities --- the
Born-chain sector ratio $r$ and the spectral/bridge two-layer
separation --- are independently determined and their relative
magnitudes produce a non-trivial ordering.

\textbf{Classification:} Near-Formal Lemma (the inequality
holds for $x = 0.0604$ exact from T-canonical; a formal algebraic
proof for all $x$ in the physical range is deferred to GAP-ZC83-02).
\end{formalbox}

%%=============================================================
\subsection{THM-ZC83-01\quad Floor Lattice Ordering Theorem}
\label{sec:thm01}
%%=============================================================

\needspace{12\baselineskip}

\begin{formalbox}[title={THM-ZC83-01: Floor Lattice Ordering
  ($R_0$, Near-Formal)}]

\textbf{Statement.}
At T-canonical gauge ($d_s = 3$, $n_m = 5$, $n = 8$):
\begin{equation}
  \boxed{
    \epsEM \;=\; \epsS \;<\; \epscan \;<\; \epsuniv \;<\; \epsW.
  }
  \label{eq:ordering}
\end{equation}

\textbf{Proof.}

\textit{Step 1 ($\epsEM < \epsuniv < \epsW$).}
Follows from LEM-ZC83-01: the ordering is forced by $r > 1$,
which is forced by $x = n\epscan > 0$.

\textit{Step 2 ($\epscan < \epsuniv$).}
From THM-ZC71-01 (Two-Layer Floor Irreducibility): the spectral and
bridge self-consistency equations probe distinct invariants of the
$\mathbb{Z}_{15}$ Laplacian. Their gap $\epsuniv - \epscan = 0.66\%$
is irreducible. In particular, $\epscan < \epsuniv$.

\textit{Step 3 ($\epsEM < \epscan$).}
From LEM-ZC83-02: the condition $\epsuniv/\epscan < r^{1/4}$ holds
at T-canonical, hence $\epsEM = \epsuniv\cdot r^{-1/4} < \epscan$.

\textit{Step 4 ($\epsS = \epsEM$).}
From THM-ZC81-01 (Strong Floor Identity): the Strong and EM sectors
share the same spatial-type crossing cost, hence the same floor.

Combining Steps 1--4:
$\epsEM = \epsS < \epscan < \epsuniv < \epsW$.
\hfill$\square$

\textbf{Complete ordering summary:}
\begin{center}
{\small\renewcommand{\arraystretch}{1.35}
\begin{tabular}{@{}p{2.8cm}p{1.6cm}p{1.8cm}p{4.8cm}@{}}
\toprule
\textbf{Floor} & \textbf{Value} & \textbf{Layer} & \textbf{Origin} \\
\midrule
$\epsEM = \epsS$ & $0.007494$ & bridge
  & $\epsuniv\cdot r^{-1/4}$; spatial-type (ZC70/81) \\
$\epscan = \epsTLS$ & $0.007550$ & spectral
  & spectral quadratic (ZC79/THM-ZC79-01) \\
$\epsuniv$ & $0.007599$ & bridge
  & WaveChain fixed point (ZC69/COR-ZC69-02) \\
$\epsW$ & $0.007705$ & bridge
  & $\epsuniv\cdot r^{+1/4}$; matter-type (ZC70) \\
\bottomrule
\end{tabular}}
\end{center}

\textbf{Why the ordering is non-trivial.}
The four floors come from three distinct sub-chains:
$\epscan$ from the TLS spectral quadratic (ZC79),
$\epsuniv$ from the WaveChain K14 fixed point (ZC69),
and $(\epsEM, \epsW)$ from the Born-chain sector weights applied to
$\epsuniv$ (ZC67/ZC70). The cross-layer inequality $\epsEM < \epscan$
mixes two sub-chains: it holds because the sector-weight suppression
($r^{-1/4} = 0.9862$) exceeds the two-layer fractional gap
($1 + 0.65\% = 1.00646$) in the right direction.

\textbf{Classification:} Near-Formal Theorem (Step~3 is numerical
at T-canonical; full algebraic generality is GAP-ZC83-02).
\end{formalbox}

%%=============================================================
\subsection{FIND-ZC83-01\quad Geometric-Mean Identity}
\label{sec:find01}
%%=============================================================

\needspace{10\baselineskip}

\begin{giftbox}[title={FIND-ZC83-01: $\sqrt{\epsEM\cdot\epsW}
  = \epsuniv$ Exactly ($R_0$, Exact --- Gift)}]

\textbf{Statement.}
\begin{equation}
  \boxed{
    \sqrt{\epsEM\cdot\epsW} \;=\; \epsuniv \quad\text{(exact).}
  }
  \label{eq:gm}
\end{equation}
The bridge floor is the geometric mean of the two sector floors.

\textbf{Proof.}
From LEM-ZC83-01, equations \eqref{eq:floor-formula}:
\begin{equation}
  \sqrt{\epsEM\cdot\epsW}
  = \sqrt{(\epsuniv\cdot r^{-1/4})\cdot(\epsuniv\cdot r^{+1/4})}
  = \sqrt{\epsuniv^{2}\cdot r^{\,0}}
  = \epsuniv.
  \label{eq:gm-proof}
\end{equation}
The sector-weight exponents $-1/4$ and $+1/4$ cancel exactly.
No approximation is involved. \hfill$\square$

\textbf{Why this was not obvious.}
The five floor values span three distinct derivation sub-chains.
That $\epsuniv$ --- derived from the WaveChain K14 fixed point (ZC69)
--- equals the geometric mean of $\epsEM$ and $\epsW$ --- derived
from the Born-chain sector weights (ZC67/ZC70) --- is a non-trivial
cross-chain identity. Its proof requires only the cancellation of
exponents in LEM-ZC67-01, but this cancellation was never exhibited
before ZC83.

\textbf{Relation to COR-ZC67-01.}
The headline value of COR-ZC67-01 in ZC67 was stated as
$\sqrt{\epsEM^{\rm def}\cdot\epsW^{\rm def}} = \epscan$ (using
the ZC66/67 defined input pair). FIND-ZC83-01 uses the
ZC70-exact input pair $(\epsEM, \epsW) = \epsuniv(r^{-1/4},
r^{+1/4})$ and recovers $\epsuniv$ exactly. These are different
evaluations of the same algebraic identity on different input pairs
(cf.\ HANDOFF v17.14, Section~0).

\textbf{T-canonical:} $\sqrt{0.007494226 \times 0.007704704}
= 0.007598736 = \epsuniv$. Gap: $0.000\%$ (machine precision).

\textbf{Classification:} Exact Finding (Gift).
\physmeaning{the universal bridge floor is not an average chosen to fit
  the data --- it is the exact geometric mean of the two sector floors it
  sits between, forced by the same $\pm\tfrac14$ exponent structure that
  orders them. The middle floor is locked to the outer two.}
\end{giftbox}

%%=============================================================
\subsection{FIND-ZC83-02\quad Arithmetic-Mean Inequality}
\label{sec:find02}
%%=============================================================

\needspace{8\baselineskip}

\begin{derivedbox}[title={FIND-ZC83-02: $\AMfloor(\epsEM,\epsW)
  > \epsuniv$ ($R_0$, Exact)}]

\textbf{Statement.}
\begin{equation}
  \frac{\epsEM + \epsW}{2}
  \;>\; \epsuniv
  \;=\; \sqrt{\epsEM\cdot\epsW},
  \label{eq:am-gm}
\end{equation}
with equality iff $\epsEM = \epsW$ (which would require $r = 1$,
i.e.\ $x = 0$, which is excluded).

\textbf{Proof.}
The AM--GM inequality gives $(\epsEM + \epsW)/2 \geq
\sqrt{\epsEM\cdot\epsW}$ with strict inequality when
$\epsEM \neq \epsW$. Since $r > 1$ (LEM-ZC83-01), we have
$\epsEM \neq \epsW$, and FIND-ZC83-01 gives
$\sqrt{\epsEM\cdot\epsW} = \epsuniv$, hence the result.

The exact gap has closed form:
\begin{equation}
  \AMfloor - \GMfloor
  = \epsuniv\!\left(\frac{r^{1/4}+r^{-1/4}}{2} - 1\right)
  = \epsuniv\!\left(\cosh\!\tfrac{\ln r}{4} - 1\right),
  \label{eq:am-gm-closed}
\end{equation}
evaluated at T-canonical: $7.29\times10^{-7}$ ($+0.0096\%$ above
$\epsuniv$).

\textbf{T-canonical:}
\begin{center}
{\small\renewcommand{\arraystretch}{1.25}
\begin{tabular}{@{}lrl@{}}
\toprule
Quantity & Value & \\
\midrule
$\AMfloor = (\epsEM+\epsW)/2$ & $0.007599465$ & \\
$\GMfloor = \sqrt{\epsEM\cdot\epsW}$ & $0.007598736$ & $= \epsuniv$ \\
$\AMfloor - \GMfloor$         & $7.29\times10^{-7}$ & \\
$\AMfloor / \epsuniv - 1$     & $+0.0096\%$ & \\
\bottomrule
\end{tabular}}
\end{center}

\textbf{Classification:} Exact Finding.
\end{derivedbox}

%%=============================================================
\subsection{Phase Misalignment and Open Gap}
\label{sec:pm}
%%=============================================================

\begin{openqbox}[title={PM-ZC83-01: $\epsSC \neq \AMfloor(\epsEM,\epsW)$
  --- Probe Refuted
  (pattern: tautology test)}]

\textbf{Probe.}
After establishing FIND-ZC83-01 (GM) and FIND-ZC83-02 (AM), the
natural question was: does $\epsSC = 0.007605$ (the observed
bracket of ZC68) equal the AM of the floor pair?

\textbf{Result.}
\[
  \AMfloor(\epsEM,\epsW) = 0.007599465
  \neq 0.007605 = \epsSC.
\]
Gap: $\epsSC - \AMfloor = 5.54\times10^{-6}$ ($+0.073\%$).
The observed bracket $\epsSC$ lies strictly above both the GM
($= \epsuniv$, $+0.080\%$) and the AM ($+0.073\%$ above AM).

\textbf{Consequence.}
$\epsSC$ is not the AM, GM, or any simple algebraic function of
the defined floor pair $(\epsEM, \epsW)$ under the ZC67/70
sector-weight scheme. Its ontology (what invariant of the
$\mathbb{Z}_{15}$ structure it probes) remains open.

\textbf{Status:} Scar. Opens GAP-ZC83-01.
\end{openqbox}

\begin{openqbox}[title={GAP-ZC83-01: $\epsSC$ Ontology and Closed
  Form (Priority 1)}]

\textbf{Statement.}
Find a closed-form expression for $\epsSC = 0.007605$ or prove
it is transcendental over the known CRF floor extensions.

\textbf{Known constraints (from PM-ZC83-01):}
\begin{itemize}[leftmargin=2em, noitemsep]
  \item $\epsSC = \sqrt{\epsEM^{\rm obs}\cdot\epsW^{\rm obs}}$
    (observed bracket; ZC68/THM-ZC68-01b).
  \item $\epsSC > \AMfloor(\epsEM,\epsW) = 0.007599465$.
  \item $\epsSC > \epsuniv = \GMfloor(\epsEM,\epsW) = 0.007598736$.
  \item Offset from $\epsuniv$: $+0.08\%$ --- the two-layer separation
    of ZC71 plus an additional contribution.
\end{itemize}

\textbf{Approach candidates:}
\begin{enumerate}[leftmargin=2em, noitemsep]
  \item Express $\epsSC$ as a function of the observed
    $\delta\alpha/\alpha$ and $\delta(\sin^2\theta_W)$ via
    LEM-ZC68-01 fixed-point equation.
  \item Check whether $\epsSC \in \mathbb{Q}(\sqrt{5})\setminus
    \mathbb{Q}$ (same field extension as $\epsuniv$, THM-ZC71-01)
    or requires a different extension.
  \item If no closed form exists: register $\epsSC$ as
    ``observationally bounded'' and state the bounding theorem.
\end{enumerate}

\textbf{Status:} Open. Target: ZC84 or later.
\end{openqbox}

\begin{openqbox}[title={GAP-ZC83-02: Algebraic Generality of
  LEM-ZC83-02}]

\textbf{Statement.}
LEM-ZC83-02 proves $\epsEM < \epscan$ for $x = 0.0604$ (T-canonical)
by numerical comparison. A full algebraic proof that
$\epsuniv/\epscan < r^{1/4}$ holds for \emph{all} $x$ in the
physical range (e.g.\ $x \in (0, 0.1)$) would promote LEM-ZC83-02
to Exact and thereby promote THM-ZC83-01 to Exact.

\textbf{Approach:} Expand $r^{1/4} = ((1+2x)/(1+x))^{1/4}$ and
$\epsuniv/\epscan$ as functions of $x$; prove the inequality holds
for all $x > 0$ (or all $x$ in the relevant range).

\textbf{Status:} Open (low priority; T-canonical is sufficient
for CRF).
\end{openqbox}

%%=============================================================
\subsection{Complete Floor Lattice: Unified Picture}
\label{sec:lattice}
%%=============================================================

\begin{keybox}[title={Complete Floor Lattice (T-canonical, $R_0$)
  --- Post-ZC83 Unified Picture}]

\begin{center}
{\small\renewcommand{\arraystretch}{1.35}
\begin{tabular}{@{}p{2.8cm}p{1.6cm}p{1.8cm}p{5.3cm}@{}}
\toprule
\textbf{Floor} & \textbf{Value} & \textbf{Layer} & \textbf{Derivation} \\
\midrule
$\epsEM = \epsS$ & $0.007494$ & bridge
  & $\epsuniv\cdot r^{-1/4}$ (ZC67/70/81) \\
$\epscan = \epsTLS$ & $0.007550$ & spectral
  & quadratic root (ZC79) \\
$\epsuniv$       & $0.007599$ & bridge
  & K14 fixed point (ZC69) \\
$\epsW$          & $0.007705$ & bridge
  & $\epsuniv\cdot r^{+1/4}$ (ZC67/70) \\
\midrule
$\AMfloor$       & $0.007599$ & ---
  & $(\epsEM+\epsW)/2 > \epsuniv$ by $7.3\times10^{-7}$ \\
$\epsSC$         & $0.007605$ & observed
  & ZC68 bracket; $>\AMfloor$ by $5.5\times10^{-6}$ (GAP-ZC83-01) \\
\bottomrule
\end{tabular}}
\end{center}

\medskip
\textbf{Exact identities (ZC83 additions):}
\begin{align}
  \sqrt{\epsEM\cdot\epsW} &= \epsuniv
    && \text{(FIND-ZC83-01, exact)} \label{eq:gm-id} \\
  \frac{\epsEM+\epsW}{2}   &> \epsuniv
    && \text{(FIND-ZC83-02, exact; gap $7.29\times10^{-7}$)}
    \label{eq:am-id}
\end{align}

\textbf{Complete ordering (THM-ZC83-01):}
\[
  \epsEM = \epsS \;<\; \epscan \;<\; \epsuniv \;<\; \epsW.
\]

\textbf{Physical reading.}
The four derived floors span two layers and three sub-chains:
\begin{itemize}[leftmargin=2em, noitemsep]
  \item \emph{Spectral layer} ($\epscan$): set by the TLS quadratic
    at $R_0$ (ZC79). This is the lowest bridge-adjacent value.
  \item \emph{Bridge layer} ($\epsEM$, $\epsuniv$, $\epsW$):
    set by the Born-chain sector weights. $\epsuniv$ sits at the
    geometric center of the sector pair.
  \item The two layers are separated by the irreducible $0.66\%$
    gap of THM-ZC71-01.
\end{itemize}
The ordering $\epsEM < \epscan$ is the non-trivial statement: it
says the lower sector floor crosses below the spectral layer boundary.
\end{keybox}

%%=============================================================
\subsection{Atomic Registry}
\label{sec:registry}
%%=============================================================

\begin{formalbox}[title={Atomic Registry --- ZC83}]
\begin{center}
\renewcommand{\arraystretch}{1.3}
\small
\begin{tabular}{lp{7.8cm}l}
\toprule
\textbf{ID} & \textbf{Description} & \textbf{Status} \\
\midrule
LEM-ZC83-01
  & $\epsEM < \epsuniv < \epsW$ forced by $r > 1$ ($x > 0$);
    two-step proof from LEM-ZC67-01
  & Exact \\[3pt]
LEM-ZC83-02
  & $\epsEM < \epscan$; condition $\epsuniv/\epscan < r^{1/4}$;
    holds at T-canonical ($1.00646 < 1.01395$)
  & Near-Formal \\[3pt]
THM-ZC83-01
  & Complete ordering $\epsEM = \epsS < \epscan < \epsuniv < \epsW$;
    three-step proof from LEM-ZC83-01/02 + THM-ZC71-01 + THM-ZC81-01
  & Near-Formal \\[3pt]
FIND-ZC83-01
  & $\sqrt{\epsEM\cdot\epsW} = \epsuniv$ exactly;
    exponent cancellation in LEM-ZC67-01 (Gift)
  & Exact \\[3pt]
FIND-ZC83-02
  & $\AMfloor(\epsEM,\epsW) > \epsuniv$; strict AM--GM;
    gap $7.29\times10^{-7}$; closed form eq.~\eqref{eq:am-gm-closed}
  & Exact \\[3pt]
PM-ZC83-01
  & $\epsSC \neq \AMfloor(\epsEM,\epsW)$;
    probe refuted ($+0.073\%$ mismatch)
    (pattern: tautology test)
  & Scar \\[3pt]
GAP-ZC83-01
  & $\epsSC = 0.007605$ ontology; not AM or GM of floor pair;
    closed form or field classification open; target ZC84
  & Open \\[3pt]
GAP-ZC83-02
  & LEM-ZC83-02 algebraic generality; promote to Exact
    when $\epsuniv/\epscan < r^{1/4}$ proved for all $x > 0$
  & Open \\[3pt]
\bottomrule
\end{tabular}
\end{center}
\end{formalbox}

%%=============================================================
\subsection{Master Status Table}
\label{sec:status}
%%=============================================================

{\small
\setlength{\LTpre}{4pt}\setlength{\LTpost}{4pt}
\renewcommand{\arraystretch}{1.30}
\begin{longtable}{>{\raggedright\arraybackslash}p{5.0cm}
                  >{\centering\arraybackslash}p{2.4cm}
                  >{\raggedright\arraybackslash}p{5.2cm}}
\toprule
\textbf{Item} & \textbf{Status} & \textbf{Notes} \\
\midrule
\endfirsthead
\multicolumn{3}{l}{\small\textit{(continued)}}\\[4pt]
\toprule\textbf{Item}&\textbf{Status}&\textbf{Notes}\\\midrule
\endhead
\midrule\multicolumn{3}{r}{\small\textit{(continues)}}\\
\endfoot
\bottomrule
\endlastfoot
%%
LEM-ZC83-01 (sector-weight ordering)
  & Exact
  & $\epsEM < \epsuniv < \epsW$; from $r>1$ iff $x>0$ \\[4pt]
LEM-ZC83-02 (cross-layer ordering)
  & Near-Formal
  & $\epsEM < \epscan$; holds at T-canonical;
    GAP-ZC83-02 for algebraic generality \\[4pt]
THM-ZC83-01 (complete ordering)
  & Near-Formal
  & $\epsEM=\epsS<\epscan<\epsuniv<\epsW$;
    three-step; inherits LEM-ZC83-02 near-formal \\[4pt]
FIND-ZC83-01 (GM identity)
  & Exact (Gift)
  & $\sqrt{\epsEM\cdot\epsW}=\epsuniv$ exactly;
    cross-chain algebraic identity \\[4pt]
FIND-ZC83-02 (AM inequality)
  & Exact
  & $\AMfloor > \epsuniv$; strict; gap $7.29\times10^{-7}$ \\[4pt]
PM-ZC83-01 ($\epsSC \neq \AMfloor$)
  & Scar
  & tautology test; probe refuted; opens GAP-ZC83-01 \\[4pt]
GAP-ZC83-01 ($\epsSC$ ontology)
  & Open
  & neither AM nor GM of defined pair; target ZC84 \\[4pt]
GAP-ZC83-02 (algebraic generality)
  & Open
  & LEM-ZC83-02 for all $x>0$; low priority \\[4pt]
\midrule
Floor lattice ordering (no prior theorem)
  & Closed $\checkmark$
  & via THM-ZC83-01; first algebraic proof \\
\bottomrule
\end{longtable}
}

%%=============================================================
\subsection{Canonical $\varepsilon$ Reference Table}
\label{sec:eps-table}
%%=============================================================

\begin{keybox}[title={Canonical $\varepsilon$ Ladder ---
  v17.14 + ZC83 (FIND-ZC83-01 addition)}]
\sloppy
T-canonical: $n=8$, $\ds=3$, $\nm=5$,
$\Ks = 1-e^{-1/8} = 0.117503$.

{\small\renewcommand{\arraystretch}{1.35}
\begin{center}
\begin{tabular}{@{}p{3.5cm}p{1.5cm}p{1.5cm}p{5.5cm}@{}}
\toprule
\textbf{Symbol} & \textbf{Value} & \textbf{Layer} & \textbf{Source} \\
\midrule
$\epsEM = \epsS$ & $0.007494$ & bridge
  & ZC70/81; $\epsuniv\cdot r^{-1/4}$ \\[3pt]
$\epscan = \epsTLS$ & $0.007550$ & spectral
  & ZC79/THM-ZC79-01; exact \\[3pt]
$\epsuniv$ & $0.007599$ & bridge
  & ZC69/COR-ZC69-02; exact \\[3pt]
$\epsW$ & $0.007705$ & bridge
  & ZC70/THM-ZC70-01; $\epsuniv\cdot r^{+1/4}$ \\[4pt]
\midrule
\multicolumn{4}{@{}p{12.5cm}}{\textit{Ordering (THM-ZC83-01):}
  $\epsEM=\epsS < \epscan < \epsuniv < \epsW$}\\[3pt]
\multicolumn{4}{@{}p{12.5cm}}{\textbf{NEW (FIND-ZC83-01):}
  $\sqrt{\epsEM\cdot\epsW} = \epsuniv$ (exact;
  bridge floor = GM of sector floors)}\\[3pt]
\multicolumn{4}{@{}p{12.5cm}}{\textit{AM--GM (FIND-ZC83-02):}
  $(\epsEM+\epsW)/2 = 0.007599465 > \epsuniv$ by $7.29\times10^{-7}$}\\[3pt]
\multicolumn{4}{@{}p{12.5cm}}{\textit{Two-layer separation (ZC71):}
  $\epsuniv - \epscan = 0.66\%$; irreducible}\\
\bottomrule
\end{tabular}
\end{center}}
\end{keybox}

%%=============================================================
\subsection{R-Layer Research Note}
\label{sec:rlayer}
%%=============================================================

\begin{layerbox}[title={R-Layer Assignments for ZC83}]

{\small\renewcommand{\arraystretch}{1.25}
\begin{center}
\begin{tabular}{@{}llll@{}}
\toprule
Atomic unit & R-layer & Cost $T$ & Source \\
\midrule
LEM-ZC83-01 & $R_0$ & algebraic & ZC67 + ZC57 \\
LEM-ZC83-02 & $R_0$ & algebraic & ZC67 + ZC71 + ZC79 \\
THM-ZC83-01 & $R_0$ & algebraic & LEM-ZC83-01/02 + ZC71 + ZC81 \\
FIND-ZC83-01 & $R_0$ & $0$ & algebraic; exponent cancellation \\
FIND-ZC83-02 & $R_0$ & $0$ & AM--GM inequality \\
GAP-ZC83-01 & $R_0$ & TBD & ZC68 observed bracket \\
\bottomrule
\end{tabular}
\end{center}}

\textbf{Two-layer floor (PWA-8):}
\begin{itemize}[leftmargin=2em, noitemsep]
  \item Spectral layer ($\epscan = 0.007550$): ZC79. Used only
    as the lower boundary in LEM-ZC83-02. Not modified.
  \item Bridge layer ($\epsuniv = 0.007599$): ZC69. The exact
    value appears in FIND-ZC83-01 and FIND-ZC83-02. Not modified.
  \item Separation $0.66\%$ irreducible (THM-ZC71-01): inherited.
\end{itemize}

All results are $R_0$; no bridge crossing formula enters.
The two-layer structure is used structurally (THM-ZC71-01 supplies
$\epscan < \epsuniv$ for THM-ZC83-01 Step~2) but not modified.
\end{layerbox}

%%=============================================================
\subsection{Genealogy Map}
\label{sec:genealogy}
%%=============================================================

\begin{genealogybox}[title={How ZC83's Key Objects Trace Back}]

\textbf{Branch A --- Sector-floor formula chain:}
\begin{itemize}[leftmargin=2em, noitemsep]
  \item \textbf{ZC57/THM-ZC57-01}: $C_{\rm EM} = 1+x$,
    $C_W = 1+2x$; Born-chain weights from $\Zth\times\Zfi$.
  \item $\to$ \textbf{ZC67/LEM-ZC67-01}: sector-floor formula
    $\varepsilon_i = \epsuniv(C_i/C_{\rm ref})^{1/4}$.
  \item $\to$ \textbf{ZC70/THM-ZC70-01}: $\epsEM$, $\epsW$ exact.
  \item $\to$ \textbf{ZC83/LEM-ZC83-01}: ordering from $r > 1$.
  \item $\to$ \textbf{ZC83/FIND-ZC83-01}: GM identity
    (exponent cancellation).
\end{itemize}

\textbf{Branch B --- Two-layer spectral chain:}
\begin{itemize}[leftmargin=2em, noitemsep]
  \item \textbf{ZC71/THM-ZC71-01}: $\epscan < \epsuniv$;
    irreducible $0.66\%$ gap.
  \item \textbf{ZC79/THM-ZC79-01}: $\epscan$ exact from
    spectral quadratic; $x = n\epscan$ exact.
  \item $\to$ \textbf{ZC83/LEM-ZC83-02}: $\epsEM < \epscan$
    (cross-layer inequality).
  \item $\to$ \textbf{ZC83/THM-ZC83-01}: Step~2 and Step~3
    of complete ordering.
\end{itemize}

\textbf{Branch C --- Bridge fixed-point chain:}
\begin{itemize}[leftmargin=2em, noitemsep]
  \item \textbf{ZC69/COR-ZC69-02}: $\epsuniv$ exact closed form.
  \item $\to$ \textbf{ZC83/FIND-ZC83-01}: $\epsuniv$ identified
    as exact GM of floor pair.
  \item $\to$ \textbf{ZC83/FIND-ZC83-02}: AM--GM strict inequality.
\end{itemize}

\textbf{Branch D --- Strong floor chain:}
\begin{itemize}[leftmargin=2em, noitemsep]
  \item \textbf{ZC81/THM-ZC81-01}: $\epsS = \epsEM$
    (spatial-type class).
  \item $\to$ \textbf{ZC83/THM-ZC83-01}: Step~4; Strong sector
    included in ordering without new proof.
\end{itemize}

\textbf{Convergence at ZC83:}
All four branches converge in THM-ZC83-01. Branches~A and~B
supply the two independent inequalities that sandwich $\epscan$
between $\epsEM$ and $\epsuniv$. Branch~C supplies the GM identity.
Branch~D extends the ordering to the Strong sector for free.

ZC83's contribution: the first paper to prove the complete floor
ordering from first principles, and to identify $\epsuniv$ as the
exact geometric mean of the sector floor pair.
\end{genealogybox}

%%=============================================================
\subsection{Closing Sentence}
%%=============================================================
%% TONE A -- TRIUMPHANT

\bigskip
\begin{center}
\textit{%
  The five instability floors were always ordered.\\[0.3em]
  The ordering was forced by the sector-weight arithmetic
  of $\mathbb{Z}_{15}$\\
  and the irreducible two-layer gap ---\\
  waiting only to be stated.\\[0.5em]
  $\displaystyle
    \sqrt{\varepsilon_{\rm EM}\cdot\varepsilon_W}
    \;=\;
    \varepsilon_{\rm univ}
    \qquad\text{(exactly).}$
}
\end{center}

% Governance Note: This document follows Upgrade Protocol v1.2
% and CRF Style Guide v3.5.


\newpage
%%─── VERBATIM EMBED: ZC84 ───────────────────────────────────────────────
\addcontentsline{toc}{section}{ZC84: The Complete Sector Synthesis}

% [BRIDGE-PROSE: integration-added, Extension Freedom narrative, v3.6 §31.3]
A theory built one result at a time accumulates a hidden risk: each piece
may be sound on its own while the whole quietly fails to cohere. By the
time CRF reached this point it had produced a dozen separate results
across four force sectors --- floors, couplings, mirror relations, a
cosmological term --- each proved in its own paper, each checked against
its own data. What had never been done was to lay them side by side and
ask whether they tell one story or several.

ZC84 is that reckoning. It introduces nothing new; that is the point. It
takes the eight standing consequences of $\delta$ and writes them as a
single object, so that any reader --- or any critic --- can see at one
glance what the framework commits to and where a single wrong number
would bring the rest down with it. This is the moment a collection of
findings is asked to behave like a theory.

The honest content of this paper is therefore its scope, not its novelty,
and the sections below say so plainly: a synthesis is only as strong as
the weakest result it inherits, and ZC84 puts every one of them on the
table at once, alongside a falsification record that names the conditions
under which the whole assembly would fail.



%%=============================================================
\begin{abstract}
\sloppy
ZC80 through ZC83 completed four independent arcs: the electroweak
completeness theorem, the Strong sector floor, the cosmological
ladder identity, and the floor lattice ordering. No single paper
had assembled them into one picture.

\textbf{THM-ZC84-01} collects eight simultaneous consequences of
the single primitive $\delta$ and the Key Identity $3d_s = 2^{d_s}
+ 1$, with zero free parameters: the Weinberg angle, the EM and
gravitational fine structure constants, the complete floor lattice
ordering, the geometric-mean identity, the ladder mirror symmetry,
the cosmological lock, and the uniqueness of $\mathbb{Z}_{15}$.

\textbf{FIND-ZC84-01} displays the complete $\delta \to \Gamma_{\rm CRF}
\to \{\text{all sectors}\}$ chain in one diagram.

\textbf{COR-ZC84-01} assembles the four-sector falsification table.

No new derivations are performed. ZC84 is a synthesis paper; its
content is the assembly.
\end{abstract}


%%=============================================================
\subsection{Background and Scope}
\label{sec:background}
%%=============================================================

\begin{keybox}[title={Numerical Summary --- ZC84 (T-canonical, $R_0$)}]
{\small\renewcommand{\arraystretch}{1.25}
\begin{center}
\begin{tabular}{@{}p{4.8cm}p{2.2cm}p{1.5cm}p{3.5cm}@{}}
\toprule
\textbf{Observable} & \textbf{CRF} & \textbf{Gap} & \textbf{Source} \\
\midrule
$\swsq = 1 - \ln 8/\ln 15$    & $0.232126$ & $0.39\%$ & ZC20/THM-ZC20-01 \\
$\aEM(R_0) = (n/15)\GammaCRF$ & $1/137.22$ & $0.13\%$ & ZC29/ZC64 \\
$\agrav(R_0) = \GammaCRF/15$  & $0.000911$ & exact     & ZC72/THM-ZC72-01 \\
$\aEM/\agrav = n = 8$          & $8.000000$ & $0.000\%$ & ZC72/ZC75 \\
\midrule
$\epsEM = \epsS$              & $0.007494$ & exact     & ZC70/ZC81 \\
$\epscan = \epsTLS$           & $0.007550$ & exact     & ZC79 \\
$\epsuniv$                    & $0.007599$ & exact     & ZC69 \\
$\epsW$                       & $0.007705$ & exact     & ZC70 \\
\midrule
$\sqrt{\epsEM\cdot\epsW}$     & $0.007599$ & $0.000\%$ & ZC83/FIND-ZC83-01 \\
$\GammaEM + \GammaS - 2\GammaW$ & $0$      & $0.000\%$ & ZC81/FIND-ZC81-01 \\
$\wDE + 3\DeltaGam - \GammaCRF$ & $0$      & $0.000\%$ & ZC82/THM-ZC82-01 \\
\bottomrule
\end{tabular}
\end{center}}
\end{keybox}

\subsubsection{The Assembly Motivation}

ZC72 through ZC83 produced twelve papers across four arcs:
\begin{itemize}[leftmargin=2em, noitemsep]
  \item \textbf{Gravity arc} (ZC72--75): vacuum sector identification,
    $\Gamma$-ladder, $\mathbb{Z}_{15}$ uniqueness, gravity dimensional bridge.
  \item \textbf{Floor arc} (ZC79, ZC83): exact $\epscan$ from the
    TLS quadratic; complete ordering theorem with GM identity.
  \item \textbf{Electroweak completeness} (ZC80--81): seven $R_0$-layer
    EW observables; Strong sector floor and ladder mirror symmetry.
  \item \textbf{Cosmological lock} (ZC82): $\wDE + 3\DeltaGam =
    \GammaCRF$ connecting the cosmological and electroweak arcs.
\end{itemize}

No single paper had placed all four arcs in one derivation chain.
ZC84 supplies that assembly without introducing any new derivation.

\begin{keybox}[title={What ZC84 Does NOT Do}]
\begin{itemize}[leftmargin=2em, noitemsep]
  \item Does not re-derive any pillar theorem; all results used verbatim.
  \item Does not open new gaps beyond those already registered in ZC80--83.
  \item Does not close GAP-ZC83-01 ($\epsSC$ ontology),
    GAP-ZC81-01 (Strong bridge), or GAP-ZC76-01 ($\tau_{\rm sweep}$).
  \item Does not claim any result beyond what the pillar papers state.
\end{itemize}
\end{keybox}

\begin{keybox}[title={Companion Papers --- ZC84}]
{\small\renewcommand{\arraystretch}{1.15}
\begin{center}
\begin{tabular}{@{}p{3.8cm}p{8.5cm}@{}}
\toprule
\textbf{Paper} & \textbf{Pillar provided} \\
\midrule
ZC20/THM-ZC20-01 & $\swsq = 1 - \ln 8/\ln 15$ \\
ZC29/THM-ZC29-01 & $\alpha_i = \varphi(d_i)/15\cdot\GammaCRF$ \\
ZC57/THM-ZC57-01 & $C_{\rm EM} = 1+x$, $C_W = 1+2x$ \\
ZC69/COR-ZC69-02 & $\epsuniv$ exact; K14 exact \\
ZC70/THM-ZC70-01 & $\epsEM$, $\epsW$ exact \\
ZC71/THM-ZC71-01 & two-layer gap irreducible \\
ZC72/THM-ZC72-01 & $\agrav = \GammaCRF/15$; vacuum sector \\
ZC74/THM-ZC74-01 & $\mathbb{Z}_{15}$ unique squarefree semiprime \\
ZC75/THM-ZC75-01 & $\agrav\cdot|\Ztf| = \GammaCRF$ \\
ZC79/THM-ZC79-01 & $\epscan$ exact from quadratic \\
ZC80/THM-ZC80-01 & 7 EW observables at $R_0$, zero free params \\
ZC81/THM-ZC81-01, FIND-ZC81-01 & $\epsS = \epsEM$; mirror symmetry \\
ZC82/THM-ZC82-01 & $\wDE + 3\DeltaGam = \GammaCRF$ \\
ZC83/THM-ZC83-01, FIND-ZC83-01 & ordering; $\sqrt{\epsEM\epsW}=\epsuniv$ \\
\bottomrule
\end{tabular}
\end{center}}
\end{keybox}

\subsubsection{Pre-Work Audit (PWA-1 through PWA-10)}

\begin{formalbox}[title={Pre-Work Audit Record --- ZC84}]
Scanned before writing (Upgrade Protocol v1.2):
\begin{itemize}[leftmargin=2em, noitemsep]
  \item All pillar papers ZC20--ZC83 scanned via Index v17.14.
  \item No duplicate atomic unit introduced; ZC84 introduces only
    THM-ZC84-01, COR-ZC84-01, FIND-ZC84-01, and PM-ZC84-SCOPE-01.
\end{itemize}
\textbf{PWA-5} (tautology): THM-ZC84-01 is a collection theorem,
not an independent derivation. Registered as near-formal; each
clause inherits the status of its source paper.\\
\textbf{PWA-6} (macro): all new macros via \verb|\providecommand|;
no conflicts with v17.14 macro chain.\\
\textbf{PWA-7} (sector-floor): ordering clause (iv) inherits
THM-ZC83-01; no new sector assignment.\\
\textbf{PWA-8} (layer): two-layer separation clause inherited from
ZC71/ZC83; not modified.\\
\textbf{No duplicate atomic units found.}
\end{formalbox}

\newpage

%%=============================================================
\subsection{THM-ZC84-01\quad The Complete Sector Synthesis}
\label{sec:thm01}
%%=============================================================

\needspace{12\baselineskip}

\begin{formalbox}[title={THM-ZC84-01: Complete Sector Synthesis
  ($R_0$, Near-Formal)}]

\textbf{Statement.}
At T-canonical gauge ($d_s = 3$ exact, $n_m = 5$ exact, derived
solely from the Key Identity $3d_s = 2^{d_s}+1$), the following
eight results hold simultaneously with zero free parameters:

\medskip
\begin{enumerate}[leftmargin=2.5em, label=(\roman*)]

\item \textbf{Weinberg angle (ZC20, Exact):}
\begin{equation}
  \swsq = 1 - \frac{\ln 8}{\ln 15} = 0.232126.
  \label{eq:sw}
\end{equation}

\item \textbf{EM fine structure (ZC29/ZC64, Near-Formal):}
\begin{equation}
  \aEM(R_0) = \frac{n}{15}\,\GammaCRF
  = \frac{\varphi(15)}{15}\!\left(\frac{\Gcoup}{d_s} - \swsq\right).
  \label{eq:alpha-em}
\end{equation}
Gap vs CODATA: $0.13\%$ at $R_0$.

\item \textbf{Gravitational fine structure (ZC72/ZC75, Exact):}
\begin{equation}
  \agrav(R_0) = \frac{\GammaCRF}{15},
  \qquad
  \frac{\aEM}{\agrav} = n = 8,
  \qquad
  \agrav\cdot|\Ztf| = \GammaCRF.
  \label{eq:alpha-grav}
\end{equation}

\item \textbf{Floor lattice ordering (ZC83, Near-Formal):}
\begin{equation}
  \epsEM = \epsS \;<\; \epscan \;<\; \epsuniv \;<\; \epsW.
  \label{eq:ordering}
\end{equation}

\item \textbf{Geometric-mean identity (ZC83, Exact):}
\begin{equation}
  \sqrt{\epsEM\cdot\epsW} = \epsuniv.
  \label{eq:gm}
\end{equation}
The bridge floor is the exact geometric mean of the sector floors.

\item \textbf{Ladder mirror symmetry (ZC81, Exact):}
\begin{equation}
  \GammaEM + \GammaS = 2\,\GammaW.
  \label{eq:mirror}
\end{equation}
Gap: $0.0000000000\%$.

\item \textbf{Cosmological lock (ZC82, Exact):}
\begin{equation}
  \wDE + 3\,\DeltaGam = \GammaCRF.
  \label{eq:cosmo}
\end{equation}
Gap: $0.000\%$.

\item \textbf{Group uniqueness (ZC74, Exact):}
\begin{equation}
  \mathbb{Z}_{15} = \mathbb{Z}_3 \times \mathbb{Z}_5
  \;\text{ is the unique such group with }
  \varphi\text{-set } \{1,2,4,8\}.
  \label{eq:unique}
\end{equation}
The binary ladder is exclusive. All cascades above are forced.

\end{enumerate}

\textbf{Proof.}
Each clause is proved in its source paper (cited in the statement).
The collection is consistent: all eight share the same T-canonical
cascade $\delta \to d_s = 3 \to |\Ztf| = 15 \to \varphi(15) = n = 8
\to \Gcoup \to \GammaCRF$. No parameter is introduced twice or
inconsistently. \hfill$\square$

\textbf{Classification:} Near-Formal Theorem (inherits the weakest
status of its clauses; clauses (i), (iii), (v)--(viii) are Exact;
clauses (ii) and (iv) are Near-Formal from their source papers).
\end{formalbox}

%%=============================================================
\subsection{FIND-ZC84-01\quad Complete Derivation Chain}
\label{sec:find01}
%%=============================================================

\begin{keybox}[title={FIND-ZC84-01: $\delta \to \GammaCRF \to
  \{\text{all sectors}\}$ --- Complete Chain ($R_0$, Exact)}]

\textbf{Foundation layer:}
\[
  \delta
  \;\xrightarrow{\text{Axiom 0}}\; \mathbb{Z}_{15}
  \;\xrightarrow{3d_s = 2^{d_s}+1}\; d_s = 3,\; n = 8
  \;\xrightarrow{\text{ZC74}}\; \text{unique binary ladder}
\]

\textbf{Coupling layer:}
\[
  d_s = 3,\; |\Ztf| = 15,\; \varphi(15) = 8
  \;\xrightarrow{\text{ZC20}}\; \swsq = 1 - \tfrac{\ln 8}{\ln 15}
  \;\xrightarrow{\text{ZC28}}\; \GammaCRF = \tfrac{\Gcoup}{d_s} - \swsq
\]

\textbf{Sector rate layer (ZC73/THM-ZC73-01):}
\[
  \Gamma_k = \frac{\Gcoup}{d_s} - 1 + k\DeltaGam,
  \quad k = 0,1,2,3,
  \quad \DeltaGam = \log_{15} 2
\]
\[
  \GammaVac \;<\; \GammaS \;<\; \GammaW \;<\; \GammaEM,
  \qquad
  \GammaEM + \GammaS = 2\GammaW \;\text{(ZC81, exact)}
\]

\textbf{Coupling constant layer (ZC29/THM-ZC29-01):}
\[
  \aEM = \tfrac{8}{15}\GammaCRF, \quad
  \alpha_W = \tfrac{4}{15}\GammaCRF, \quad
  \alpha_S = \tfrac{2}{15}\GammaCRF, \quad
  \agrav = \tfrac{1}{15}\GammaCRF
\]

\textbf{Floor lattice layer (ZC67/ZC70/ZC79/ZC83):}
\[
  \epsEM = \epsS
  \;<\; \epscan \;<\; \epsuniv \;<\; \epsW,
  \qquad
  \sqrt{\epsEM\cdot\epsW} = \epsuniv \;\text{(exact)}
\]

\textbf{Cosmological layer (ZC76/ZC82):}
\[
  \wDE = \frac{\Gcoup}{d_s} - 1,
  \qquad
  \wDE + 3\DeltaGam = \GammaCRF \;\text{(exact)}
\]

\textbf{Gravity dimensional bridge (ZC75/THM-ZC75-01):}
\[
  \agrav\cdot|\Ztf| = \GRzero = \frac{\ln 2}{D_0} = \GammaCRF,
  \qquad
  \frac{\aEM}{\agrav} = n = 8 \;\text{(exact)}
\]

\medskip
\textbf{Summary:} Every entry above traces back to $\delta$ through
the T-canonical cascade. The number of adjustable parameters is zero.
\end{keybox}

%%=============================================================
\subsection{COR-ZC84-01\quad Four-Sector Falsification Table}
\label{sec:cor01}
%%=============================================================

\begin{nobelbox}[title={COR-ZC84-01: Four-Sector Nobel-Table
  Falsification Signatures}]

Each row below constitutes a falsification signature (DEF-ZC64-02)
for the corresponding sector. A single $>5\sigma$ discrepancy in
\emph{any} row falsifies the indicated axiom chain.

\medskip
{\small\renewcommand{\arraystretch}{1.40}
\begin{center}
\begin{tabular}{@{}p{1.5cm}p{3.5cm}p{2.2cm}p{1.5cm}p{3.5cm}@{}}
\toprule
\textbf{Sector} & \textbf{CRF prediction} & \textbf{Obs.} & \textbf{Gap} & \textbf{Falsifies} \\
\midrule
EM
  & $\aEM(R_0) = (n/15)\GammaCRF$
  & $1/137.036$
  & $0.13\%$
  & THM-ZC29-01 + ZC20 \\[4pt]
EM+W
  & $\delta\aEM/\alpha : \delta(\swsq) = 1.4590$
  & (obs.\,ratio)
  & $0.007\%$
  & THM-ZC57-01 + ZC70 \\[4pt]
Strong
  & $\epsS = \epsEM = 0.007494$
  & (floor obs.)
  & exact
  & THM-ZC81-01 + ZC30 \\[4pt]
Grav
  & $\aEM/\agrav = n = 8$
  & (laboratory)
  & exact at $R_0$
  & THM-ZC72-01 + ZC75 \\[4pt]
Cosmo
  & $\wDE + 3\DeltaGam = \GammaCRF$
  & (CMB $w$)
  & exact
  & THM-ZC82-01 + ZC76 \\
\bottomrule
\end{tabular}
\end{center}}

\medskip
\textbf{Does NOT falsify} (if any single row fails):
the primitive $\delta$, Axioms 0--10, and all results not in the
indicated chain. CRF is modular; each sector chain is independently
falsifiable.

\textbf{Nobel-table criterion (DEF-ZC70/FIND-ZC70-02):}
The EM+W ratio row ($\delta\aEM/\alpha : \delta(\swsq) = 1.4590$)
tests a ratio of two independent experimental observables
simultaneously, is resistant to correlated systematics, and derives
from a zero-free-parameter chain. This row meets the Nobel-table
criterion of DEF-ZC64-02.
\end{nobelbox}

%%=============================================================
\subsection{Phase Misalignment}
\label{sec:pm}
%%=============================================================

\begin{openqbox}[title={PM-ZC84-SCOPE-01: Synthesis Content is
  Assembly, Not Derivation (pre-registered)
  (pattern: scope creep)}]

\textbf{Pre-registered stance.}
THM-ZC84-01 and FIND-ZC84-01 collect and display existing
theorems. They do not introduce new physics beyond what the
pillar papers state. In particular:
\begin{itemize}[leftmargin=2em, noitemsep]
  \item The eight clauses of THM-ZC84-01 are individually
    proved in their cited source papers.
  \item FIND-ZC84-01 is a display of the derivation chain, not
    a new step in the chain.
  \item COR-ZC84-01 collects existing falsification signatures
    (DEF-ZC64-02); it does not introduce new ones.
\end{itemize}

\textbf{Correct reading.}
ZC84's contribution is the \emph{demonstration of coherence}:
that all eight simultaneous results are consistent with the same
T-canonical cascade and carry no free parameters collectively.
This is valuable as a milestone record, not as a new derivation.
\end{openqbox}

%%=============================================================
\subsection{Complete Floor Lattice Reference}
\label{sec:lattice}
%%=============================================================

\begin{keybox}[title={Complete $\varepsilon$ Floor Lattice ---
  ZC84 Unified Reference}]

{\small\renewcommand{\arraystretch}{1.35}
\begin{center}
\begin{tabular}{@{}p{3.2cm}p{1.5cm}p{1.8cm}p{4.8cm}@{}}
\toprule
\textbf{Floor} & \textbf{Value} & \textbf{Layer} & \textbf{Source} \\
\midrule
$\epsEM = \epsS$ & $0.007494$ & bridge
  & ZC70/ZC81; $\epsuniv\cdot r^{-1/4}$; spatial-type \\
$\epscan = \epsTLS$ & $0.007550$ & spectral
  & ZC79/THM-ZC79-01; quadratic root \\
$\epsuniv$ & $0.007599$ & bridge
  & ZC69/COR-ZC69-02; K14 fixed point \\
$\epsW$ & $0.007705$ & bridge
  & ZC70/THM-ZC70-01; $\epsuniv\cdot r^{+1/4}$; matter-type \\
\midrule
$\sqrt{\epsEM\cdot\epsW}$ & $= \epsuniv$ & ---
  & ZC83/FIND-ZC83-01; exact \\
$(\epsEM+\epsW)/2$ & $0.007599$ & ---
  & ZC83/FIND-ZC83-02; $> \epsuniv$ by $7.3\times10^{-7}$ \\
\midrule
\multicolumn{4}{@{}p{12.5cm}}{\textit{Ordering (THM-ZC83-01):}
  $\epsEM = \epsS < \epscan < \epsuniv < \epsW$; forced by $x > 0$}\\[3pt]
\multicolumn{4}{@{}p{12.5cm}}{\textit{Two-layer separation (ZC71):}
  $\epsuniv - \epscan = 0.66\%$; irreducible;
  spectral layer vs bridge layer}\\
\bottomrule
\end{tabular}
\end{center}}
\end{keybox}

%%=============================================================
\subsection{Complete $\Gamma$-Ladder Reference}
\label{sec:ladder}
%%=============================================================

\begin{keybox}[title={Complete $\Gamma$-Ladder and Sector
  Summary (T-canonical, $R_0$)}]

From ZC73/THM-ZC73-01: $\Gamma_k = \Gcoup/d_s - 1 + k\DeltaGam$,
$\DeltaGam = \log_{15}2$.

\medskip
{\small\renewcommand{\arraystretch}{1.35}
\begin{center}
\begin{tabular}{@{}p{1.3cm}p{1.1cm}p{2.5cm}p{1.8cm}p{1.8cm}p{3.0cm}@{}}
\toprule
\textbf{Rung} & \textbf{Sector} & \textbf{$\Gamma_k$} &
\textbf{$\alpha_i(R_0)$} & \textbf{$\varepsilon$ floor} & \textbf{Type} \\
\midrule
$k=3$ & EM      & $+0.01366$ & $\tfrac{8}{15}\GammaCRF$
  & $\epsEM = 0.007494$ & spatial \\
$k=2$ & Weak    & $-0.24229$ & $\tfrac{4}{15}\GammaCRF$
  & $\epsW = 0.007705$  & matter \\
$k=1$ & Strong  & $-0.49825$ & $\tfrac{2}{15}\GammaCRF$
  & $\epsS = 0.007494$  & spatial \\
$k=0$ & Vac     & $-0.75421$ & $\tfrac{1}{15}\GammaCRF$
  & no floor          & vacuum \\
\midrule
\multicolumn{6}{@{}p{12.5cm}}{\textit{Mirror (ZC81):}
  $\GammaEM + \GammaS = 2\GammaW$ ($k=1$ and $k=3$ equidistant from $k=2$)}\\[2pt]
\multicolumn{6}{@{}p{12.5cm}}{\textit{Ladder bit (ZC82):}
  $\DeltaGam = \log_{15}2$; each rung adds one bit of distinction}\\[2pt]
\multicolumn{6}{@{}p{12.5cm}}{\textit{Cosmo lock (ZC82):}
  $\wDE + 3\DeltaGam = \GammaCRF$ (exact); cosmological arc meets EW arc}\\
\bottomrule
\end{tabular}
\end{center}}
\end{keybox}

%%=============================================================
\subsection{Atomic Registry}
\label{sec:registry}
%%=============================================================

\begin{formalbox}[title={Atomic Registry --- ZC84}]
\begin{center}
\renewcommand{\arraystretch}{1.3}
\small
\begin{tabular}{lp{8.0cm}l}
\toprule
\textbf{ID} & \textbf{Description} & \textbf{Status} \\
\midrule
THM-ZC84-01
  & 8-clause synthesis: $\swsq$, $\aEM$, $\agrav$, ordering,
    GM identity, mirror symmetry, cosmo lock, $\Ztf$ uniqueness;
    all from $\delta$, zero free params
  & Near-Formal \\[3pt]
FIND-ZC84-01
  & Complete $\delta \to \GammaCRF \to \{\text{all sectors}\}$
    chain; first paper to display all four sectors + cosmo arc
    in one diagram
  & Exact \\[3pt]
COR-ZC84-01
  & Four-sector Nobel-table falsification table;
    EM+W ratio row meets Nobel-table criterion (DEF-ZC64-02)
  & Near-Formal \\[3pt]
PM-ZC84-SCOPE-01
  & Synthesis content = assembly, not derivation;
    (pattern: scope creep; pre-registered)
  & Scar \\[3pt]
\bottomrule
\end{tabular}
\end{center}
\end{formalbox}

%%=============================================================
\subsection{Master Status Table}
\label{sec:status}
%%=============================================================

{\small
\setlength{\LTpre}{4pt}\setlength{\LTpost}{4pt}
\renewcommand{\arraystretch}{1.30}
\begin{longtable}{>{\raggedright\arraybackslash}p{5.0cm}
                  >{\centering\arraybackslash}p{2.4cm}
                  >{\raggedright\arraybackslash}p{5.2cm}}
\toprule
\textbf{Item} & \textbf{Status} & \textbf{Notes} \\
\midrule
\endfirsthead
\multicolumn{3}{l}{\small\textit{(continued)}}\\[4pt]
\toprule\textbf{Item}&\textbf{Status}&\textbf{Notes}\\\midrule
\endhead
\midrule\multicolumn{3}{r}{\small\textit{(continues)}}\\
\endfoot
\bottomrule
\endlastfoot
%%
THM-ZC84-01 (Complete Sector Synthesis)
  & Near-Formal
  & 8 clauses; weakest clause = Near-Formal; all from $\delta$ \\[4pt]
FIND-ZC84-01 (Full derivation chain)
  & Exact
  & $\delta \to$ all sectors; first unified display \\[4pt]
COR-ZC84-01 (Falsification table)
  & Near-Formal
  & EM+W ratio row = Nobel-table; assembled from ZC64/70/72/75/81/82 \\[4pt]
PM-ZC84-SCOPE-01 (scope guard)
  & Scar
  & scope creep; pre-registered; ZC84 = assembly only \\[4pt]
\midrule
Unified 4-sector + floor + cosmo display
  & Closed $\checkmark$
  & first time all arcs assembled in one paper \\
\bottomrule
\end{longtable}
}

%%=============================================================
\subsection{Canonical $\varepsilon$ Reference Table}
\label{sec:eps-table}
%%=============================================================

\begin{keybox}[title={Canonical $\varepsilon$ Ladder ---
  v17.14 (ZC84: no new values)}]
\sloppy
T-canonical: $n=8$, $\ds=3$, $\nm=5$,
$\Ks = 1-e^{-1/8} = 0.117503$.

{\small\renewcommand{\arraystretch}{1.35}
\begin{center}
\begin{tabular}{@{}p{3.2cm}p{1.5cm}p{1.5cm}p{5.5cm}@{}}
\toprule
\textbf{Symbol} & \textbf{Value} & \textbf{Layer} & \textbf{Source} \\
\midrule
$\epsEM = \epsS$ & $0.007494$ & bridge
  & ZC70/ZC81; unchanged \\[3pt]
$\epscan = \epsTLS$ & $0.007550$ & spectral
  & ZC79; unchanged \\[3pt]
$\epsuniv$ & $0.007599$ & bridge
  & ZC69; unchanged \\[3pt]
$\epsW$ & $0.007705$ & bridge
  & ZC70; unchanged \\[4pt]
\midrule
\multicolumn{4}{@{}p{12.5cm}}{\textit{Ordering (THM-ZC83-01):}
  $\epsEM = \epsS < \epscan < \epsuniv < \epsW$}\\[2pt]
\multicolumn{4}{@{}p{12.5cm}}{\textit{GM (FIND-ZC83-01):}
  $\sqrt{\epsEM\cdot\epsW} = \epsuniv$ (exact)}\\[2pt]
\multicolumn{4}{@{}p{12.5cm}}{\textit{Two-layer (ZC71):}
  $\epsuniv - \epscan = 0.66\%$; irreducible}\\
\bottomrule
\end{tabular}
\end{center}}
\end{keybox}

%%=============================================================
\subsection{R-Layer Research Note}
\label{sec:rlayer}
%%=============================================================

\begin{layerbox}[title={R-Layer Assignments for ZC84}]

All results in ZC84 are at $R_0$. No bridge crossing enters.

{\small\renewcommand{\arraystretch}{1.25}
\begin{center}
\begin{tabular}{@{}llll@{}}
\toprule
Atomic unit & R-layer & Cost $T$ & Source \\
\midrule
THM-ZC84-01 & $R_0$ & $0$ (synthesis) & all pillar papers \\
FIND-ZC84-01 & $R_0$ & $0$ & chain display \\
COR-ZC84-01 & $R_0$ & $0$ & assembled FS \\
\bottomrule
\end{tabular}
\end{center}}

\textbf{Two-layer floor (PWA-8):}
Spectral ($\epscan$) and bridge ($\epsuniv$) layers inherited from
ZC71/ZC79/ZC83. ZC84 makes no new layer claim; both layers appear
in the floor lattice display only as pillars.

Open gaps at $R_0{\to}R_{\max}$: GAP-ZC81-01 (Strong bridge),
GAP-ZC76-01 ($\tau_{\rm sweep}$) --- out of scope for ZC84.
\end{layerbox}

%%=============================================================
\subsection{Genealogy Map}
\label{sec:genealogy}
%%=============================================================

\begin{genealogybox}[title={How ZC84's Key Objects Trace Back}]

\textbf{Branch A --- Electroweak arc:}
\begin{itemize}[leftmargin=2em, noitemsep]
  \item \textbf{ZC20}: $\swsq$ exact.
  \item $\to$ \textbf{ZC28/ZC29}: $\GammaCRF$, $\alpha_i$ hierarchy.
  \item $\to$ \textbf{ZC57/ZC67/ZC70}: sector weights, floor formula.
  \item $\to$ \textbf{ZC79/ZC83}: $\epscan$ exact; ordering + GM.
  \item $\to$ \textbf{ZC80}: 7 EW observables at $R_0$.
  \item $\to$ \textbf{ZC84/THM-ZC84-01}: clauses (i)--(v).
\end{itemize}

\textbf{Branch B --- Strong + Ladder arc:}
\begin{itemize}[leftmargin=2em, noitemsep]
  \item \textbf{ZC30}: spatial-type class.
  \item $\to$ \textbf{ZC73}: $\Gamma$-ladder, $\DeltaGam = \log_{15}2$.
  \item $\to$ \textbf{ZC81}: $\epsS = \epsEM$; mirror symmetry.
  \item $\to$ \textbf{ZC84/THM-ZC84-01}: clauses (iv) and (vi).
\end{itemize}

\textbf{Branch C --- Cosmological arc:}
\begin{itemize}[leftmargin=2em, noitemsep]
  \item \textbf{ZC76}: $\wDE = \Gcoup/d_s - 1$.
  \item $\to$ \textbf{ZC82}: $\wDE + 3\DeltaGam = \GammaCRF$.
  \item $\to$ \textbf{ZC84/THM-ZC84-01}: clause (vii).
\end{itemize}

\textbf{Branch D --- Gravity arc:}
\begin{itemize}[leftmargin=2em, noitemsep]
  \item \textbf{ZC72}: vacuum sector; $\agrav = \GammaCRF/15$.
  \item $\to$ \textbf{ZC74}: $\mathbb{Z}_{15}$ uniqueness.
  \item $\to$ \textbf{ZC75}: $\agrav\cdot|\Ztf| = \GammaCRF$.
  \item $\to$ \textbf{ZC84/THM-ZC84-01}: clauses (iii) and (viii).
\end{itemize}

\textbf{Convergence at ZC84:}
All four branches converge in THM-ZC84-01. The convergence is not
algebraic coincidence: all branches share the root
$\delta \to d_s = 3 \to |\Ztf| = 15$. The eight simultaneous
results are forced by that single cascade.

ZC84's contribution: the first paper to display all four arcs in
one derivation chain, demonstrating that the zero-free-parameter
property holds collectively, not merely per-arc.
\end{genealogybox}

%%=============================================================
\subsection{Closing Sentence}
%%=============================================================
%% TONE C — AUSTERE

\bigskip
\begin{center}
\[
  \delta
  \;\xrightarrow{\;3d_s = 2^{d_s}+1\;}\;
  \mathbb{Z}_{15}
  \;\xrightarrow{\;\varphi(15)=8\;}\;
  \left\{\,
    \sin^2\!\theta_W,\;
    \alpha_{\rm EM},\;
    \alpha_{\rm grav},\;
    \varepsilon_{\rm EM} < \varepsilon_{\rm can}
      < \varepsilon_{\rm univ} < \varepsilon_W,\;
    w_{\rm DE}
  \,\right\}_{\!\delta}
\]
\end{center}

% Governance Note: This document follows Upgrade Protocol v1.2
% and CRF Style Guide v3.5.


\newpage
%%=============================================================
\part{ZC85--86 Cluster: The $\varepsilon_{\rm SC}$ Ontology,
  Opened and Closed (v17.15)}
\label{part:LIV}
%%=============================================================

\begin{volumebox}[title={Part LIV --- ZC85--86: From a Gated Gap to a
  Closed Bracket (v17.15)}]
\sloppy
\textbf{Part LIV} is the heart of the closure arc. ZC83 left GAP-ZC83-01
open: what \emph{is} the observed bracket
$\varepsilon_{\rm SC}=0.007605$, the $+0.08\%$ sitting just above
$\varepsilon_{\rm univ}$? These two papers answer it --- first
structurally, then in magnitude.

\textbf{ZC85} (The Observed-Bridge-Bracket Theorem) shows that
$\varepsilon_{\rm SC}$ is the geometric mean of the two per-observable
bridge floors of FIND-ZC65-02 (FIND-ZC85-01) --- the \emph{same}
$\sqrt{\varepsilon\varepsilon}$ structure as $\varepsilon_{\rm univ}$,
evaluated one $R$-layer up. The $+0.08\%$ was never a new constant
(FIND-ZC85-02). COR-ZC85-01 refines GAP-ZC83-01: the structural ontology
is \emph{closed}, but the magnitude --- which depends on the two running
shifts --- is \emph{gated} on GAP-ZC81-01 (the Strong/EW running bridge),
not an independent open problem. PM-ZC85-01 records honestly that
``$\varepsilon_{\rm SC}=$ GM of the bridge pair'' is a restatement of
existing structure (FIND-ZC65-02 + COR-ZC67-01), not a fresh derivation.

\textbf{ZC86} (The Strong Running Bridge Theorem) closes the gate.
THM-ZC86-01 derives the Strong-sector $R_0\!\to\!R_{\max}$ running bridge
and shows it is identical in form \emph{and} magnitude to the EM bridge,
because $\varepsilon_{\rm strong}=\varepsilon_{\rm EM}$ (THM-ZC81-01,
spatial-type) and the $R_1$ crossing cost is identical (LEM-ZC81-01).
LEM-ZC86-01 (confinement--coupling decoupling) removes the only worry:
the bridge perturbs a positive coupling, never the negative confinement
decay rate --- confinement lives in $\Gamma$, not in $\alpha$. With
GAP-ZC81-01 closed, the gate of COR-ZC85-01 lifts: FIND-ZC86-01 (gift)
identifies the running-floor layer as having exactly two floors, whose
geometric mean is $\varepsilon_{\rm SC}$ --- the $+0.08\%$ magnitude is
now closed. PM-ZC86-01 preserves a failed universal-coupling propagation
path that lit the way.

\textbf{Key closures:}
GAP-ZC81-01 (Strong running bridge) --- CLOSED (ZC86/THM-ZC86-01,
near-formal). GAP-ZC83-01 ($\varepsilon_{\rm SC}$ ontology and closed
form) --- CLOSED (structural via ZC85/FIND-ZC85-01; magnitude via
ZC86/FIND-ZC86-01).

\textbf{Phase Misalignments preserved:}
PM-ZC85-01 (restatement, not derivation; pattern: tautology test);
PM-ZC86-01 (universal-coupling propagation; failed probe, magnitudes
$0.15\%/0.41\%$, preserved as scar).
\end{volumebox}

\newpage
%%─── VERBATIM EMBED: ZC85 ───────────────────────────────────────────────
\addcontentsline{toc}{section}{ZC85: The Observed-Bridge-Bracket Theorem}

% [BRIDGE-PROSE: integration-added, Extension Freedom narrative, v3.6 §31.3]
The previous cluster ended with a loose thread it could not tie. Among
the floors there is one value that refuses to fall into the clean pattern
of the others: an observed bracket sitting about eight hundredths of a
percent above the universal floor, close enough to look like a rounding
artefact, stubborn enough that it is not. The earlier work had already
ruled out the easy explanations, which left an uncomfortable question
hanging --- is this a genuine feature of the framework, or a number CRF
simply does not account for?

This paper answers it by reversing the usual direction of work. Instead
of starting from the axioms and deriving forward, it starts from the
measured quantity and reasons backward to ask what structure could
produce exactly it. The answer, when it comes, is almost deflating in its
simplicity: the bracket is the same geometric-mean relation already seen
among the floors, only evaluated one layer further out. Nothing new had
to be invented; the structure was already present, one step away.

That tidiness is also the paper's caution, and the sections below are
careful about it: showing that a number equals a mean of two others you
already had is a restatement, not a fresh derivation. What ZC85 genuinely
settles is \emph{what kind of thing} the bracket is. Its exact size still
hangs on a separate question about how the Strong sector runs --- the
thread the next paper picks up.



%%=============================================================
\begin{abstract}
\sloppy
ZC83 closed the floor-lattice ordering but left GAP-ZC83-01 open: the
ontology and closed form of the observed bridge bracket
$\epsSC = 0.007605$, which sits a structural $+0.08\%$ above
$\epsuniv = 0.007599$. This paper resolves the ontology and reframes
the remaining open question.

\textbf{FIND-ZC85-01} (Observed-Bridge-Bracket Identity):
\[
  \epsSC \;=\; \GMfloor\!\left(\epszA,\ \epszS\right)
  \;=\; \sqrt{\,0.007501 \times 0.007710\,}
  \;=\; 0.0076048,
\]
where $\epszA$, $\epszS$ are the two per-observable bridge floors of
FIND-ZC65-02 (one solved from $\delta\alpha/\alpha$, one from
$\delta(\sin^2\theta_W)$). This is the \emph{same} COR-ZC67-01
geometric-mean structure as $\epsuniv = \GMfloor(\epsEM,\epsW)$
(FIND-ZC83-01), evaluated one $R$-layer up: $\epsuniv$ at $R_0$,
$\epsSC$ at $R_0{\to}R_{\max}$.

\textbf{FIND-ZC85-02}: the $+0.08\%$ factors into two independent
running channels --- EM ($+0.09\%$) and Weak ($+0.07\%$) --- and
$\epsSC/\epsuniv$ is the geometric mean of the two channel ratios.

\textbf{COR-ZC85-01}: GAP-ZC83-01 is thereby refined. Its structural
ontology is \emph{closed}; the remaining question --- a closed form
for the \emph{magnitude} of each running shift --- is \emph{gated} on
GAP-ZC81-01 (the Strong/EW running bridge), not an independent gap.
$\epsSC$ is not a new cascade constant.
\end{abstract}


%%=============================================================
\subsection{Background and Scope}
\label{sec:background}
%%=============================================================

\begin{keybox}[title={Numerical Summary --- ZC85 (T-canonical)}]
{\small\renewcommand{\arraystretch}{1.25}
\begin{center}
\begin{tabular}{@{}p{5.6cm}p{2.4cm}p{1.5cm}p{2.6cm}@{}}
\toprule
\textbf{Quantity} & \textbf{Value} & \textbf{Gap} & \textbf{Source} \\
\midrule
$\epszA$ ($\alpha$ bridge floor)   & $0.007501$ & --- & FIND-ZC65-02 \\
$\epszS$ ($\sin^2\theta_W$ bridge) & $0.007710$ & --- & FIND-ZC65-02 \\
$\GMfloor(\epszA,\epszS)$ (NEW)    & $0.0076048$ & --- & \textbf{FIND-ZC85-01} \\
$\epsSC$ (observed bracket)        & $0.007605$ & $8{\times}10^{-8}$ & ZC68/THM-ZC68-01b \\
\midrule
$\epsuniv = \GMfloor(\epsEM,\epsW)$ & $0.007599$ & exact & ZC83/FIND-ZC83-01 \\
$\epsSC/\epsuniv - 1$               & $+0.080\%$ & --- & two-layer offset \\
\midrule
EM channel $\epszA/\epsEM$          & $+0.090\%$ & --- & FIND-ZC85-02 \\
Weak channel $\epszS/\epsW$         & $+0.069\%$ & --- & FIND-ZC85-02 \\
\bottomrule
\end{tabular}
\end{center}}
\end{keybox}

\subsubsection{The Open Item Inherited from ZC83}

ZC83/GAP-ZC83-01 stated the open problem precisely: find a closed form
for $\epsSC = 0.007605$, or prove it transcendental over the known CRF
floor extensions. PM-ZC83-01 had already refuted the natural probe
``$\epsSC = \AMfloor(\epsEM,\epsW)$'' (the arithmetic mean falls
$0.073\%$ short). The question left standing was: \emph{what invariant
of the $\mathbb{Z}_{15}$ structure does $\epsSC$ probe?}

This paper answers it by working backward --- from the observed
quantity to its origin --- rather than forward from axioms.

\subsubsection{What Was Already Known}

\begin{keybox}[title={Pre-ZC85 Status}]
\textbf{Known (exact):}
\begin{itemize}[leftmargin=2em, noitemsep]
  \item $\epsuniv = \GMfloor(\epsEM,\epsW)$ exactly --- the bridge
    floor is the geometric mean of the \emph{defined} sector pair
    at $R_0$ (ZC83/FIND-ZC83-01).
  \item $\epsSC = \sqrt{\epsEM^{\rm obs}\,\epsW^{\rm obs}}$ --- the
    \emph{observed} bracket (ZC68/THM-ZC68-01b), $+0.08\%$ above
    $\epsuniv$ (ZC71 two-layer separation).
  \item COR-ZC67-01 is \emph{one} identity
    $\sqrt{\varepsilon_{\rm EM}\varepsilon_W}=\varepsilon_{\rm base}$
    on three input pairs (Part LII verified $\varepsilon$ notice).
\end{itemize}

\textbf{Known but not connected:}
\begin{itemize}[leftmargin=2em, noitemsep]
  \item FIND-ZC65-02 had already solved the two bridge equations
    \emph{separately} for the bridge floor implied by each observable:
    $\epszA = 0.007501$ (from $\delta\alpha/\alpha$) and
    $\epszS = 0.007710$ (from $\delta(\sin^2\theta_W)$), noting both
    point to $\varepsilon_0^{\rm exact}\approx 0.00761$.
  \item No paper had taken the geometric mean of these two values.
\end{itemize}

\textbf{Open before ZC85:}
\begin{itemize}[leftmargin=2em, noitemsep]
  \item The ontology of $\epsSC$ (GAP-ZC83-01).
  \item Whether the $+0.08\%$ has a closed form.
\end{itemize}
\end{keybox}

\begin{keybox}[title={What ZC85 Does NOT Do}]
\begin{itemize}[leftmargin=2em, noitemsep]
  \item Does not re-derive $\epsuniv$, $\epsEM$, $\epsW$, or the
    FIND-ZC65-02 bridge floors (used as exact pillars).
  \item Does not derive the \emph{magnitude} of the running shifts
    from $\delta$ alone --- this is declared gated on GAP-ZC81-01,
    not attempted here.
  \item Does not modify COR-ZC67-01, FIND-ZC65-02, or the ZC81/82
    retraction scars (read-only).
  \item Does not re-open the verified $\varepsilon$-floor notice
    (Part LII); this paper is consistent with it.
\end{itemize}
\end{keybox}

\begin{derivedbox}[title={Discovery Note}]
The Council was tasked to attack GAP-ZC83-01 ``to settle all doubt.''
Following the bidirectional method (observed $\to$ origin), a probe
decomposed the $+0.08\%$ offset into two running channels, then asked
the reverse question: \emph{which existing CRF objects, taken as a
geometric mean, reproduce $\epsSC$?} The answer was immediate ---
$\sqrt{0.007501\times 0.007710} = 0.0076048$, the two FIND-ZC65-02
bridge floors, matching $\epsSC$ to $8\times10^{-8}$. The gap had been
hiding a restatement: $\epsSC$ is the COR-ZC67-01 geometric mean
evaluated one $R$-layer above $\epsuniv$.
\end{derivedbox}

\subsubsection{Pre-Work Audit (PWA-1 through PWA-10)}

\begin{formalbox}[title={Pre-Work Audit Record --- ZC85}]
Scanned before writing (Upgrade Protocol v1.2):
\begin{itemize}[leftmargin=2em, noitemsep]
  \item ZC65/FIND-ZC65-02 --- \textbf{KEY PILLAR} (the two
    per-observable bridge floors $\epszA = 0.007501$,
    $\epszS = 0.007710$; Part~I \S\,for FIND-ZC65-02).
  \item ZC67/COR-ZC67-01 --- \textbf{KEY PILLAR} (GM structure;
    one identity, three input pairs; Part LII notice).
  \item ZC68/THM-ZC68-01b --- pillar ($\epsSC$ observed bracket).
  \item ZC69/COR-ZC69-02, ZC70/THM-ZC70-01 --- pillars
    ($\epsuniv$, $\epsEM$, $\epsW$ exact).
  \item ZC83/FIND-ZC83-01 --- pillar ($\epsuniv = \GMfloor$ at $R_0$;
    the $R_0$-layer twin of FIND-ZC85-01).
  \item ZC81/GAP-ZC81-01 --- the gate for the residual open part.
  \item Part LII verified $\varepsilon$-notice --- read; ZC85
    is consistent with it (no reconciliation attempted).
\end{itemize}

\textbf{PWA-5} (tautology --- \emph{critical for this paper}):
FIND-ZC85-01 is an \emph{Exact restatement}, not a new derivation.
The bridge floors $\epszA$, $\epszS$ are independently fitted to
\emph{observed} gaps (not constructed by multiplying a chosen
$\varepsilon_0$ back through $(\ds/2)\Ks$ and $\Ks$); the observed
values are genuinely above the defined ones (verified, not rounding).
Hence their GM equalling $\epsSC$ is a real cross-layer identity, but
its only novelty is the cross-layer \emph{reading}. Registered as
Exact-restatement (Gift), \emph{not} promoted to Theorem, to avoid the
substitution-tautology pattern retracted in COR-ZC81-02/COR-ZC82-02
(PM-ZC81-RETRACT-01, PM-ZC82-RETRACT-02).\\
\textbf{PWA-6} (macro): $\epszA$, $\epszS$, $\epsSC$ added via
\verb|\providecommand|; no conflict.\\
\textbf{PWA-7} (domain): EM/Weak channel split stays within the
electroweak running domain; no Spontaneous/Agentive conflation.\\
\textbf{PWA-8} (two-layer): $R_0$ vs $R_0{\to}R_{\max}$ distinction
maintained throughout; this is the load-bearing structure.\\
\textbf{No duplicate atomic units found.}
\end{formalbox}

\newpage

%%=============================================================
\subsection{FIND-ZC85-01\quad The Observed-Bridge-Bracket Identity}
\label{sec:find01}
%%=============================================================

\needspace{12\baselineskip}

\begin{giftbox}[title={FIND-ZC85-01: $\epsSC = \GMfloor(\epszA,\epszS)$
  ($R_0{\to}R_{\max}$, Exact restatement --- Gift)}]

\textbf{Statement.}
\begin{equation}
  \boxed{
    \epsSC \;=\; \GMfloor\!\left(\epszA,\ \epszS\right)
    \;=\; \sqrt{\,\epszA \cdot \epszS\,},
  }
  \label{eq:find85-01}
\end{equation}
where $\epszA$ and $\epszS$ are the two per-observable bridge floors
of FIND-ZC65-02:
\begin{align}
  \epszA &= \frac{\delta\alpha/\alpha}{(\ds/2)\,\Ks} = 0.007501
    && \text{(EM bridge, ZC65/FIND-ZC65-01 inverted),}
    \label{eq:epszA} \\
  \epszS &= \frac{\delta(\sin^2\theta_W)}{\Ks} = 0.007710
    && \text{(Weak bridge, CONJ-ZC66-01 inverted).}
    \label{eq:epszS}
\end{align}

\textbf{Proof.}
THM-ZC68-01b defines the observed bracket as
$\epsSC = \sqrt{\epsEM^{\rm obs}\cdot\epsW^{\rm obs}}$. By the bridge
equations of FIND-ZC65-02, the observed EM-sector floor solved from
$\delta\alpha/\alpha$ is exactly $\epszA$, and the observed
Weak-sector floor solved from $\delta(\sin^2\theta_W)$ is exactly
$\epszS$. Substituting,
\begin{equation}
  \epsSC
  = \sqrt{\epsEM^{\rm obs}\cdot\epsW^{\rm obs}}
  = \sqrt{\epszA\cdot\epszS}
  = \sqrt{0.007501 \times 0.007710}
  = 0.0076048.
  \label{eq:find85-proof}
\end{equation}
This matches $\epsSC = 0.007605$ (ZC68, four significant figures) to
$8\times10^{-8}$. \hfill$\square$

\textbf{The cross-layer reading.}
ZC83/FIND-ZC83-01 proved $\epsuniv = \GMfloor(\epsEM,\epsW)$ --- the
geometric mean of the \emph{defined} sector floors at $R_0$.
FIND-ZC85-01 shows $\epsSC = \GMfloor(\epszA,\epszS)$ --- the
geometric mean of the \emph{observed} bridge floors fitted across
$R_0{\to}R_{\max}$. Both are the COR-ZC67-01 geometric-mean structure
$\sqrt{\varepsilon\varepsilon}=\varepsilon_{\rm base}$; they differ
only in which $R$-layer the input pair lives on:

\begin{center}
{\small\renewcommand{\arraystretch}{1.35}
\begin{tabular}{@{}p{3.2cm}p{4.0cm}p{2.0cm}p{2.6cm}@{}}
\toprule
\textbf{Base} & \textbf{Input pair} & \textbf{Layer} & \textbf{Value} \\
\midrule
$\epsuniv$ & $(\epsEM,\epsW)$ defined & $R_0$ & $0.007599$ \\
$\epsSC$   & $(\epszA,\epszS)$ observed & $R_0{\to}R_{\max}$ & $0.007605$ \\
\bottomrule
\end{tabular}}
\end{center}

\textbf{Why this dissolves the gap.}
GAP-ZC83-01 asked ``what invariant does $\epsSC$ probe?'' The answer:
none new. $\epsSC$ probes the \emph{same} geometric-mean invariant as
$\epsuniv$; the $+0.08\%$ is entirely the shift in the input pair from
the defined ($R_0$) layer to the observed (running) layer. No new
cascade constant exists or is needed.

\textbf{Classification:} Exact restatement (Gift). Registered as
Finding, \emph{not} Theorem (PWA-5; see PM-ZC85-01).
\end{giftbox}

%%=============================================================
\subsection{FIND-ZC85-02\quad Two-Channel Decomposition of the $+0.08\%$}
\label{sec:find02}
%%=============================================================

\needspace{10\baselineskip}

\begin{derivedbox}[title={FIND-ZC85-02: The Offset Factors into
  Two Running Channels ($R_0{\to}R_{\max}$, Exact)}]

\textbf{Statement.}
The two-layer offset is the geometric mean of two independent
per-observable running-shift ratios:
\begin{equation}
  \frac{\epsSC}{\epsuniv}
  \;=\;
  \sqrt{\frac{\epszA}{\epsEM}\cdot\frac{\epszS}{\epsW}},
  \label{eq:find85-02}
\end{equation}
with
\begin{align}
  \frac{\epszA}{\epsEM} &= \frac{0.007501}{0.007494} = 1.00090
    && (+0.090\%,\ \text{EM channel}), \\
  \frac{\epszS}{\epsW}  &= \frac{0.007710}{0.007705} = 1.00069
    && (+0.069\%,\ \text{Weak channel}).
\end{align}

\textbf{Proof.}
Divide \eqref{eq:find85-01} by $\epsuniv = \sqrt{\epsEM\epsW}$
(FIND-ZC83-01):
\[
  \frac{\epsSC}{\epsuniv}
  = \frac{\sqrt{\epszA\,\epszS}}{\sqrt{\epsEM\,\epsW}}
  = \sqrt{\frac{\epszA}{\epsEM}\cdot\frac{\epszS}{\epsW}}
  = \sqrt{1.00090 \times 1.00069}
  = 1.00080.
\]
\hfill$\square$

\textbf{Reading.}
The $+0.08\%$ is not a single number to be derived. It is the
geometric mean of two independent running shifts --- how far the EM
and Weak observables run from $R_0$ to $R_{\max}$. Each shift is a
per-sector bridge effect, not a property of the floor lattice itself.

\textbf{Classification:} Exact Finding.
\end{derivedbox}

%%=============================================================
\subsection{COR-ZC85-01\quad GAP-ZC83-01 Refined to Gated}
\label{sec:cor01}
%%=============================================================

\needspace{10\baselineskip}

\begin{formalbox}[title={COR-ZC85-01: The Residual of GAP-ZC83-01 Is
  Gated on GAP-ZC81-01}]

\textbf{Statement.}
GAP-ZC83-01 splits cleanly into two parts:
\begin{enumerate}[leftmargin=2.2em, label=(\arabic*)]
  \item \textbf{Structural ontology --- CLOSED.} By FIND-ZC85-01,
    $\epsSC$ is the COR-ZC67-01 geometric mean of the observed bridge
    pair; it is not a new constant. The question ``what is $\epsSC$''
    is answered.
  \item \textbf{Magnitude closed form --- GATED.} A closed form for
    the size of each running shift ($\epszA/\epsEM$, $\epszS/\epsW$)
    from $\delta$ alone requires the running bridge to be derived for
    each sector. The EM running bridge is ZC65/THM-ZC65-01 (near-formal,
    observation-bounded); the analogous Weak/Strong bridge is the
    content of GAP-ZC81-01.
\end{enumerate}
Therefore the residual open part of GAP-ZC83-01 is \emph{gated} on
GAP-ZC81-01: it is not an independent problem, and closing
GAP-ZC81-01 closes it.

\textbf{Consequence.}
\begin{itemize}[leftmargin=2em, noitemsep]
  \item GAP-ZC83-01 is downgraded from ``open, Priority 1'' to
    ``structurally closed; magnitude gated on GAP-ZC81-01.''
  \item No new gap is opened.
  \item The next natural target is GAP-ZC81-01 (Strong/EW running
    bridge), which now subsumes the $\epsSC$ magnitude question.
\end{itemize}

\textbf{Classification:} Corollary (gap refinement, Near-Formal).
\end{formalbox}

%%=============================================================
\subsection{Phase Misalignment}
\label{sec:pm}
%%=============================================================

\begin{openqbox}[title={PM-ZC85-01: FIND-ZC85-01 Is a Restatement,
  Not a New Derivation (pattern: tautology test)}]

\textbf{Probe.}
The natural temptation, on finding
$\sqrt{0.007501\times 0.007710} = \epsSC$, is to promote it to a
Theorem ``$\epsSC$ derived.''

\textbf{Guard.}
This would repeat the failure pattern of COR-ZC81-02 and COR-ZC82-02
(retracted as substitution tautologies; PM-ZC81-RETRACT-01,
PM-ZC82-RETRACT-02). FIND-ZC85-01 reuses two values
($\epszA$, $\epszS$) that were themselves fitted to the observed
gaps. Their geometric mean reproducing $\epsSC$ is a genuine
cross-layer identity --- the inputs live on the running layer, not
constructed circularly --- but it \emph{derives} nothing new about the
magnitude of those gaps.

\textbf{Resolution.}
Register FIND-ZC85-01 as an \emph{Exact restatement} (Gift): it
closes the \emph{ontology} of $\epsSC$ (a real result) while making no
claim to derive its \emph{magnitude} from axioms (that is gated,
COR-ZC85-01). The distinction is the whole point of the scar.

\textbf{Status:} Scar (pre-registered; pattern: tautology test).
No retraction --- the restatement is correctly scoped from the start.
\end{openqbox}

%%=============================================================
\subsection{Complete Floor Lattice: $\epsSC$ Located}
\label{sec:lattice}
%%=============================================================

\begin{keybox}[title={Floor Lattice with $\epsSC$ Ontology Resolved
  (T-canonical)}]

\begin{center}
{\small\renewcommand{\arraystretch}{1.35}
\begin{tabular}{@{}p{2.6cm}p{1.7cm}p{2.6cm}p{4.6cm}@{}}
\toprule
\textbf{Floor} & \textbf{Value} & \textbf{Layer} & \textbf{Structure} \\
\midrule
$\epsEM = \epsS$ & $0.007494$ & $R_0$ bridge
  & $\epsuniv\cdot r^{-1/4}$ (ZC70/81) \\
$\epscan = \epsTLS$ & $0.007550$ & $R_0$ spectral
  & quadratic root (ZC79) \\
$\epsuniv$ & $0.007599$ & $R_0$ bridge
  & $\GMfloor(\epsEM,\epsW)$ (ZC83) \\
$\epsW$ & $0.007705$ & $R_0$ bridge
  & $\epsuniv\cdot r^{+1/4}$ (ZC70) \\
\midrule
$\epsSC$ & $0.007605$ & $R_0{\to}R_{\max}$
  & $\GMfloor(\epszA,\epszS)$ \textbf{(ZC85)} \\
\bottomrule
\end{tabular}}
\end{center}

\medskip
\textbf{The unifying statement (ZC85):}
\begin{align}
  \epsuniv &= \GMfloor(\epsEM,\epsW)
    && \text{defined pair, $R_0$ \quad (FIND-ZC83-01)} \\
  \epsSC   &= \GMfloor(\epszA,\epszS)
    && \text{observed pair, $R_0{\to}R_{\max}$ \quad (FIND-ZC85-01)}
\end{align}
Both are the single COR-ZC67-01 geometric-mean identity. The $+0.08\%$
between them is the geometric mean of two per-observable running shifts
(FIND-ZC85-02), whose magnitude is gated on GAP-ZC81-01 (COR-ZC85-01).

\medskip
\textbf{Consistency with Part LII notice.}
This is exactly the ``one identity, three input pairs'' reading of the
verified $\varepsilon$-floor notice. ZC85 adds the explicit input pair
for the observed layer ($\epszA,\epszS$) and the cross-layer reading;
it reconciles nothing and changes no floor value.
\end{keybox}

%%=============================================================
\subsection{Atomic Registry}
\label{sec:registry}
%%=============================================================

\begin{formalbox}[title={Atomic Registry --- ZC85}]
\begin{center}
\renewcommand{\arraystretch}{1.3}
\small
\begin{tabular}{lp{7.0cm}l}
\toprule
\textbf{ID} & \textbf{Description} & \textbf{Status} \\
\midrule
FIND-ZC85-01
  & $\epsSC = \GMfloor(\epszA,\epszS)$
    ${} = \sqrt{0.007501\cdot0.007710}$
    ${} = 0.0076048$;
    COR-ZC67-01 GM at the $R_{\max}$ layer
  & Exact (restatement, Gift) \\[3pt]
FIND-ZC85-02
  & $\epsSC/\epsuniv = \sqrt{(\epszA/\epsEM)(\epszS/\epsW)}$;
    $+0.08\%$ = GM of EM ($+0.09\%$) and Weak ($+0.07\%$) channels
  & Exact \\[3pt]
COR-ZC85-01
  & GAP-ZC83-01 refined: ontology closed (FIND-ZC85-01);
    magnitude gated on GAP-ZC81-01
  & Near-Formal \\[3pt]
PM-ZC85-01
  & FIND-ZC85-01 is restatement not derivation; guarded against
    the COR-ZC81-02/82-02 tautology pattern
  & Scar \\[3pt]
\bottomrule
\end{tabular}
\end{center}
\end{formalbox}

%%=============================================================
\subsection{R-Layer Research Note}
\label{sec:rlayer}
%%=============================================================

\begin{layerbox}[title={R-Layer Assignments for ZC85}]

{\small\renewcommand{\arraystretch}{1.25}
\begin{center}
\begin{tabular}{@{}llll@{}}
\toprule
Atomic unit & R-layer & Cost $T$ & Source \\
\midrule
FIND-ZC85-01 & $R_0{\to}R_{\max}$ & algebraic & ZC65 + ZC67 + ZC68 \\
FIND-ZC85-02 & $R_0{\to}R_{\max}$ & algebraic & ZC83 + ZC85/FIND-01 \\
COR-ZC85-01  & --- & --- & gap refinement; gates on ZC81 \\
\bottomrule
\end{tabular}
\end{center}}

\textbf{Two-layer discipline (PWA-8):} This paper is fundamentally a
\emph{cross-layer} statement. $\epsuniv$ lives at $R_0$ (defined
sector floors); $\epsSC$ lives at $R_0{\to}R_{\max}$ (observed bridge
floors). The $+0.08\%$ is precisely the layer shift, not a property of
either layer alone. No bridge \emph{magnitude} is derived here; that
is gated on GAP-ZC81-01.
\end{layerbox}

%%=============================================================
\subsection{Genealogy Map}
\label{sec:genealogy}
%%=============================================================

\begin{genealogybox}[title={How ZC85's Key Objects Trace Back}]

\textbf{Branch A --- Geometric-mean structure:}
\begin{itemize}[leftmargin=2em, noitemsep]
  \item \textbf{ZC67/COR-ZC67-01}: one identity
    $\sqrt{\varepsilon_{\rm EM}\varepsilon_W}=\varepsilon_{\rm base}$.
  \item $\to$ \textbf{ZC83/FIND-ZC83-01}: $\epsuniv = \GMfloor$ of
    defined pair at $R_0$.
  \item $\to$ \textbf{ZC85/FIND-ZC85-01}: $\epsSC = \GMfloor$ of
    observed pair at $R_0{\to}R_{\max}$ (same identity, next layer).
\end{itemize}

\textbf{Branch B --- Bridge-floor chain:}
\begin{itemize}[leftmargin=2em, noitemsep]
  \item \textbf{ZC65/FIND-ZC65-01}: $\delta\alpha/\alpha
    = (\ds/2)\varepsilon_0\Ks$ (the EM bridge equation).
  \item \textbf{ZC65/FIND-ZC65-02}: solving each observable
    separately gives $\epszA = 0.007501$, $\epszS = 0.007710$.
  \item $\to$ \textbf{ZC85/FIND-ZC85-01}: their GM is $\epsSC$.
  \item $\to$ \textbf{ZC85/FIND-ZC85-02}: the two ratios are the
    two running channels.
\end{itemize}

\textbf{Branch C --- The gate:}
\begin{itemize}[leftmargin=2em, noitemsep]
  \item \textbf{ZC68/THM-ZC68-01b}: $\epsSC$ observed bracket defined.
  \item \textbf{ZC83/GAP-ZC83-01}: $\epsSC$ ontology opened.
  \item $\to$ \textbf{ZC85/COR-ZC85-01}: ontology closed; magnitude
    gated on \textbf{ZC81/GAP-ZC81-01} (Strong/EW running bridge).
\end{itemize}

\textbf{Convergence at ZC85:}
Branch~A supplies the geometric-mean identity, Branch~B supplies the
observed input pair that FIND-ZC65-02 had already computed but never
combined, and Branch~C identifies the gate. ZC85's contribution: the
gap GAP-ZC83-01 was never an independent problem --- $\epsSC$ is the
COR-ZC67-01 mean one $R$-layer up, and its magnitude belongs to the
running-bridge question of GAP-ZC81-01.
\end{genealogybox}

%%=============================================================
\subsection{Closing Sentence}
%%=============================================================
%% TONE B -- REFLECTIVE

\bigskip
\begin{center}
\textit{%
  The bracket that sat $0.08\%$ above the bridge floor\\
  was never a stranger to the lattice.\\[0.3em]
  It was the same geometric mean,\\
  taken one layer higher --- where the constants run.\\[0.5em]
  $\displaystyle
    \varepsilon_{\rm SC}
    \;=\;
    \sqrt{\varepsilon_0^{\alpha}\cdot\varepsilon_0^{\rm sw}}
    \qquad\text{(the observed pair).}$\\[0.3em]
  What remains open is not its name, but its size ---\\
  and that waits at GAP-ZC81-01.
}
\end{center}

% Governance Note: This document follows Upgrade Protocol v1.2
% and CRF Style Guide v3.5.


\newpage
%%─── VERBATIM EMBED: ZC86 ───────────────────────────────────────────────
\addcontentsline{toc}{section}{ZC86: The Strong Running Bridge Theorem}

% [BRIDGE-PROSE: integration-added, Extension Freedom narrative, v3.6 §31.3]
There was a real reason to worry. The technique that lets CRF carry a
coupling from one scale to another had only ever been proved on the
friendliest possible case: electromagnetism, whose strength grows as you
zoom in. The Strong force is the opposite kind of beast --- it confines,
its characteristic rung is negative, it behaves at short distances in a
way that has historically broken naive arguments. Borrowing a result
proved on the easy sector and assuming it survives a sign flip is exactly
the kind of move that quietly wrecks a framework. The honest question was
whether the bridge would hold, or whether the floor lattice was secretly
an electromagnetism-only result wearing a universal disguise.

ZC86 shows it holds, and the reason is cleaner than the worry. The thing
that makes the Strong sector dangerous --- confinement --- lives in a
different quantity from the one the bridge actually moves. The bridge acts
on the coupling; confinement sits in the decay rate. They do not touch.
Once that separation is made explicit, the Strong bridge turns out to be
identical to the electromagnetic one, in form and in size, and the worry
dissolves.

With that settled, the loose thread from the previous paper is finally
pulled tight: the observed bracket's exact magnitude was waiting on this
result, and the sections below close it. What looked like two separate
open problems were, all along, one question asked twice.



%%=============================================================
\begin{abstract}
\sloppy
GAP-ZC81-01 asked for the Strong sector $R_0{\to}R_{\max}$ running
bridge and flagged a worry: the Strong $\Gamma$-ladder rung is negative
($\GammaS = -0.498 < 0$, confined), so it was unclear whether the
EM bridge of ZC65 applies. This paper closes the gap.

\textbf{LEM-ZC86-01} (Confinement--Coupling Decoupling): the bridge
perturbs the \emph{coupling} $\aS = (2/15)\Gamma_{\rm CRF} > 0$;
confinement is a property of the \emph{decay rate} $\GammaS < 0$.
These are distinct objects, and the bridge never references $\GammaS$.
The confinement worry dissolves.

\textbf{THM-ZC86-01} (Strong Running Bridge):
\[
  \aS(R_{\max}) \;=\; \aS(R_0)\left[1 + \tfrac{d_s}{2}\,\epsS\,K^*\right],
\]
\emph{identical in form and magnitude} to the EM bridge, because
$\epsS = \epsEM$ (THM-ZC81-01, spatial-type) and the $R_1$ crossing
cost is identical (LEM-ZC81-01). ZC65's near-formal analysis transfers
verbatim. \textbf{GAP-ZC81-01 closed} (Near-Formal).

\textbf{FIND-ZC86-01} (Gift): the running floor layer has only
\emph{two} independent values --- $\epszA = 0.007501$ (EM\,=\,Strong,
spatial) and $\epszS = 0.007710$ (Weak, matter). Strong shares EM's
floor. This closes the $\epsSC$ magnitude question left gated by
COR-ZC85-01: the $+0.08\%$ is the geometric mean of exactly these two.

\textbf{PM-ZC86-01} preserves a failed path: propagating the bridge
through the shared universal coupling predicted the wrong channel
shifts ($0.15\%/0.41\%$ vs the observed $0.09\%/0.07\%$). The failure
is what revealed the per-sector mechanism.
\end{abstract}


%%=============================================================
\subsection{Background and Scope}
\label{sec:background}
%%=============================================================

\begin{keybox}[title={Numerical Summary --- ZC86 (T-canonical)}]
{\small\renewcommand{\arraystretch}{1.25}
\begin{center}
\begin{tabular}{@{}p{5.6cm}p{2.4cm}p{1.6cm}p{2.6cm}@{}}
\toprule
\textbf{Quantity} & \textbf{Value} & \textbf{Gap} & \textbf{Source} \\
\midrule
$\delta\aEMc/\aEMc = (d_s/2)\epsEM K^*$ & $0.001321$ & --- & ZC65/THM-ZC65-01 \\
$\delta\aS/\aS = (d_s/2)\epsS K^*$ (NEW) & $0.001321$ & $0\%$ vs EM & \textbf{THM-ZC86-01} \\
$\delta\alpha_W/\alpha_W = (d_s/2)\epsW K^*$ & $0.001358$ & --- & ZC65 (Weak) \\
\midrule
$\aS(R_0) = (2/15)\Gamma_{\rm CRF}^{\rm EM}$ & $0.001822$ & --- & ZC29/THM-ZC29-01 \\
$\aS(R_{\max})$                              & $0.001824$ & --- & THM-ZC86-01 \\
$\GammaS$ (decay rate, $k=1$)                & $-0.498252$ & --- & ZC73/THM-ZC73-01 \\
\midrule
$\epszA$ (EM\,=\,Strong floor)               & $0.007501$ & --- & FIND-ZC65-02 \\
$\epszS$ (Weak floor)                        & $0.007710$ & --- & FIND-ZC65-02 \\
\midrule
\textit{Failed probe} ($G^{\rm eff}$ path)   & $0.15\%/0.41\%$ & --- & \textbf{PM-ZC86-01} \\
\bottomrule
\end{tabular}
\end{center}}
\end{keybox}

\subsubsection{The Open Item}

ZC81 derived the Strong floor $\epsS = \epsEM$ (THM-ZC81-01) and the
ladder mirror $\GammaEM + \GammaS = 2\Gamma_W$ (FIND-ZC81-01), but
explicitly deferred the Strong running bridge:

\begin{quote}
\itshape
``The bridge formula of ZC65 is exact for the EM sector; whether the
same formula applies to a confined sector ($\GammaS < 0$) requires
additional analysis of the crossing integral in the confinement
regime.'' (GAP-ZC81-01)
\end{quote}

The worry was concrete: ZC65's bridge was derived for EM, whose ladder
rung $\GammaEM = +0.0137 > 0$ is the only positive (deconfined) rung.
The Strong rung is negative. Does the bridge survive?

\subsubsection{What Was Already Known}

\begin{keybox}[title={Pre-ZC86 Status}]
\textbf{Known (exact):}
\begin{itemize}[leftmargin=2em, noitemsep]
  \item $\epsS = \epsEM = 0.007494$ (THM-ZC81-01, spatial-type).
  \item $\TRone_{\rm strong} = \TRone_{\rm EM} = (d_s/n_m)\epsuniv^2$
    (LEM-ZC81-01; both spatial-type, FIND-ZC30-02).
  \item $\aS(R_0) = (2/15)\,\Gamma_{\rm CRF}$ (ZC29/THM-ZC29-01,
    $\varphi(3)=2$); positive.
  \item $\GammaS = G/d_s - 1 + 1\cdot\Delta_\Gamma = -0.498$
    (ZC73/THM-ZC73-01, rung $k=1$); negative (confined).
  \item EM bridge $\delta\aEMc/\aEMc = (d_s/2)\epsEM K^*$
    (ZC65/FIND-ZC65-01), near-formal.
\end{itemize}

\textbf{Open before ZC86:}
\begin{itemize}[leftmargin=2em, noitemsep]
  \item Whether the EM bridge applies to the confined Strong sector.
  \item The Strong $R_0{\to}R_{\max}$ running form and magnitude.
\end{itemize}
\end{keybox}

\begin{keybox}[title={What ZC86 Does NOT Do}]
\begin{itemize}[leftmargin=2em, noitemsep]
  \item Does not re-derive $\epsS$, $\epsEM$, $\TRone$ (exact pillars).
  \item Does not claim Exact: it inherits ZC65's near-formal status
    (the bridge carries a near-formal $\varepsilon_0$).
  \item Does not model confinement \emph{dynamics} (how $\GammaS < 0$
    evolves); it only shows confinement does not obstruct the bridge.
  \item Does not modify the $\alpha$ definition layers; it pins the
    layer used and flags the two-layer hazard (PWA-3, \S\ref{sec:pm}).
\end{itemize}
\end{keybox}

\begin{derivedbox}[title={Discovery Note}]
The Council was tasked to close GAP-ZC81-01 ``full system.'' The first
move tried to propagate the bridge through the \emph{shared universal
coupling} $G^{\rm eff} = \ln 2/D_0^{\rm eff}$ (the ZC65 mechanism) and
then read off each sector's shift from
$\alpha_i \propto (G^{\rm eff}/d_s - \sin^2\theta_W)$. This predicted
channel shifts of $0.15\%$ (spatial) and $0.41\%$ (matter) --- nowhere
near the $0.09\%/0.07\%$ channels established in ZC85. The mismatch was
the gift: it showed the bridge acts \emph{per sector} through that
sector's own floor $\varepsilon_i$, not through the residual coupling
with its $-\sin^2\theta_W$ subtraction. Once that was clear, and once
$\aS > 0$ was separated from $\GammaS < 0$, the Strong bridge collapsed
onto the EM bridge by spatial-type inheritance.
\end{derivedbox}

\subsubsection{Pre-Work Audit (PWA-1 through PWA-10)}

\begin{formalbox}[title={Pre-Work Audit Record --- ZC86}]
Scanned before writing (Upgrade Protocol v1.2):
\begin{itemize}[leftmargin=2em, noitemsep]
  \item ZC81/THM-ZC81-01, LEM-ZC81-01 --- \textbf{KEY PILLARS}
    ($\epsS=\epsEM$; $\TRone$ identical, spatial-type).
  \item ZC65/THM-ZC65-01, FIND-ZC65-01/02, LEM-ZC65-01 ---
    \textbf{KEY PILLARS} (the EM bridge and its mechanism).
  \item ZC30/FIND-ZC30-02 --- pillar (spatial vs matter-type class).
  \item ZC73/THM-ZC73-01 --- pillar ($\GammaS$ at rung $k=1$).
  \item ZC85/COR-ZC85-01 --- the gate this paper feeds
    ($\epsSC$ magnitude).
  \item Part~J GAP-ZC81-01 + PM-ZC81-RETRACT-01 --- read; this paper
    is the declared successor (PWA-3 PM-registry scan complete).
\end{itemize}

\textbf{PWA-3} (\emph{critical here}): the $\alpha$ symbol carries
\emph{two distinct layer conventions} in the corpus ---
$\alpha_i = (\varphi(d_i)/15)\,G$ (coupling-scale, $G = 0.737$;
ZC29 line, $\aEMc^{R_0}=0.3697$) and
$\alpha_i = (\varphi(d_i)/15)\,\Gamma_{\rm CRF}^{\rm EM}$
(residual-scale, $\Gamma_{\rm CRF}^{\rm EM}=0.0137$; ZC84).
This paper uses the \emph{fractional} bridge $\delta\alpha/\alpha$,
which is layer-independent (the ratio cancels the scale), so both
conventions give the same running shift. The hazard is flagged so no
downstream paper propagates an absolute $\alpha$ across layers.\\
\textbf{PWA-5} (tautology): THM-ZC86-01 is not a substitution
tautology --- it makes a physical claim (the confined sector runs)
resting on the spatial-type \emph{mechanism}, not on relabelling.\\
\textbf{PWA-6} (macro): $\epsS$, $\aS$, $\GammaS$, $\epszA$, $\epszS$
via \verb|\providecommand|; no conflict.\\
\textbf{PWA-7} (domain): Spontaneous-mode physics only; no mind-domain
crossing.\\
\textbf{No duplicate atomic units found.}
\end{formalbox}

\newpage

%%=============================================================
\subsection{LEM-ZC86-01\quad Confinement--Coupling Decoupling}
\label{sec:lem01}
%%=============================================================

\needspace{12\baselineskip}

\begin{formalbox}[title={LEM-ZC86-01: The Bridge Perturbs $\aS > 0$,
  Not $\GammaS < 0$ ($R_0$, Exact)}]

\textbf{Statement.}
The Strong running bridge perturbs the coupling
$\aS = (2/15)\,\Gamma_{\rm CRF}^{\rm EM} > 0$. Confinement is a
property of the decay rate $\GammaS = G/d_s - 1 + \Delta_\Gamma < 0$.
The two are distinct objects; the bridge never references $\GammaS$.
Hence $\GammaS < 0$ places no obstruction on the bridge.

\textbf{Proof.}

\textit{Step 1 (Two distinct objects).}
The coupling $\aS$ is the $\mathbb{Z}_{15}$ charge weight
$\varphi(3)/15 = 2/15$ times the universal rate
$\Gamma_{\rm CRF}^{\rm EM}$ (ZC29/THM-ZC29-01). The decay rate
$\GammaS$ is the rung-$k{=}1$ value of the $\Gamma$-ladder
(ZC73/THM-ZC73-01). They have different units and different roles:
$\aS$ sets coupling strength, $\GammaS$ sets the
instability/decay direction.

\textit{Step 2 ($\aS > 0$ always).}
Since $\varphi(3)/15 = 2/15 > 0$ and
$\Gamma_{\rm CRF}^{\rm EM} = G/d_s - \sin^2\theta_W = +0.0137 > 0$,
the product $\aS = 0.001822 > 0$. The coupling is positive regardless
of the sign of any ladder rung.

\textit{Step 3 (The bridge references only $\aS$ and $\epsS$).}
The ZC65 bridge mechanism (LEM-ZC65-01) modifies the resolution
capacity $\Dzeroeff = D_0 - \beta\,\TRone$ and reads the shift as a
\emph{fractional} change in the coupling
$\delta\alpha/\alpha = (d_s/2)\,\varepsilon\,K^*$. Neither $\Dzeroeff$
nor the fractional shift contains $\GammaS$. The decay rate is simply
absent from the bridge.

Therefore $\GammaS < 0$ cannot obstruct what does not reference it.
\hfill$\square$

\textbf{Physical reading.}
Confinement says the Strong \emph{decay channel} runs the wrong way
(the sector cannot freely radiate); it does not say the Strong
\emph{coupling} vanishes or that the floor crossing changes character.
The floor crossing is spatial-type (LEM-ZC81-01) --- the same physical
event as EM --- so the same bridge applies.

\textbf{Classification:} Exact Lemma.
\end{formalbox}

%%=============================================================
\subsection{THM-ZC86-01\quad The Strong Running Bridge}
\label{sec:thm01}
%%=============================================================

\needspace{12\baselineskip}

\begin{formalbox}[title={THM-ZC86-01: Strong Running Bridge
  ($R_0{\to}R_{\max}$, Near-Formal)}]

\textbf{Statement.}
\begin{equation}
  \boxed{
    \aS(R_{\max})
    \;=\;
    \aS(R_0)\left[1 + \frac{d_s}{2}\,\epsS\,K^*\right],
  }
  \label{eq:strong-bridge}
\end{equation}
with $\epsS = \epsEM$, so the bridge is \emph{identical in form and
magnitude} to the EM bridge (ZC65/THM-ZC65-01).

\textbf{Proof.}

\textit{Step 1 (Spatial-type crossing).}
By LEM-ZC81-01, $\TRone_{\rm strong} = \TRone_{\rm EM}
= (d_s/n_m)\,\epsuniv^2$ --- the Strong and EM sectors undergo the
\emph{same} $R_1$ crossing. The ZC65 bridge is built entirely from
this crossing cost (LEM-ZC65-01: $\Dzeroeff = D_0 - \beta\,\TRone$).

\textit{Step 2 (Same mechanism, same depletion).}
Since the crossing cost is identical, the resolution-capacity
depletion $\beta\,\TRone$ is identical, hence the modified coupling
$G^{\rm eff}$ and the fractional shift are identical to EM.

\textit{Step 3 (Per-sector floor).}
The fractional shift reads
$\delta\aS/\aS = (d_s/2)\,\epsS\,K^*$ (FIND-ZC65-01 applied to the
Strong sector). Since $\epsS = \epsEM$ (THM-ZC81-01),
\[
  \frac{\delta\aS}{\aS}
  = \frac{d_s}{2}\,\epsEM\,K^*
  = 0.001321
  = \frac{\delta\aEMc}{\aEMc}.
\]

\textit{Step 4 (No confinement obstruction).}
By LEM-ZC86-01, $\GammaS < 0$ does not enter the bridge. Steps 1--3
hold unchanged. \hfill$\square$

\textbf{Numerical (T-canonical):}
\begin{center}
{\small\renewcommand{\arraystretch}{1.25}
\begin{tabular}{@{}lrl@{}}
\toprule
Quantity & Value & Source \\
\midrule
$\aS(R_0) = (2/15)\Gamma_{\rm CRF}^{\rm EM}$ & $0.001822$ & ZC29 \\
$(d_s/2)\,\epsS\,K^*$                        & $0.001321$ & THM-ZC81-01 + ZC65 \\
$\aS(R_{\max})$                              & $0.001824$ & this theorem \\
\bottomrule
\end{tabular}}
\end{center}

\textbf{Classification:} Near-Formal Theorem. The form is exact (it
is the EM bridge by spatial-type inheritance); the \emph{near-formal}
label is inherited from ZC65, whose $\varepsilon_0$ is a near-formal
TLS value. \textbf{Closes GAP-ZC81-01.}
\end{formalbox}

%%=============================================================
\subsection{FIND-ZC86-01\quad Spatial-Floor Sharing (Gift)}
\label{sec:find01}
%%=============================================================

\needspace{10\baselineskip}

\begin{giftbox}[title={FIND-ZC86-01: The Running Floor Layer Has Only
  Two Independent Values ($R_0{\to}R_{\max}$, Exact --- Gift)}]

\textbf{Statement.}
At the $R_0{\to}R_{\max}$ running layer, there are exactly two
independent bridge floors:
\begin{equation}
  \epszA = 0.007501\ \text{(EM\,=\,Strong, spatial-type)},
  \qquad
  \epszS = 0.007710\ \text{(Weak, matter-type)}.
  \label{eq:two-floors}
\end{equation}
The Strong sector adds \emph{no} new bridge floor; it shares the EM
floor because both are spatial-type with $\epsS = \epsEM$.

\textbf{Proof.}
The running floor of a sector is fixed by its $R_1$ crossing cost
$\TRone$, which takes only two values (LEM-ZC81-01, FIND-ZC30-02):
$(d_s/n_m)\epsuniv^2$ for spatial-type (EM, Strong) and
$(n_m/d_s)\epsuniv^2$ for matter-type (Weak). EM and Strong therefore
solve the \emph{same} bridge equation and share the floor $\epszA$;
Weak solves the matter-type equation, giving $\epszS$. \hfill$\square$

\textbf{Consequence --- closes the COR-ZC85-01 magnitude question.}
ZC85/FIND-ZC85-01 showed
$\epsSC = \GMfloor(\epszA, \epszS)$, and COR-ZC85-01 gated the
\emph{magnitude} of the $+0.08\%$ on this gap (GAP-ZC81-01). With the
Strong bridge now closed and the running layer shown to carry exactly
the spatial floor $\epszA$ and the matter floor $\epszS$, the $+0.08\%$
is precisely the geometric mean of these two structurally-determined
running floors:
\begin{equation}
  \frac{\epsSC}{\epsuniv}
  = \GMfloor\!\left(\frac{\epszA}{\epsEM},\ \frac{\epszS}{\epsW}\right)
  = \sqrt{(1{+}\tfrac{d_s}{2}\epsEM K^*)(1{+}\tfrac{d_s}{2}\epsW K^*)}\ \text{-type}.
  \label{eq:closes-85}
\end{equation}
The magnitude is no longer gated: it is set by the two spatial/matter
running shifts, both now derived (near-formal, inheriting ZC65).

\textbf{Classification:} Exact Finding (Gift). The two-value structure
is exact; the floor \emph{values} carry ZC65's near-formal status.
\end{giftbox}

%%=============================================================
\subsection{Phase Misalignment --- The Path That Failed}
\label{sec:pm}
%%=============================================================

\begin{openqbox}[title={PM-ZC86-01: Universal-Coupling Propagation
  (Failed Path --- pattern: phase misalignment, wrong layer)}]

\textbf{The probe.}
The first attempt propagated the bridge through the \emph{shared
universal coupling}. The ZC65 mechanism modifies
$\Dzeroeff = D_0 - \beta\,\TRone$, hence $G^{\rm eff} = \ln 2/\Dzeroeff$.
Reading each sector's shift from its absolute coupling
$\alpha_i \propto (G^{\rm eff}/d_s - \sin^2\theta_W)$, with
spatial-type $\TRone$ for EM/Strong and matter-type $\TRone$ for Weak,
predicted fractional shifts:
\[
  \text{spatial channel: } +0.15\%, \qquad
  \text{matter channel: } +0.41\%.
\]

\textbf{Why it failed.}
These do not match the ZC85 channels ($+0.09\%$ EM, $+0.07\%$ Weak;
FIND-ZC85-02). The error: propagating through the \emph{residual}
coupling $G/d_s - \sin^2\theta_W$ injects the $-\sin^2\theta_W$
subtraction into the fractional shift, which does not belong there.
The bridge does not act on the residual; it acts \emph{per sector}
through that sector's own floor $\varepsilon_i$, as
$\delta\alpha_i/\alpha_i = (d_s/2)\,\varepsilon_i\,K^*$.

\textbf{What the failure revealed (the gift inside the scar).}
The mismatch was diagnostic. It forced the separation of two layers:
the bridge propagates per-sector via $\varepsilon_i$ (correct),
\emph{not} through the universal residual coupling (wrong). This is
the same PWA-3 two-layer $\alpha$ hazard, surfaced empirically. Once
corrected, the per-sector shifts ($0.1321\%$ EM\,=\,Strong, $0.1358\%$
Weak) are consistent with ZC85, and THM-ZC86-01 follows.

\textbf{Pattern.} Phase misalignment --- wrong propagation layer
(not an arithmetic error; the arithmetic of the failed path was
correct, but the layer was wrong).

\textbf{Status:} Scar. Preserved as visible scar tissue. The failed
path is retained because it is what isolated the per-sector mechanism;
it is the reason THM-ZC86-01 is correctly scoped.
\end{openqbox}

%%=============================================================
\subsection{Atomic Registry}
\label{sec:registry}
%%=============================================================

\begin{formalbox}[title={Atomic Registry --- ZC86}]
\begin{center}
\renewcommand{\arraystretch}{1.3}
\small
\begin{tabular}{lp{7.4cm}l}
\toprule
\textbf{ID} & \textbf{Description} & \textbf{Status} \\
\midrule
LEM-ZC86-01
  & Bridge perturbs $\aS>0$, not $\GammaS<0$; confinement and
    coupling are distinct objects
  & Exact \\[3pt]
THM-ZC86-01
  & $\aS(R_{\max}) = \aS(R_0)[1+(d_s/2)\epsS K^*]$; = EM bridge by
    spatial-type inheritance; closes GAP-ZC81-01
  & Near-Formal \\[3pt]
FIND-ZC86-01
  & Running layer has two floors $\epszA$ (EM=Strong) and $\epszS$
    (Weak); Strong shares EM's; closes COR-ZC85-01 magnitude (Gift)
  & Exact \\[3pt]
PM-ZC86-01
  & Universal-coupling propagation path failed ($0.15/0.41\%$ vs
    $0.09/0.07\%$); revealed per-sector mechanism
    (pattern: phase misalignment, wrong layer)
  & Scar \\[3pt]
\bottomrule
\end{tabular}
\end{center}
\end{formalbox}

%%=============================================================
\subsection{R-Layer Research Note}
\label{sec:rlayer}
%%=============================================================

\begin{layerbox}[title={R-Layer Assignments for ZC86}]

{\small\renewcommand{\arraystretch}{1.25}
\begin{center}
\begin{tabular}{@{}llll@{}}
\toprule
Atomic unit & R-layer & Cost $T$ & Source \\
\midrule
LEM-ZC86-01 & $R_0$ & algebraic & ZC29 + ZC73 \\
THM-ZC86-01 & $R_0{\to}R_{\max}$ & $\beta\,\TRone$ & ZC65 + ZC81 \\
FIND-ZC86-01 & $R_0{\to}R_{\max}$ & algebraic & ZC30 + ZC65 + ZC85 \\
PM-ZC86-01 & $R_0{\to}R_{\max}$ & --- & failed $G^{\rm eff}$ path \\
\bottomrule
\end{tabular}
\end{center}}

\textbf{Two-layer discipline (PWA-8, PWA-3).}
The load-bearing distinction is between the \emph{coupling} $\aS$
(perturbed by the bridge) and the \emph{decay rate} $\GammaS$
(carries confinement). The failed path PM-ZC86-01 conflated the
per-sector floor layer with the universal-coupling layer; the
correct bridge is per-sector via $\varepsilon_i$. The Strong bridge
is the EM bridge because the crossing layer ($\TRone$, spatial-type)
is shared, not because the sectors are otherwise identical.
\end{layerbox}

%%=============================================================
\subsection{Genealogy Map}
\label{sec:genealogy}
%%=============================================================

\begin{genealogybox}[title={How ZC86's Key Objects Trace Back}]

\textbf{Branch A --- Spatial-type crossing:}
\begin{itemize}[leftmargin=2em, noitemsep]
  \item \textbf{ZC30/FIND-ZC30-02}: spatial vs matter-type class;
    $\TRone = (d_s/n_m)\epsuniv^2$ for spatial.
  \item $\to$ \textbf{ZC81/LEM-ZC81-01}: $\TRone_{\rm strong}
    = \TRone_{\rm EM}$.
  \item $\to$ \textbf{ZC86/THM-ZC86-01}: same crossing $\Rightarrow$
    same bridge.
\end{itemize}

\textbf{Branch B --- The EM bridge:}
\begin{itemize}[leftmargin=2em, noitemsep]
  \item \textbf{ZC65/LEM-ZC65-01}: $\Dzeroeff = D_0-\beta\TRone$.
  \item \textbf{ZC65/FIND-ZC65-01}: $\delta\alpha/\alpha
    = (d_s/2)\varepsilon_0 K^*$.
  \item $\to$ \textbf{ZC86/THM-ZC86-01}: transferred verbatim to
    Strong (Step 2--3).
\end{itemize}

\textbf{Branch C --- Confinement vs coupling:}
\begin{itemize}[leftmargin=2em, noitemsep]
  \item \textbf{ZC73/THM-ZC73-01}: $\GammaS = -0.498$ (rung $k=1$).
  \item \textbf{ZC29/THM-ZC29-01}: $\aS = (2/15)\Gamma_{\rm CRF}>0$.
  \item $\to$ \textbf{ZC86/LEM-ZC86-01}: distinct objects; $\GammaS<0$
    no obstruction.
\end{itemize}

\textbf{Branch D --- The gate it feeds:}
\begin{itemize}[leftmargin=2em, noitemsep]
  \item \textbf{ZC85/FIND-ZC85-01}: $\epsSC = \GMfloor(\epszA,\epszS)$.
  \item \textbf{ZC85/COR-ZC85-01}: magnitude gated on GAP-ZC81-01.
  \item $\to$ \textbf{ZC86/FIND-ZC86-01}: gate opened; two running
    floors fixed; $\epsSC$ magnitude no longer gated.
\end{itemize}

\textbf{Convergence at ZC86:}
Branches A and B close the bridge (Strong = EM via shared crossing);
Branch C dissolves the confinement worry (distinct objects); Branch D
feeds the result back to ZC85, closing the $\epsSC$ magnitude. The
failed path PM-ZC86-01 sits across all four: it is the wrong turn that
forced the per-sector mechanism into view. ZC86's contribution: the
last electroweak/strong running bridge, closed --- and a scar that
shows how a wrong propagation layer can be the most informative probe.
\end{genealogybox}

%%=============================================================
\subsection{Closing Sentence}
%%=============================================================
%% TONE A -- TRIUMPHANT

\bigskip
\begin{center}
\textit{%
  The confined sector was never barred from running.\\[0.3em]
  Confinement lived in the decay rate;\\
  the bridge only ever touched the coupling.\\[0.3em]
  And the wrong road --- through the universal coupling ---\\
  was the one that showed us the right one.\\[0.5em]
  $\displaystyle
    \alpha_{\rm strong}(R_{\max})
    \;=\;
    \alpha_{\rm strong}(R_0)\!\left[1 + \tfrac{d_s}{2}\,\varepsilon_{\rm strong}K^*\right]
    \qquad\text{(exactly the EM bridge).}$
}
\end{center}

% Governance Note: This document follows Upgrade Protocol v1.2
% and CRF Style Guide v3.5, with an additive XeLaTeX bold-font fix
% (fontspec + Latin Modern). No prior style standard reduced.


\newpage
%%=============================================================
\part{ZC87: The Temporal Continuum Bridge (v17.15)}
\label{part:LV}
%%=============================================================

\begin{volumebox}[title={Part LV --- ZC87: CRF's Two Times Meet in the
  Middle (v17.15)}]
\sloppy
\textbf{Part LV} carries a single paper that bridges CRF's discrete
$\delta$-sweep tick to the continuum $dt$ of relativistic time.

\textbf{ZC87} (The Temporal Continuum Bridge, v2) shows that CRF's two
notions of time meet in the middle. FIND-ZC87-01 (recovery) notes the
upper half already existed and was forgotten: ZC33/DEF-ZC33-01 defines
sector physical time. THM-ZC87-01 supplies the lower half --- continuum
$dt$ is the $\delta$-sweep tick made dense, the $\varepsilon_0\!\to\!0$
density limit of ZC51's machinery transferred to the time axis with the
same fixed-observable construction --- closing the structural part of
GAP-ZC76-01. LEM-ZC87-01 proves the bridge is non-circular. FIND-ZC87-02
(bonus, not load-bearing) reads gravitational dilation as a sweep-rate
slowdown.

This is a v2 paper: it simultaneously closes GAP-ZC87-01 (joint-limit
well-definedness, now trivially satisfied because $\tau_{\rm sweep}$ is
held fixed while tick-density $T_{\rm avg}\!\to\!\infty$), and it
preserves two scars. PM-ZC87-01 records an early probe that read
$N_{\rm sweep}=t_{\rm Hubble}/t_{\rm Planck}$ and matched the Dirac large
number $\sim 10^{61}$ --- a coincidence that fed its own input. PM-ZC87-02
preserves v1's spurious $dt=\tau_{\rm sweep}/T_{\rm avg}$ (a sub-tick
quantity below one sweep), corrected by the v2 fixed-observable probe.

\textbf{Key closures:}
GAP-ZC87-01 (joint-limit well-definedness) --- CLOSED in v2.
GAP-ZC76-01 ($\tau_{\rm sweep}$) --- structural bridge CLOSED
(near-formal); seconds-calibration (GtRate-1) remains externally blocked.

\textbf{Phase Misalignments preserved:}
PM-ZC87-01 (Dirac-large-number probe that fed its own input; pattern:
circular anchor);
PM-ZC87-02 (v1 sub-tick formula $dt=\tau_{\rm sweep}/T_{\rm avg}$;
corrected in v2).
\end{volumebox}

\newpage
%%─── VERBATIM EMBED: ZC87 ───────────────────────────────────────────────
\addcontentsline{toc}{section}{ZC87: The Temporal Continuum Bridge}

% [BRIDGE-PROSE: integration-added, Extension Freedom narrative, v3.6 §31.3]
CRF has always carried a quiet tension about time. At its foundation, time
is discrete: reality advances in distinguishing steps, one sweep at a
time, like the ticks of a clock that has no in-between. But the physics it
hopes to recover --- relativity, dilation, the smooth flow of a worldline
--- is written in continuous time, the $dt$ of the textbooks. A framework
cannot claim both and leave the seam unexamined. Either the discrete ticks
and the continuous flow are the same thing seen at two resolutions, or CRF
has two notions of time that do not meet, which would be a crack running
straight through the foundation.

This paper closes the seam, and it does so in two halves that turn out to
already fit. The upper half --- how a sector's physical time is defined ---
had in fact been built several papers earlier and then half-forgotten;
part of the work here was simply recovering it. The lower half is new: the
continuous $dt$ is what the discrete tick becomes when you let the ticks
grow dense, the same limiting move CRF uses elsewhere, now applied to the
time axis. The two halves meet in the middle, and the clock that had no
in-between acquires one.

The sections below also keep two scars in plain view, and they are worth
the honesty: an early version of this argument produced a seductive
numerical coincidence that, on inspection, was feeding on its own input,
and a first formula that placed an event below a single tick --- both
preserved here as the wrong turns that mapped the right road.



%%=============================================================
\begin{abstract}
\sloppy
CRF carries \emph{two} notions of time, used side by side but never
derived from each other: a \textbf{lower}, discrete one --- the
$\delta$-sweep count, with $\Tav$ measured in sweeps --- and an
\textbf{upper}, continuum one --- the $dt$ of the $R_1$ wave layer
(DEF-RN02: $T(R_1) = \varepsilon_0^2\,\Sinv\,dt$; the rate laws
$\dot H/H$, $\dot P/P$, $d\tau/dt$). The middle was missing.

\textbf{FIND-ZC87-01}: the upper half already exists. ZC33/DEF-ZC33-01
defines sector physical time $t_i = (\Sinv\,dt)\,\Tav$ with
$\Sinv\,dt = d_s/n_m$ fixed by ZC30 --- a \emph{dimensionless} CRF
number; the continuum $dt$ is fully absorbed and $t_i$ is in sweeps.

\textbf{THM-ZC87-01} (Temporal Continuum Bridge): the lower half is the
continuum limit of the renewal count. The continuum increment \emph{is}
the tick, $dt = \tsweep$; what makes time \emph{smooth} is that the
number of ticks per renewal window,
$\Tav = 2nK^*/[(n-1)\varepsilon_0]$, diverges as $\varepsilon_0\to 0$:
\[
  \Tav \xrightarrow{\varepsilon_0\to 0} \infty
  \quad\Longrightarrow\quad
  \{k\,\tsweep\}_{k} \text{ becomes dense},\qquad dt = \tsweep .
\]
This is the THM-ZC51-01 machinery on the time axis, with the \emph{same}
fixed-observable structure ($\{\text{fixed: }\tsweep\}
\leftrightarrow$ ZC51's $\{\text{fixed: }\lambda_2\Tav\}$). It closes
the \emph{structural} part of GAP-ZC76-01.

\textbf{LEM-ZC87-01}: the bridge is not circular --- the $dt$ in
DEF-RN02 appears only inside the absorbed product $\Sinv\,dt = d_s/n_m$,
while the emergent $dt = \tsweep$ is the calibration tick itself,
built from $dt$-free quantities. The bridge converts exactly one
unknown ($\tsweep$) into $dt$.

\textbf{FIND-ZC87-02} (bonus): beside the Part A relation
$d\tau/dt = 1/\sqrt{\Ccal/\alpha}$, gravitational time dilation reads as
a \emph{slowdown of the sweep-resolution rate} near mass --- testable.

\textbf{GAP-ZC87-01} (closed in v2): the joint limit is well-defined
once $dt = \tsweep$ --- $\tsweep$ is held fixed and only the
tick-density $\Tav$ diverges; no pathological $\tsweep\to 0$ is needed.
\textbf{PM-ZC87-01} preserves a wrong turn: an earlier probe read
$t_{\rm Hubble}/t_{\rm Planck}\sim 10^{61}$ as a CRF \emph{derivation}
of the Dirac number --- it was the ratio of two observed inputs.
\textbf{PM-ZC87-02} preserves a second: v1 wrote $dt = \tsweep/\Tav$, a
sub-tick quantity (below one sweep); the GAP-ZC87-01 probe corrected it
to $dt = \tsweep$.
\end{abstract}


%%=============================================================
\subsection{Background: CRF's Two Times}
\label{sec:background}
%%=============================================================

\begin{keybox}[title={Numerical Summary --- ZC87 (T-canonical)}]
{\small\renewcommand{\arraystretch}{1.25}
\begin{center}
\begin{tabular}{@{}p{6.0cm}p{2.6cm}p{1.4cm}p{2.2cm}@{}}
\toprule
\textbf{Quantity} & \textbf{Value} & \textbf{Unit} & \textbf{Source} \\
\midrule
$\Tav(\varepsilon_{\rm univ}) = 2nK^*/[(n{-}1)\varepsilon_0]$ & $35.345$ & sweeps & ZC62 \\
$(\Sinv\,dt)_{\rm spatial} = d_s/n_m$ & $0.600$ & --- & ZC30 \\
$t_i(\text{spatial}) = (d_s/n_m)\Tav$ & $21.21$ & sweeps & ZC33 \\
$t_i(\text{matter}) = (n_m/d_s)\Tav$ & $58.91$ & sweeps & ZC33 \\
ratio $t_i^{\rm m}/t_i^{\rm s} = (n_m/d_s)^2$ & $25/9$ & --- & ZC30/33 \\
\midrule
$\Tav(\varepsilon_0{=}10^{-6})$ & $268579$ & ticks/ren. & density limit \\
$dt = \tsweep$ & --- & (one calib.) & \textbf{THM-ZC87-01} \\
\bottomrule
\end{tabular}
\end{center}}
\end{keybox}

\subsubsection{The Two Times, Stated Plainly}

\begin{keybox}[title={Lower Time and Upper Time}]
\textbf{Lower time (discrete, derived).}
The $\delta$-cascade advances in \emph{sweeps}. The renewal scale
$\Tav = 2nK^*/[(n-1)\varepsilon_0]$ (COR-ZC62-02) is a pure number of
sweeps; the single-mode fire probability $K^* = 1-e^{-1/n}$ and the
Axiom 10 fixed point $(1-K^*)^{-n}=e$ define the tick internally. No
external datum enters. But the tick has no physical \emph{duration}:
$\Tav$ is dimensionless.

\textbf{Upper time (continuum, used as primitive).}
The $R_1$ wave layer carries a continuum $dt$:
DEF-RN02 gives $T(R_1) = \varepsilon_0^2\,\Sinv\,dt$; the dynamics use
$\dot H/H$, $\dot P/P = \mathcal{G}H_0$ (FIND-ZC76-01),
$d\tau/dt = 1/\sqrt{\Ccal/\alpha}$ (Part A). All take $dt$ as given.

\textbf{The missing middle.}
No equation in the corpus writes $dt = (\text{something})\times$ sweep.
The lower time is counted; the upper time is assumed; the bridge that
turns a count into a duration was never built. GAP-ZC76-01 (GtRate-1,
$\tsweep = t_{\rm Planck}$) is the visible symptom of this gap.
\end{keybox}

\subsubsection{What ZC87 Does NOT Do}

\begin{keybox}[title={Scope Boundaries}]
\begin{itemize}[leftmargin=2em, noitemsep]
  \item Does not calibrate $\tsweep$ in seconds --- that is the
    external (SNSPD) part of GAP-ZC76-01, here \emph{refined} to a
    single constant, not closed.
  \item Does not modify ZC33, ZC51, ZC30, DEF-RN02, or the Part A
    metric derivation (pillars).
  \item FIND-ZC87-02 (dilation) is a bonus reading, not load-bearing;
    the core bridge stands without it.
\end{itemize}
\end{keybox}

\begin{derivedbox}[title={Discovery Note}]
The session began with a question --- is CRF time ``point-like''
(accumulated $\delta$-events) rather than the standard continuum arrow?
--- which sharpened into a structural diagnosis: CRF has two times, a
lower discrete one that ``hasn't reached all the way up'' to a duration
and an upper continuum one that ``hasn't reached down'' to the sweep,
with a missing middle. A corpus scan found the upper half already built
in ZC33 (and forgotten); transferring the THM-ZC51-01 continuum limit
($\varepsilon_0\to 0$) from the instability axis to the time axis
produced the lower half. The Dirac-number probe that failed first
(\S\ref{sec:pm}) is what forced the constant/rate/epoch partition into
view.
\end{derivedbox}

\subsubsection{Pre-Work Audit (PWA-1 through PWA-10)}

\begin{formalbox}[title={Pre-Work Audit Record --- ZC87}]
Scanned before writing (Upgrade Protocol v1.2):
\begin{itemize}[leftmargin=2em, noitemsep]
  \item ZC33/DEF-ZC33-01, THM-ZC33-02 --- \textbf{KEY PILLAR} (the
    upper half: sector physical time $t_i = (\Sinv\,dt)\Tav$; semigroup
    clock $A = \Tav^{-1}M$).
  \item ZC51/THM-ZC51-01 --- \textbf{KEY PILLAR} (continuum-limit
    machinery $\varepsilon_0\to 0$ + LEM-ZC51-02 finite correction).
  \item ZC30/FIND-ZC30-02 --- \textbf{KEY PILLAR}
    ($\Sinv\,dt = d_s/n_m$ spatial, $n_m/d_s$ matter; dimensionless).
  \item ZC62/COR-ZC62-02 --- pillar ($\Tav = 2nK^*/[(n-1)\varepsilon_0]$).
  \item ZC76/FIND-ZC76-01, GAP-ZC76-01 --- the gap refined here.
  \item Part A proper time $d\tau/dt = 1/\sqrt{\Ccal/\alpha}$,
    $\Ccal(r)$ continuum (Instability~\#17) --- pillar for FIND-ZC87-02.
  \item DEF-RN02 (Lexicon Lock v4) --- the $R$-layer table; the
    \emph{only} place $dt$ enters, and only as $\Sinv\,dt$.
\end{itemize}

\textbf{PWA-5} (\emph{critical}): the bridge must not be circular ---
checked explicitly in LEM-ZC87-01 ($dt_{\rm int}$ absorbed vs
$dt_{\rm em}$ built $dt$-free). PM-ZC87-01 records a genuine
fed-vs-derived tautology that was caught.\\
\textbf{PWA-6} (macro): $\tsweep$, $\dtem$, $\dtint$, $\Sinv$, $\Tav$
via \verb|\providecommand|; no conflict.\\
\textbf{PWA-7} (domain): Spontaneous-mode (physics) time only; the
mind-domain S\=ant\=ana continuum (Part F) is a separate object, not
conflated.\\
\textbf{PWA-8} (two-layer): lower (sweep) vs upper (dt) is the entire
subject; the $R_0/R_1$ distinction is load-bearing.\\
\textbf{No duplicate atomic units found} (ZC33 recovered, not
duplicated --- cited as pillar).
\end{formalbox}

\newpage

%%=============================================================
\subsection{FIND-ZC87-01\quad The Upper Half Already Exists (ZC33)}
\label{sec:find01}
%%=============================================================

\needspace{10\baselineskip}

\begin{giftbox}[title={FIND-ZC87-01: ZC33 Sector Physical Time Is the
  Upper Half of the Bridge ($R_0\!\to\!R_1$, Exact --- recovery)}]

\textbf{Statement.}
The sweep-to-clock half of the temporal bridge was built in ZC33 and
not recognised as such. DEF-ZC33-01 defines sector physical time
\begin{equation}
  t_i \;=\; (\Sinv^{(i)}\,dt)\cdot\Tav,
  \label{eq:zc33-time}
\end{equation}
and by ZC30/FIND-ZC30-02 the factor $\Sinv^{(i)}\,dt$ is a
\emph{dimensionless} CRF number:
\begin{equation}
  (\Sinv\,dt)_{\rm spatial} = \frac{d_s}{n_m} = 0.6, \qquad
  (\Sinv\,dt)_{\rm matter} = \frac{n_m}{d_s} = \frac{5}{3}.
  \label{eq:sinv-dt}
\end{equation}
Hence $t_i$ is measured in \emph{sweeps}, and the semigroup generator
$A_i^{R_1} = \Tav^{-1}M_i$ (THM-ZC33-01) runs on this sweep clock
(THM-ZC33-02).

\textbf{Why it is the upper half.}
The bridge we seek spans \{sweep count\} $\leftrightarrow$ \{continuum
$dt$\}. Equation~\eqref{eq:zc33-time} maps the renewal sweep count
$\Tav$ onto a sector clock $t_i$ --- the half that ascends from the
discrete count to a clock. What it does \emph{not} do is give that
clock a continuum unit; $t_i$ remains a sweep count scaled by
$d_s/n_m$.

\textbf{Numerical (T-canonical, $\varepsilon_{\rm univ}$):}
$\Tav = 35.345$ sweeps;
$t_i^{\rm spatial} = 0.6\times 35.345 = 21.21$ sweeps;
$t_i^{\rm matter} = (5/3)\times 35.345 = 58.91$ sweeps;
ratio $25/9$ (the same spatial/matter ratio as $T(R_1)$, ZC81).

\textbf{Classification:} Exact Finding (recovery of an existing object;
no new content, re-purposed as the bridge's upper half).
\end{giftbox}

%%=============================================================
\subsection{THM-ZC87-01\quad The Temporal Continuum Bridge}
\label{sec:thm01}
%%=============================================================

\needspace{12\baselineskip}

\begin{formalbox}[title={THM-ZC87-01: $dt$ Is the Tick, Made Dense
  ($\varepsilon_0\!\to\!0$, Near-Formal)}]

\textbf{Statement.}
\begin{equation}
  \boxed{\;
    dt \;=\; \tsweep,
    \qquad\text{continuum}\ \Longleftarrow\
    \Tav(\varepsilon_0) = \frac{2nK^*}{(n-1)\varepsilon_0}
    \xrightarrow{\varepsilon_0\to 0} \infty .
  \;}
  \label{eq:bridge}
\end{equation}
The continuum time increment \emph{is} the physical sweep tick
$\tsweep$. What makes time \emph{smooth} (continuum) is that the number
of ticks per renewal window, $\Tav$, diverges as $\varepsilon_0\to 0$:
the discrete tick-sequence $\{k\,\tsweep\}_{k\in\mathbb{N}}$ becomes
dense relative to any fixed observable window.

\textbf{Proof (near-formal).}

\textit{Step 1 (One tick advances time by $\tsweep$).}
A renewal window of physical duration $D_R$ contains $\Tav$ sweeps. The
time advanced per sweep is $D_R/\Tav$, which is by definition $\tsweep$
(the duration of one sweep). Hence the elementary increment is
$dt = \tsweep$ --- nothing is resolved below one sweep, so no smaller
increment is physical.

\textit{Step 2 (Tick-density diverges; transfer of THM-ZC51-01).}
$\Tav = 2nK^*/[(n-1)\varepsilon_0]\to\infty$ as $\varepsilon_0\to 0$
(numerically $\Tav(10^{-3})=268.6$, $\Tav(10^{-4})=2686$,
$\Tav(10^{-6})=2.69\times10^{5}$). THM-ZC51-01 established that CRF
discrete observables converge to continuum values as $\varepsilon_0\to
0$, with finite offset LEM-ZC51-02. Here the refining observable is the
tick \emph{density} per renewal: as it diverges, the step function
built from $\{k\,\tsweep\}$ converges to a smooth time line.

\textit{Step 3 (Fixed-observable structure matches ZC51 exactly).}
ZC51 holds $\lambda_2\Tav = (n-1)\cdot 29/30$ fixed while
$\lambda_2\to 0$ and $\Tav\to\infty$. The time-axis analogue holds
$\tsweep$ fixed (it is the single calibration constant) while the tick
rate per window $\to\infty$. The fixed object is $\tsweep$; the
diverging object is $\Tav$. No quantity is sent pathologically to a
degenerate value.

\textit{Step 4 (One unknown in, one out).}
$\Tav$ is derived and dimensionless; $\tsweep$ is the single
uncalibrated duration, and $dt = \tsweep$ carries it directly. No
second unknown is introduced (LEM-ZC87-01). \hfill$\square$

\textbf{Why ``Near-Formal''.}
Step 2 invokes the THM-ZC51-01 continuum-convergence result by
transfer; a fully rigorous time-axis proof would carry a
LEM-ZC51-02-type finite-$\varepsilon_0$ density correction, which is
expected to be small and is not computed here. The \emph{relation}
$dt = \tsweep$ and the divergence of $\Tav$ are exact.

\textbf{What this closes.}
The \emph{structural} half of GAP-ZC76-01: the sweep tick and the
continuum $dt$ are identified, and the continuum character is derived
from tick density. What remains is a single calibration ($\tsweep$ in
seconds), gated on external experiment --- the gap is refined from
``derive the bridge'' to ``measure one constant.''
\end{formalbox}

\begin{derivedbox}[title={Why $dt = \tsweep$ and not $\tsweep/\Tav$
  (the v1 correction)}]
Version~1 of this paper wrote $dt = \lim \tsweep/\Tav$. That quantity
is \emph{sub-tick}: smaller than one sweep, hence below the resolution
floor where no time is defined. The error was conflating the
\emph{duration of a renewal window} ($\tsweep\Tav$) with the
\emph{per-tick increment} ($\tsweep$). The continuum does not come from
shrinking the tick; it comes from packing more fixed-size ticks into
each window ($\Tav\to\infty$). This is recorded as PM-ZC87-02 and is
what closed GAP-ZC87-01.
\end{derivedbox}

%%=============================================================
\subsection{LEM-ZC87-01\quad The Bridge Is Not Circular}
\label{sec:lem01}
%%=============================================================

\needspace{10\baselineskip}

\begin{formalbox}[title={LEM-ZC87-01: $dt_{\rm int}$ (absorbed) and
  $dt_{\rm em}$ (built $dt$-free) Are Distinct ($R_0$, Exact)}]

\textbf{Statement.}
The continuum $dt$ inside DEF-RN02 and the emergent $dt$ of
THM-ZC87-01 are different objects; no feedback loop exists.

\textbf{Proof.}

\textit{Step 1 ($dt_{\rm int}$ is absorbed).}
In DEF-RN02, $T(R_1) = \varepsilon_0^2\,\Sinv\,dt$. By ZC30
(equating with $T(R_1) = (d_s/n_m)\varepsilon_0^2$ for spatial-type),
\[
  \Sinv\,dt_{\rm int} = \frac{d_s}{n_m} = 0.6.
\]
The symbol $dt_{\rm int}$ never appears alone --- only in the product
$\Sinv\,dt_{\rm int}$, which is a pure CRF number. It is an internal
bookkeeping infinitesimal, already eliminated.

\textit{Step 2 ($dt_{\rm em}$ is built $dt$-free).}
$dt_{\rm em} = \tsweep$, the physical sweep tick, which is the single
calibration constant of the cascade. It is not assembled from
$dt_{\rm int}$ at all; it is an independent physical duration whose
value is set (eventually) by external calibration, while its
\emph{role} as the time increment follows from THM-ZC87-01. In
particular it does not contain $dt_{\rm int}$.

\textit{Step 3 (No loop).}
Since $dt_{\rm em}$ is assembled entirely from $dt$-free quantities and
$dt_{\rm int}$ is absorbed before any rate law uses it, neither feeds
the other. The bridge does not consume its own output. \hfill$\square$

\textbf{Consequence.}
The bridge converts exactly \emph{one} unknown ($\tsweep$) into $dt$;
it creates no second unknown. This is why GAP-ZC76-01 reduces to a
single calibration constant rather than expanding.

\textbf{Classification:} Exact Lemma.
\end{formalbox}

%%=============================================================
\subsection{FIND-ZC87-02\quad Gravitational Dilation as Sweep Slowdown}
\label{sec:find02}
%%=============================================================

\needspace{10\baselineskip}

\begin{giftbox}[title={FIND-ZC87-02: Time Dilation Is a Slowdown of the
  $\delta$-Resolution Rate ($R_1$, Near-Formal --- Gift, bonus)}]

\textbf{Statement.}
Placing $dt_{\rm em}$ beside the Part A proper-time relation
$d\tau/dt = 1/\sqrt{\Ccal/\alpha}$: in the flat region $\Ccal=\alpha$,
$d\tau = dt_{\rm em}$ (one sweep-window per proper-time unit). Near a
mass, $\Ccal\gg\alpha$, so $d\tau < dt_{\rm em}$: fewer sweeps are
resolved per unit proper time. Gravitational time dilation \emph{is}
the slowdown of the $\delta$-resolution (sweep) rate in regions of high
constraint density.

\textbf{Reading.}
A clock deep in a gravitational well ticks slowly, in CRF, because the
$\delta$-cascade resolves fewer sweeps there per unit of the emergent
$dt$. The metric factor $\sqrt{\Ccal/\alpha}$ is the local sweep-rate
suppression. This recasts a kinematic GR effect as a resolution-rate
statement --- and it is falsifiable: any setting where CRF predicts a
sweep-rate change must match the measured time-dilation factor.

\textbf{Status note.}
This is a \emph{bonus} reading, registered separately from the core
bridge. It rests on the Part A metric derivation (near-formal); if it
were wrong, THM-ZC87-01 and LEM-ZC87-01 would stand unchanged. It is
\emph{not} load-bearing.

\textbf{Classification:} Near-Formal Finding (Gift, bonus).
\end{giftbox}

%%=============================================================
\subsection{Phase Misalignment and Open Gap}
\label{sec:pm}
%%=============================================================

\begin{openqbox}[title={PM-ZC87-01: The Dirac-Number Probe That Fed Its
  Own Answer (pattern: tautology test)}]

\textbf{The probe.}
While attacking GAP-ZC76-01, a probe computed
\[
  \Nsweep \;=\; \frac{t_{\rm Hubble}}{t_{\rm Planck}}
  \;\approx\; 8.4\times10^{60},
\]
and noted it matched the Dirac large number $\sim 10^{61}$. For a
moment this looked like CRF \emph{predicting} the Dirac number.

\textbf{Why it failed.}
$t_{\rm Hubble}$ is the observed cosmic age and $t_{\rm Planck}$ an
observed constant; their ratio is the ratio of two \emph{measured}
inputs. No CRF derivation occurred --- the number was fed in and read
back out. A search for any path to $10^{60.9}$ from CRF constants alone
($n$, $K^*$, $\mathcal{G}$, $\varphi$, $e$, $|\mathbb{Z}_{15}|$)
returned only non-integer, non-structural exponents: forcing a fit
would be numerology, the pattern retracted in COR-ZC81-02/82-02.

\textbf{What the failure revealed (the gift inside the scar).}
The audit forced a partition of cosmic quantities by predictability:
\begin{itemize}[leftmargin=2em, noitemsep]
  \item \textbf{Constant} ($\tsweep$, the tick): derivable in
    \emph{structure} (Axiom 10), gated only on calibration.
  \item \textbf{Rate} ($\dot G/G = \mathcal{G}H_0$): derivable
    (FIND-ZC76-01).
  \item \textbf{Epoch} ($\Nsweep \approx 10^{61}$, the accumulated
    sweep count): \emph{not} derivable --- it is a clock reading that
    grows every tick, not a structural constant. Its
    non-derivability is a feature: a theory that derives constants and
    rates should \emph{not} output the present count.
\end{itemize}
This partition is precisely why THM-ZC87-01 closes the \emph{structural}
bridge but not the epoch: $dt$ emerges, the rate is fixed, but the
accumulated $\Nsweep$ remains an observer's clock reading.

\textbf{Pattern.} Tautology test --- fed an observed quantity, read it
back as derived. Caught by provenance audit (PWA-5).

\textbf{Status:} Scar. Preserved: it is what produced the
constant/rate/epoch partition.
\end{openqbox}

\begin{formalbox}[title={COR-ZC87-01 (was GAP-ZC87-01): Joint Limit Is
  Well-Defined --- CLOSED}]

\textbf{Resolution.}
The v1 worry --- that the limit needed $\tsweep\to 0$ with
$\Tav\to\infty$ on a constrained surface --- dissolves under the
corrected relation $dt = \tsweep$ (THM-ZC87-01). The limit holds
$\tsweep$ \emph{fixed} (it is the single calibration constant) and
sends only the tick-density $\Tav\to\infty$. This is exactly the ZC51
fixed-observable structure: ZC51 fixes $\lambda_2\Tav$ while
$\lambda_2\to 0$; ZC87 fixes $\tsweep$ while the per-window tick count
$\to\infty$. No quantity is driven to a degenerate value, so the limit
is trivially well-defined and yields the finite, nonzero increment
$dt = \tsweep$.

\textbf{Verification (four consistency tests, all pass).}
\begin{enumerate}[leftmargin=2em, noitemsep]
  \item Part A $d\tau/dt = 1/\sqrt{\Ccal/\alpha}$: flat region gives
    $d\tau = \tsweep$; near mass $d\tau < \tsweep$ (dilation,
    FIND-ZC87-02). Consistent.
  \item Tick-density: $\Tav(\varepsilon_0)\to\infty$ packs fixed-size
    ticks densely; the step sequence $\{k\tsweep\}$ becomes smooth.
    Consistent.
  \item Dimensions: $dt = \tsweep$ is time $=$ time; $\tsweep/\Tav$
    would be sub-tick (rejected).
  \item ZC33: $t_i^{\rm phys} = (d_s/n_m)\Tav\,\tsweep$ has per-tick
    increment $\tsweep$. Consistent.
\end{enumerate}

\textbf{Classification:} Corollary (gap closure, Near-Formal). A fully
rigorous time-axis density-convergence proof (LEM-ZC51-02-type
correction) would upgrade this to Exact; the structure is settled.
\end{formalbox}

\begin{openqbox}[title={PM-ZC87-02: The Sub-Tick Formula
  (pattern: phase misalignment, wrong scaling layer)}]

\textbf{The probe.}
Version~1 of this paper stated
$dt = \lim_{\varepsilon_0\to 0}\tsweep/\Tav$, reasoning that a renewal
window of duration $\tsweep\Tav$ should be divided by its sweep count
to give the increment.

\textbf{Why it failed.}
$\tsweep/\Tav$ is \emph{sub-tick}: smaller than one sweep, hence below
the resolution floor where no time is defined. The error conflated the
\emph{renewal-window duration} ($\tsweep\Tav$) with the \emph{per-tick
increment} ($\tsweep$). Three trial fixed-observable choices all sent
$dt\to 0$, exposing the mistake.

\textbf{What the failure revealed.}
The continuum does not arise from shrinking the tick but from packing
more fixed-size ticks per window ($\Tav\to\infty$). Correcting to
$dt = \tsweep$ simultaneously closed GAP-ZC87-01 (now COR-ZC87-01):
the confusing joint limit was an artifact of the wrong formula.

\textbf{Pattern.} Phase misalignment --- wrong scaling layer
(window-duration vs per-tick increment). The arithmetic was fine; the
physical identification of $dt$ was off by a factor $\Tav$.

\textbf{Status:} Scar. Preserved: it is what closed the gap.
\end{openqbox}

\begin{openqbox}[title={GAP-ZC76-01: Refined to a Single Calibration}]

\textbf{Refinement.}
With THM-ZC87-01 ($dt = \tsweep$) and COR-ZC87-01 (limit well-defined),
GAP-ZC76-01 (GtRate-1, $\tsweep = t_{\rm Planck}$) splits:
\begin{itemize}[leftmargin=2em, noitemsep]
  \item \textbf{Structural bridge sweep$\leftrightarrow dt$:} CLOSED
    (THM-ZC87-01 + COR-ZC87-01, near-formal).
  \item \textbf{Calibration $\tsweep$ in seconds:} still GATED on
    external SNSPD timing. The gap is now exactly \emph{one} constant.
\end{itemize}

\textbf{Status:} Refined (not closed). The external block is now a
single-number calibration ($\tsweep$), not a missing derivation.
\end{openqbox}

%%=============================================================
\subsection{Atomic Registry}
\label{sec:registry}
%%=============================================================

\begin{formalbox}[title={Atomic Registry --- ZC87}]
\begin{center}
\renewcommand{\arraystretch}{1.3}
\small
\begin{tabular}{lp{7.2cm}l}
\toprule
\textbf{ID} & \textbf{Description} & \textbf{Status} \\
\midrule
FIND-ZC87-01
  & ZC33 sector physical time $t_i=(\Sinv dt)\Tav=(d_s/n_m)\Tav$ is the
    bridge's upper half (sweep$\to$clock); $dt$ absorbed
  & Exact (recovery) \\[3pt]
THM-ZC87-01
  & $dt = \tsweep$; continuum from tick-density $\Tav\to\infty$;
    transfers THM-ZC51-01 fixed-observable limit to the time axis;
    closes structural GAP-ZC76-01
  & Near-Formal \\[3pt]
LEM-ZC87-01
  & $dt_{\rm int}$ (absorbed in $d_s/n_m$) $\neq$ $dt_{\rm em}$ (built
    $dt$-free); bridge not circular; one unknown in, one out
  & Exact \\[3pt]
FIND-ZC87-02
  & gravitational dilation $=$ sweep-rate slowdown
    ($d\tau/dt=1/\sqrt{\Ccal/\alpha}$); testable; bonus, not
    load-bearing
  & Near-Formal (Gift) \\[3pt]
PM-ZC87-01
  & $t_{\rm Hubble}/t_{\rm Planck}\sim10^{61}$ mistaken for derivation;
    fed observed inputs; produced constant/rate/epoch partition
  & Scar \\[3pt]
PM-ZC87-02
  & v1 wrote $dt = \tsweep/\Tav$ (sub-tick); corrected to
    $dt = \tsweep$; pattern: wrong scaling layer
  & Scar \\[3pt]
COR-ZC87-01
  & (was GAP-ZC87-01) joint limit well-defined: $\tsweep$ fixed,
    tick-density $\Tav\to\infty$; four tests pass
  & Near-Formal \\[3pt]
GAP-ZC76-01
  & refined: structural bridge closed; $\tsweep$ calibration gated
    on external
  & Refined \\[3pt]
\bottomrule
\end{tabular}
\end{center}
\end{formalbox}

%%=============================================================
\subsection{R-Layer Research Note}
\label{sec:rlayer}
%%=============================================================

\begin{layerbox}[title={R-Layer Assignments for ZC87}]

{\small\renewcommand{\arraystretch}{1.25}
\begin{center}
\begin{tabular}{@{}llll@{}}
\toprule
Atomic unit & R-layer & Cost $T$ & Source \\
\midrule
FIND-ZC87-01 & $R_0\!\to\!R_1$ & $(\Sinv dt)\Tav$ & ZC33 + ZC30 \\
THM-ZC87-01 & $R_0\!\to\!R_1$ & tick-density limit & ZC51 + ZC62 \\
LEM-ZC87-01 & $R_0$ & algebraic & DEF-RN02 + ZC30 \\
FIND-ZC87-02 & $R_1$ & metric & Part A ($d\tau/dt$) \\
COR-ZC87-01 & $R_0\!\to\!R_1$ & limit closure & ZC51 transfer \\
\bottomrule
\end{tabular}
\end{center}}

\textbf{Two-layer discipline (PWA-8).}
The entire paper lives on the $R_0\!\to\!R_1$ crossing: $R_0$ supplies
the discrete sweep (lower time), $R_1$ supplies the continuum $dt$
(upper time). The bridge is the crossing itself. The load-bearing
caution (LEM-ZC87-01) is keeping $dt_{\rm int}$ (an $R_1$ bookkeeping
infinitesimal, absorbed) distinct from $dt_{\rm em}$ (the $R_0\!\to\!R_1$
coarse-graining output). Conflating them would make the bridge
circular.
\end{layerbox}

%%=============================================================
\subsection{Genealogy Map}
\label{sec:genealogy}
%%=============================================================

\begin{genealogybox}[title={How ZC87's Key Objects Trace Back}]

\textbf{Branch A --- The upper half (sweep $\to$ clock):}
\begin{itemize}[leftmargin=2em, noitemsep]
  \item \textbf{ZC30/FIND-ZC30-02}: $\Sinv\,dt = d_s/n_m$
    (dimensionless; spatial/matter $25/9$).
  \item \textbf{ZC62/COR-ZC62-02}: $\Tav = 2nK^*/[(n-1)\varepsilon_0]$.
  \item \textbf{ZC33/DEF-ZC33-01}: $t_i = (\Sinv dt)\Tav$; semigroup
    clock (THM-ZC33-02).
  \item $\to$ \textbf{ZC87/FIND-ZC87-01}: recognised as the bridge's
    upper half.
\end{itemize}

\textbf{Branch B --- The lower half (sweep $\to$ $dt$):}
\begin{itemize}[leftmargin=2em, noitemsep]
  \item \textbf{ZC51/THM-ZC51-01}: continuum limit $\varepsilon_0\to0$
    + LEM-ZC51-02 correction.
  \item $\to$ \textbf{ZC87/THM-ZC87-01}: same machinery on the time
    axis; $dt = \tsweep$, continuum from tick-density $\Tav\to\infty$.
\end{itemize}

\textbf{Branch C --- Non-circularity:}
\begin{itemize}[leftmargin=2em, noitemsep]
  \item \textbf{DEF-RN02}: $T(R_1)=\varepsilon_0^2\Sinv\,dt$
    ($dt_{\rm int}$ enters).
  \item \textbf{ZC30}: absorbs it into $d_s/n_m$.
  \item $\to$ \textbf{ZC87/LEM-ZC87-01}: $dt_{\rm int}\neq dt_{\rm em}$;
    no loop.
\end{itemize}

\textbf{Branch D --- The bonus, the gate, and the correction:}
\begin{itemize}[leftmargin=2em, noitemsep]
  \item \textbf{Part A}: $d\tau/dt = 1/\sqrt{\Ccal/\alpha}$;
    $\Ccal(r)$ continuum (\#17).
  \item $\to$ \textbf{ZC87/FIND-ZC87-02}: dilation $=$ sweep slowdown.
  \item \textbf{ZC76/GAP-ZC76-01} $\to$ refined to one calibration.
  \item \textbf{COR-ZC87-01} (was GAP-ZC87-01): joint limit closed once
    the v1 sub-tick formula was corrected (PM-ZC87-02) to $dt=\tsweep$.
\end{itemize}

\textbf{Convergence at ZC87:}
Branch A (upper half, forgotten in ZC33) and Branch B (lower half, the
ZC51 transfer) meet in the middle --- the missing bridge between CRF's
two times. Branch C guards the join against circularity; Branch D adds
a testable dilation reading, refines the external gap to a single
constant, and --- through the corrected formula --- closes the
joint-limit question. ZC87's contribution: continuum time is no longer
a primitive in CRF --- it is the $\delta$-sweep tick made dense, and
the only thing still owed to experiment is the duration of one tick.
\end{genealogybox}

%%=============================================================
\subsection{Closing Sentence}
%%=============================================================
%% TONE B -- REFLECTIVE

\bigskip
\begin{center}
\textit{%
  CRF kept two clocks: one that counted, one that flowed.\\[0.3em]
  The counting clock was built in ZC33 and set aside;\\
  the flowing clock was borrowed, never grounded.\\[0.3em]
  We first thought the flow was the tick divided by the count ---\\
  but that is smaller than a tick, and nothing lives below a tick.\\[0.3em]
  The flow \emph{is} the tick. What makes it smooth\\
  is letting each window hold uncountably many of them.\\[0.5em]
  $\displaystyle
    dt \;=\; \tau_{\rm sweep},
    \qquad T_{\rm avg}\xrightarrow{\;\varepsilon_0\to 0\;}\infty.$\\[0.3em]
  What time \emph{is}, CRF now says.\\
  How long one tick \emph{lasts} --- that, still, the world must tell us.
}
\end{center}

% Governance Note: This document follows Upgrade Protocol v1.2
% and CRF Style Guide v3.5, with the additive XeLaTeX bold-font fix
% (fontspec + Latin Modern) carried from ZC86. No prior standard reduced.


\newpage
%%=============================================================
\part{Cross-Part Connection Map v17.15 (Part K)}
\label{part:LVI}
%%=============================================================

\begin{conmapbox}[title={Cross-Part Connection Map v17.15 --- Reading
  Conventions}]
\small
This map records, for each paper in Part K (ZC83--87):
(1) \textbf{Backward Inheritance} --- which objects from Parts A--J
are used as input; (2) \textbf{Forward Corrections} --- which
prior-Part gaps are closed or refined by Part K;
(3) \textbf{Cross-Part Phase Misalignments} --- PMs that span more
than one Part. Source files are never modified (No-Loss).

\textbf{Notation:}
\textcolor{crfgreen}{$\checkmark$ CLOSED} = gap resolved.
$\star$ PROMOTED = conjecture promoted to theorem.
\textbf{PM-V17\_15-EPSFAMILY:} local overrides
$\backslash$\texttt{epsEM}$\to\varepsilon_{\rm EM}$,
$\backslash$\texttt{epscan}$\to\varepsilon_{\rm can}$,
$\backslash$\texttt{epsSC}$\to\varepsilon_{\rm SC}$,
$\backslash$\texttt{epsS}$\to\varepsilon_{\rm strong}$
(bare sector forms, Style Guide v3.5); extends PM-V17\_13-EPSEM of
Part J. Chain and sources unmodified.
\end{conmapbox}

\begin{formalbox}[title={$\varepsilon$ FLOOR VALUES --- VERIFIED CONSISTENT --- DO NOT ``RECONCILE'' (read before editing any $\varepsilon$ line)}]
\small
If you (human or AI) see the near-equal floor values
$0.007550$, $0.007599$, $0.007605$ appearing alongside
\textbf{COR-ZC67-01} (and now \textbf{FIND-ZC83-01}) and conclude they
contradict --- \textbf{they do not.} This was audited line-by-line across
ZC66--87 and is the single most common trap in this corpus.

\medskip
\textbf{COR-ZC67-01 / FIND-ZC83-01 are ONE exact identity},
$\sqrt{\varepsilon_{\rm EM}\cdot\varepsilon_W}=\varepsilon_{\rm base}$.
It evaluates to a \emph{different number on each layer's input pair}:
\begin{itemize}[noitemsep, leftmargin=1.6em]
  \item ZC66/67 defined pair ($\varepsilon_{\rm can}$ base)
    $\to 0.007550 = \varepsilon_{\rm can}$
  \item ZC70/81/82/83 defined pair ($\varepsilon_{\rm univ}$ base)
    $\to 0.007599 = \varepsilon_{\rm univ}$ (the $R_0$-layer GM;
    FIND-ZC83-01)
  \item observed/empirical pair (the two FIND-ZC65-02 bridge floors)
    $\to 0.007605 = \varepsilon_{\rm SC}$ --- the $R_{\max}$-layer GM
    (FIND-ZC85-01); its $+0.08\%$ above $\varepsilon_{\rm univ}$ is the
    geometric mean of the two running shifts, gated on and then released
    by GAP-ZC81-01 (ZC86/FIND-ZC86-01)
\end{itemize}
These three are \emph{intentionally distinct} (different input pairs at
different $R$-layers), not rounding drift and not a contradiction.
ZC83--87 add the layer interpretation but change \emph{no} value.

\medskip
\textbf{DIRECTIVE.} Do not change any $\varepsilon$ value, do not
``reconcile'' them, do not open an audit, and do not edit the
COR-ZC67-01, FIND-ZC83-01, or FIND-ZC85-01 lines to fix a perceived
conflict. If a conflict still seems present, you are reading one
identity as if it were three values --- re-read this box, not the
papers. (ZC66--69 audited and correct; ZC81--87 correct.)
\end{formalbox}

%%─── ZC83 Connection Record ─────────────────────────────────────────────
\subsection{Connection Records for ZC83}

\begin{inheritbox}[title={ZC83 --- Backward Inheritance from Parts A--J}]
\begin{itemize}[noitemsep, leftmargin=1.5em]
  \item ZC67/LEM-ZC67-01: sector-floor formula (source of the ordering)
  % [DO-NOT-RECONCILE] COR-ZC67-01 = ONE identity on three input pairs: 0.007550 (eps_can) / 0.007599 (eps_univ, this layer) / 0.007605 (eps_SC, observed). See eps-notice box at head of Part LVI. Do NOT edit to "fix" a perceived conflict.
  \item ZC67/COR-ZC67-01: $\sqrt{\varepsilon_{\rm EM}\cdot\varepsilon_W}$ geometric-mean structure
  \item ZC69/COR-ZC69-02: $\varepsilon_{\rm univ}$ exact closed form
  \item ZC70/THM-ZC70-01: $\varepsilon_{\rm EM}$, $\varepsilon_W$ exact values
  \item ZC71/THM-ZC71-01: two-layer separation (gap irreducibility)
  \item ZC79/THM-ZC79-01: $\varepsilon_{\rm can}$ exact
  \item ZC57/THM-ZC57-01: sector-weight ratio $r=C_W/C_{\rm EM}$
\end{itemize}
\end{inheritbox}

\begin{correctionbox}[title={ZC83 --- Forward Corrections / Gaps}]
\begin{itemize}[noitemsep, leftmargin=1.5em]
  \item Floor lattice ordering $\varepsilon_{\rm EM}<\varepsilon_{\rm can}
    <\varepsilon_{\rm univ}<\varepsilon_W$: promoted to forced theorem
    (THM-ZC83-01) from prior numerical observation.
  \item FIND-ZC83-01: $\sqrt{\varepsilon_{\rm EM}\varepsilon_W}
    =\varepsilon_{\rm univ}$ exact at the $R_0$ layer (gift).
  \item GAP-ZC83-01 ($\varepsilon_{\rm SC}$ ontology): \textbf{opened}
    here; closed in Part LIV.
\end{itemize}
\end{correctionbox}

%%─── ZC84 Connection Record ─────────────────────────────────────────────
\subsection{Connection Records for ZC84}

\begin{inheritbox}[title={ZC84 --- Backward Inheritance from Parts A--J}]
\begin{itemize}[noitemsep, leftmargin=1.5em]
  \item ZC20/THM-ZC20-01: $\sin^2\theta_W = 1-\ln8/\ln15$
  \item ZC29/THM-ZC29-01: $\alpha_i = \varphi(d_i)/15\cdot\Gamma_{\rm CRF}$
  \item ZC72/THM-ZC72-01: $\alpha_{\rm grav}(R_0)=\Gamma_{\rm CRF}/15$
  \item ZC74/THM-ZC74-01: $\mathbb{Z}_{15}$ unique squarefree semiprime
  \item ZC75/THM-ZC75-01: $\alpha_{\rm grav}\cdot|Z_{15}|=\Gamma_{\rm CRF}$
  \item ZC80/THM-ZC80-01: 7 EW observables at $R_0$, zero free params
  \item ZC81/THM-ZC81-01, FIND-ZC81-01: $\varepsilon_{\rm strong}
    =\varepsilon_{\rm EM}$; mirror symmetry
  \item ZC82/THM-ZC82-01: $w_{\rm DE}+3\Delta\Gamma=\Gamma_{\rm CRF}$
  \item ZC83/THM-ZC83-01, FIND-ZC83-01: ordering and GM identity
\end{itemize}
\end{inheritbox}

\begin{correctionbox}[title={ZC84 --- Forward Corrections / Gaps}]
\begin{itemize}[noitemsep, leftmargin=1.5em]
  \item THM-ZC84-01: eight prior results assembled into one synthesis
    (no new derivation --- see PM-ZC84-SCOPE-01).
  \item COR-ZC84-01: four-sector falsification table, zero free params.
  \item Does not close GAP-ZC83-01, GAP-ZC81-01, or GAP-ZC76-01.
\end{itemize}
\end{correctionbox}

%%─── ZC85 Connection Record ─────────────────────────────────────────────
\subsection{Connection Records for ZC85}

\begin{inheritbox}[title={ZC85 --- Backward Inheritance from Parts A--J}]
\begin{itemize}[noitemsep, leftmargin=1.5em]
  \item ZC65/FIND-ZC65-02: two per-observable bridge floors
    ($\varepsilon_0^{\alpha}$, $\varepsilon_0^{\rm sw}$)
  % [DO-NOT-RECONCILE] FIND-ZC85-01 is COR-ZC67-01's GM structure at the R_max layer; result 0.007605 = eps_SC is the observed-pair reading, distinct from eps_univ (0.007599). See eps-notice box. Do NOT edit to "fix".
  \item ZC67/COR-ZC67-01, LEM-ZC67-01: GM structure
  \item ZC68/THM-ZC68-01b: $\varepsilon_{\rm SC}$ observed bracket
  \item ZC69/COR-ZC69-02: $\varepsilon_{\rm univ}$ exact
  \item ZC70/THM-ZC70-01: $\varepsilon_{\rm EM}$, $\varepsilon_W$ exact
  \item ZC81/GAP-ZC81-01: Strong/EW running bridge (the gate)
  \item ZC83/GAP-ZC83-01, FIND-ZC83-01: the gap refined; $R_0$-layer GM
\end{itemize}
\end{inheritbox}

\begin{correctionbox}[title={ZC85 --- Forward Corrections / Gaps}]
\begin{itemize}[noitemsep, leftmargin=1.5em]
  \item GAP-ZC83-01 ($\varepsilon_{\rm SC}$ ontology):
    \textbf{structurally} \textcolor{crfgreen}{$\checkmark$ CLOSED}
    (FIND-ZC85-01: $\varepsilon_{\rm SC}=$ GM of the two bridge floors).
  \item GAP-ZC83-01 magnitude: refined to \textbf{gated} on GAP-ZC81-01
    (COR-ZC85-01), no longer an independent open problem.
\end{itemize}
\end{correctionbox}

%%─── ZC86 Connection Record ─────────────────────────────────────────────
\subsection{Connection Records for ZC86}

\begin{inheritbox}[title={ZC86 --- Backward Inheritance from Parts A--J}]
\begin{itemize}[noitemsep, leftmargin=1.5em]
  \item ZC65/THM-ZC65-01, FIND-ZC65-01/02, LEM-ZC65-01: EM bridge
  \item ZC30/FIND-ZC30-02: spatial vs matter-type sector class
  \item ZC73/THM-ZC73-01: $\Gamma$ ladder; $\Gamma_{\rm strong}$ at $k=1$
  \item ZC81/THM-ZC81-01: $\varepsilon_{\rm strong}=\varepsilon_{\rm EM}$;
    LEM-ZC81-01: $T_{R_1}$ crossing cost; GAP-ZC81-01 (the gate it closes)
  \item ZC85/COR-ZC85-01: $\varepsilon_{\rm SC}$ magnitude gated on
    GAP-ZC81-01
\end{itemize}
\end{inheritbox}

\begin{correctionbox}[title={ZC86 --- Forward Corrections / Gap Closures}]
\begin{itemize}[noitemsep, leftmargin=1.5em]
  \item GAP-ZC81-01 (Strong running bridge):
    \textcolor{crfgreen}{$\checkmark$ CLOSED} (THM-ZC86-01, near-formal;
    bridge form and magnitude $=$ EM).
  \item GAP-ZC83-01 magnitude (via COR-ZC85-01 gate):
    \textcolor{crfgreen}{$\checkmark$ CLOSED} (FIND-ZC86-01: the running
    floor layer has exactly two floors; $\varepsilon_{\rm SC}=$ their GM).
  \item LEM-ZC86-01: confinement decoupling --- confinement lives in
    $\Gamma$, not in $\alpha$ (removes the only obstruction).
\end{itemize}
\end{correctionbox}

%%─── ZC87 Connection Record ─────────────────────────────────────────────
\subsection{Connection Records for ZC87}

\begin{inheritbox}[title={ZC87 --- Backward Inheritance from Parts A--J}]
\begin{itemize}[noitemsep, leftmargin=1.5em]
  \item ZC33/DEF-ZC33-01: sector physical time (the recovered upper half);
    THM-ZC33-02: semigroup clock
  \item ZC51/THM-ZC51-01: continuum-limit machinery
    ($\varepsilon_0\!\to\!0$)
  \item ZC30/FIND-ZC30-02: $\Sigma_{\rm inv}\,dt = d_s/n_m$; sector class
  \item ZC62/COR-ZC62-02: $T_{\rm avg}=2nK^*/[(n-1)\varepsilon]=eF/K^*$
  \item ZC76/FIND-ZC76-01: $\mathcal{G}H_0$ rate; GAP-ZC76-01 (GtRate-1)
  \item Part A: the $\Sigma_{\rm inv}\,dt$ relation (FIND-ZC87-02 bonus)
\end{itemize}
\end{inheritbox}

\begin{correctionbox}[title={ZC87 --- Forward Corrections / Gap Closures}]
\begin{itemize}[noitemsep, leftmargin=1.5em]
  \item GAP-ZC87-01 (joint-limit well-definedness):
    \textcolor{crfgreen}{$\checkmark$ CLOSED} in v2 ($\tau_{\rm sweep}$
    fixed, $T_{\rm avg}\!\to\!\infty$).
  \item GAP-ZC76-01 ($\tau_{\rm sweep}$): structural sweep$\leftrightarrow
    dt$ bridge \textcolor{crfgreen}{$\checkmark$ CLOSED} (THM-ZC87-01,
    near-formal); seconds-calibration remains externally blocked.
\end{itemize}
\end{correctionbox}

%%─── Cross-Part Phase Misalignment Registry (Part K) ────────────────────
\subsection{Cross-Part Phase Misalignment Registry (Part K)}

\begin{correctionbox}[title={Part K --- Phase Misalignments Preserved as
  Scar Tissue}]
\small
\begin{itemize}[noitemsep, leftmargin=1.5em]
  \item \textbf{PM-ZC83-01} --- the ``is $\varepsilon_{\rm SC}$ the
    arithmetic mean of the two floors?'' probe. Pattern: false-pattern
    test; ruled out exactly (AM $>\varepsilon_{\rm univ}$, FIND-ZC83-02).
  \item \textbf{PM-ZC84-SCOPE-01} --- synthesis paper: every identity in
    THM-ZC84-01 is inherited from a prior theorem; none is newly derived.
    Pattern: scope honesty.
  \item \textbf{PM-ZC85-01} --- ``$\varepsilon_{\rm SC}=$ GM of the bridge
    pair'' is a restatement of FIND-ZC65-02 $+$ COR-ZC67-01, not a new
    derivation. Pattern: tautology test (PWA-5).
  \item \textbf{PM-ZC86-01} --- universal-coupling propagation path:
    propagating the bridge through an effective universal coupling
    $G^{\rm eff}$ gave $0.15\%/0.41\%$ errors. Failed; preserved as the
    path that lit the correct route. Pattern: probe malformed.
  \item \textbf{PM-ZC87-01} --- the Dirac-large-number probe:
    $N_{\rm sweep}=t_{\rm Hubble}/t_{\rm Planck}\sim10^{61}$ read as a CRF
    ``prediction'' of the Dirac number, when it was the ratio of two
    observed inputs feeding itself. Pattern: circular anchor.
  \item \textbf{PM-ZC87-02} --- v1's $dt=\tau_{\rm sweep}/T_{\rm avg}$: a
    sub-tick quantity below one sweep; corrected in v2 by the
    fixed-observable probe ($dt=\tau_{\rm sweep}$, density
    $T_{\rm avg}\!\to\!\infty$). Pattern: definition drift.
\end{itemize}
\end{correctionbox}

\vfill
\begin{center}
\small\color{mutedgray}
CRF Complete Edition v17.15 (Part K of 11) ---
Ougit Jarittum (ORCID: 0009-0009-4451-6811)\\
Generated June 2026.
\end{center}

\end{document}
